% DRAFT (2026-10-25), not \input by the paper: predictor-group ablations and augmented-HAR benchmarks.
% Numbers: writeup/generated/rehaul_ablation_numbers.tex (\rhab...) and Table tab:rh_ablation
% (writeup/generated/rehaul_ablation_table.tex), both from writeup/make_rehaul_ablations.py on
% results/ablations_2014start_w3000 (RULE.txt, STATUS.md).  To use: % AUTO-GENERATED by writeup/make_rehaul_ablations.py -- do not edit; re-run the script.
% results/ablations_2014start_w3000 (RULE.txt): column-subset arms of the final 13-bar design,
% frozen hyperparameters, every bar scored as Table 1; DM on daily (average) QLIKE.
\newcommand{\rhabEnetAddEwAllQlike}{0.1871}  % QLIKE enet add_ew all
\newcommand{\rhabEnetAddEwAllMse}{1.47}  % MSE x 1e10 enet add_ew all
\newcommand{\rhabEnetAddEwAllPct}{$-$1.1}  % QLIKE change % vs har
\newcommand{\rhabEnetAddEwAllDm}{$-$2.81}  % DM vs the reference arm
\newcommand{\rhabEnetAddEwAllP}{0.005}  % p-value
\newcommand{\rhabEnetAddEwCloseQlike}{0.1249}  % QLIKE enet add_ew c16
\newcommand{\rhabEnetAddEwCloseMse}{2.10}  % MSE x 1e10 enet add_ew c16
\newcommand{\rhabEnetAddEwClosePct}{$-$1.5}  % QLIKE change % vs har
\newcommand{\rhabEnetAddEwCloseDm}{$-$1.37}  % DM vs the reference arm
\newcommand{\rhabEnetAddEwCloseP}{0.170}  % p-value
\newcommand{\rhabEnetAddEwCloseLateQlike}{0.1120}  % QLIKE enet add_ew c16_2020
\newcommand{\rhabEnetAddEwCloseLateMse}{4.56}  % MSE x 1e10 enet add_ew c16_2020
\newcommand{\rhabEnetAddEwCloseLatePct}{$-$1.5}  % QLIKE change % vs har
\newcommand{\rhabEnetAddEwCloseLateDm}{$-$0.90}  % DM vs the reference arm
\newcommand{\rhabEnetAddEwCloseLateP}{0.369}  % p-value
\newcommand{\rhabEnetAddFomcAllQlike}{0.1847}  % QLIKE enet add_fomc all
\newcommand{\rhabEnetAddFomcAllMse}{1.13}  % MSE x 1e10 enet add_fomc all
\newcommand{\rhabEnetAddFomcAllPct}{$-$2.4}  % QLIKE change % vs har
\newcommand{\rhabEnetAddFomcAllDm}{$-$5.85}  % DM vs the reference arm
\newcommand{\rhabEnetAddFomcAllP}{\ensuremath{<}0.001}  % p-value
\newcommand{\rhabEnetAddFomcCloseQlike}{0.1274}  % QLIKE enet add_fomc c16
\newcommand{\rhabEnetAddFomcCloseMse}{1.85}  % MSE x 1e10 enet add_fomc c16
\newcommand{\rhabEnetAddFomcClosePct}{0.4}  % QLIKE change % vs har
\newcommand{\rhabEnetAddFomcCloseDm}{1.66}  % DM vs the reference arm
\newcommand{\rhabEnetAddFomcCloseP}{0.097}  % p-value
\newcommand{\rhabEnetAddFomcCloseLateQlike}{0.1136}  % QLIKE enet add_fomc c16_2020
\newcommand{\rhabEnetAddFomcCloseLateMse}{3.96}  % MSE x 1e10 enet add_fomc c16_2020
\newcommand{\rhabEnetAddFomcCloseLatePct}{0.0}  % QLIKE change % vs har
\newcommand{\rhabEnetAddFomcCloseLateDm}{$-$0.09}  % DM vs the reference arm
\newcommand{\rhabEnetAddFomcCloseLateP}{0.927}  % p-value
\newcommand{\rhabEnetAddIvAllQlike}{0.1894}  % QLIKE enet add_iv all
\newcommand{\rhabEnetAddIvAllMse}{1.32}  % MSE x 1e10 enet add_iv all
\newcommand{\rhabEnetAddIvAllPct}{0.1}  % QLIKE change % vs har
\newcommand{\rhabEnetAddIvAllDm}{0.23}  % DM vs the reference arm
\newcommand{\rhabEnetAddIvAllP}{0.817}  % p-value
\newcommand{\rhabEnetAddIvCloseQlike}{0.1272}  % QLIKE enet add_iv c16
\newcommand{\rhabEnetAddIvCloseMse}{1.76}  % MSE x 1e10 enet add_iv c16
\newcommand{\rhabEnetAddIvClosePct}{0.3}  % QLIKE change % vs har
\newcommand{\rhabEnetAddIvCloseDm}{0.28}  % DM vs the reference arm
\newcommand{\rhabEnetAddIvCloseP}{0.780}  % p-value
\newcommand{\rhabEnetAddIvCloseLateQlike}{0.1152}  % QLIKE enet add_iv c16_2020
\newcommand{\rhabEnetAddIvCloseLateMse}{3.75}  % MSE x 1e10 enet add_iv c16_2020
\newcommand{\rhabEnetAddIvCloseLatePct}{1.4}  % QLIKE change % vs har
\newcommand{\rhabEnetAddIvCloseLateDm}{0.70}  % DM vs the reference arm
\newcommand{\rhabEnetAddIvCloseLateP}{0.485}  % p-value
\newcommand{\rhabEnetAddLiqAllQlike}{0.1881}  % QLIKE enet add_liq all
\newcommand{\rhabEnetAddLiqAllMse}{1.16}  % MSE x 1e10 enet add_liq all
\newcommand{\rhabEnetAddLiqAllPct}{$-$0.6}  % QLIKE change % vs har
\newcommand{\rhabEnetAddLiqAllDm}{$-$1.54}  % DM vs the reference arm
\newcommand{\rhabEnetAddLiqAllP}{0.124}  % p-value
\newcommand{\rhabEnetAddLiqCloseQlike}{0.1299}  % QLIKE enet add_liq c16
\newcommand{\rhabEnetAddLiqCloseMse}{1.99}  % MSE x 1e10 enet add_liq c16
\newcommand{\rhabEnetAddLiqClosePct}{2.4}  % QLIKE change % vs har
\newcommand{\rhabEnetAddLiqCloseDm}{1.77}  % DM vs the reference arm
\newcommand{\rhabEnetAddLiqCloseP}{0.077}  % p-value
\newcommand{\rhabEnetAddLiqCloseLateQlike}{0.1192}  % QLIKE enet add_liq c16_2020
\newcommand{\rhabEnetAddLiqCloseLateMse}{4.29}  % MSE x 1e10 enet add_liq c16_2020
\newcommand{\rhabEnetAddLiqCloseLatePct}{4.9}  % QLIKE change % vs har
\newcommand{\rhabEnetAddLiqCloseLateDm}{1.90}  % DM vs the reference arm
\newcommand{\rhabEnetAddLiqCloseLateP}{0.058}  % p-value
\newcommand{\rhabEnetAddRmAllQlike}{0.1842}  % QLIKE enet add_rm all
\newcommand{\rhabEnetAddRmAllMse}{1.72}  % MSE x 1e10 enet add_rm all
\newcommand{\rhabEnetAddRmAllPct}{$-$2.7}  % QLIKE change % vs har
\newcommand{\rhabEnetAddRmAllDm}{$-$4.25}  % DM vs the reference arm
\newcommand{\rhabEnetAddRmAllP}{\ensuremath{<}0.001}  % p-value
\newcommand{\rhabEnetAddRmCloseQlike}{0.1223}  % QLIKE enet add_rm c16
\newcommand{\rhabEnetAddRmCloseMse}{2.92}  % MSE x 1e10 enet add_rm c16
\newcommand{\rhabEnetAddRmClosePct}{$-$3.6}  % QLIKE change % vs har
\newcommand{\rhabEnetAddRmCloseDm}{$-$2.25}  % DM vs the reference arm
\newcommand{\rhabEnetAddRmCloseP}{0.024}  % p-value
\newcommand{\rhabEnetAddRmCloseLateQlike}{0.1147}  % QLIKE enet add_rm c16_2020
\newcommand{\rhabEnetAddRmCloseLateMse}{6.56}  % MSE x 1e10 enet add_rm c16_2020
\newcommand{\rhabEnetAddRmCloseLatePct}{1.0}  % QLIKE change % vs har
\newcommand{\rhabEnetAddRmCloseLateDm}{0.33}  % DM vs the reference arm
\newcommand{\rhabEnetAddRmCloseLateP}{0.741}  % p-value
\newcommand{\rhabEnetAddSentAllQlike}{0.1893}  % QLIKE enet add_sent all
\newcommand{\rhabEnetAddSentAllMse}{1.11}  % MSE x 1e10 enet add_sent all
\newcommand{\rhabEnetAddSentAllPct}{0.0}  % QLIKE change % vs har
\newcommand{\rhabEnetAddSentAllDm}{0.04}  % DM vs the reference arm
\newcommand{\rhabEnetAddSentAllP}{0.971}  % p-value
\newcommand{\rhabEnetAddSentCloseQlike}{0.1300}  % QLIKE enet add_sent c16
\newcommand{\rhabEnetAddSentCloseMse}{1.98}  % MSE x 1e10 enet add_sent c16
\newcommand{\rhabEnetAddSentClosePct}{2.5}  % QLIKE change % vs har
\newcommand{\rhabEnetAddSentCloseDm}{3.62}  % DM vs the reference arm
\newcommand{\rhabEnetAddSentCloseP}{\ensuremath{<}0.001}  % p-value
\newcommand{\rhabEnetAddSentCloseLateQlike}{0.1194}  % QLIKE enet add_sent c16_2020
\newcommand{\rhabEnetAddSentCloseLateMse}{4.27}  % MSE x 1e10 enet add_sent c16_2020
\newcommand{\rhabEnetAddSentCloseLatePct}{5.0}  % QLIKE change % vs har
\newcommand{\rhabEnetAddSentCloseLateDm}{2.88}  % DM vs the reference arm
\newcommand{\rhabEnetAddSentCloseLateP}{0.004}  % p-value
\newcommand{\rhabEnetAddVdAllQlike}{0.1900}  % QLIKE enet add_vd all
\newcommand{\rhabEnetAddVdAllMse}{1.11}  % MSE x 1e10 enet add_vd all
\newcommand{\rhabEnetAddVdAllPct}{0.4}  % QLIKE change % vs har
\newcommand{\rhabEnetAddVdAllDm}{2.92}  % DM vs the reference arm
\newcommand{\rhabEnetAddVdAllP}{0.003}  % p-value
\newcommand{\rhabEnetAddVdCloseQlike}{0.1277}  % QLIKE enet add_vd c16
\newcommand{\rhabEnetAddVdCloseMse}{1.96}  % MSE x 1e10 enet add_vd c16
\newcommand{\rhabEnetAddVdClosePct}{0.6}  % QLIKE change % vs har
\newcommand{\rhabEnetAddVdCloseDm}{1.85}  % DM vs the reference arm
\newcommand{\rhabEnetAddVdCloseP}{0.065}  % p-value
\newcommand{\rhabEnetAddVdCloseLateQlike}{0.1160}  % QLIKE enet add_vd c16_2020
\newcommand{\rhabEnetAddVdCloseLateMse}{4.21}  % MSE x 1e10 enet add_vd c16_2020
\newcommand{\rhabEnetAddVdCloseLatePct}{2.1}  % QLIKE change % vs har
\newcommand{\rhabEnetAddVdCloseLateDm}{2.57}  % DM vs the reference arm
\newcommand{\rhabEnetAddVdCloseLateP}{0.010}  % p-value
\newcommand{\rhabEnetAddVwAllQlike}{0.1872}  % QLIKE enet add_vw all
\newcommand{\rhabEnetAddVwAllMse}{1.99}  % MSE x 1e10 enet add_vw all
\newcommand{\rhabEnetAddVwAllPct}{$-$1.1}  % QLIKE change % vs har
\newcommand{\rhabEnetAddVwAllDm}{$-$2.03}  % DM vs the reference arm
\newcommand{\rhabEnetAddVwAllP}{0.042}  % p-value
\newcommand{\rhabEnetAddVwCloseQlike}{0.1256}  % QLIKE enet add_vw c16
\newcommand{\rhabEnetAddVwCloseMse}{2.37}  % MSE x 1e10 enet add_vw c16
\newcommand{\rhabEnetAddVwClosePct}{$-$1.0}  % QLIKE change % vs har
\newcommand{\rhabEnetAddVwCloseDm}{$-$0.75}  % DM vs the reference arm
\newcommand{\rhabEnetAddVwCloseP}{0.454}  % p-value
\newcommand{\rhabEnetAddVwCloseLateQlike}{0.1134}  % QLIKE enet add_vw c16_2020
\newcommand{\rhabEnetAddVwCloseLateMse}{5.22}  % MSE x 1e10 enet add_vw c16_2020
\newcommand{\rhabEnetAddVwCloseLatePct}{$-$0.2}  % QLIKE change % vs har
\newcommand{\rhabEnetAddVwCloseLateDm}{$-$0.13}  % DM vs the reference arm
\newcommand{\rhabEnetAddVwCloseLateP}{0.897}  % p-value
\newcommand{\rhabEnetFullAllQlike}{0.1778}  % QLIKE enet full all
\newcommand{\rhabEnetFullAllMse}{2.43}  % MSE x 1e10 enet full all
\newcommand{\rhabEnetFullAllPct}{$-$6.0}  % QLIKE change % vs har
\newcommand{\rhabEnetFullAllDm}{$-$6.98}  % DM vs the reference arm
\newcommand{\rhabEnetFullAllP}{\ensuremath{<}0.001}  % p-value
\newcommand{\rhabEnetFullCloseQlike}{0.1262}  % QLIKE enet full c16
\newcommand{\rhabEnetFullCloseMse}{3.94}  % MSE x 1e10 enet full c16
\newcommand{\rhabEnetFullClosePct}{$-$0.5}  % QLIKE change % vs har
\newcommand{\rhabEnetFullCloseDm}{$-$0.21}  % DM vs the reference arm
\newcommand{\rhabEnetFullCloseP}{0.832}  % p-value
\newcommand{\rhabEnetFullCloseLateQlike}{0.1234}  % QLIKE enet full c16_2020
\newcommand{\rhabEnetFullCloseLateMse}{9.01}  % MSE x 1e10 enet full c16_2020
\newcommand{\rhabEnetFullCloseLatePct}{8.5}  % QLIKE change % vs har
\newcommand{\rhabEnetFullCloseLateDm}{1.91}  % DM vs the reference arm
\newcommand{\rhabEnetFullCloseLateP}{0.056}  % p-value
\newcommand{\rhabEnetHarAllQlike}{0.1893}  % QLIKE enet har all
\newcommand{\rhabEnetHarAllMse}{1.11}  % MSE x 1e10 enet har all
\newcommand{\rhabEnetHarCloseQlike}{0.1269}  % QLIKE enet har c16
\newcommand{\rhabEnetHarCloseMse}{1.97}  % MSE x 1e10 enet har c16
\newcommand{\rhabEnetHarCloseLateQlike}{0.1137}  % QLIKE enet har c16_2020
\newcommand{\rhabEnetHarCloseLateMse}{4.24}  % MSE x 1e10 enet har c16_2020
\newcommand{\rhabEnetHarqAllQlike}{0.1893}  % QLIKE enet harq all
\newcommand{\rhabEnetHarqAllMse}{1.16}  % MSE x 1e10 enet harq all
\newcommand{\rhabEnetHarqAllPct}{0.0}  % QLIKE change % vs har
\newcommand{\rhabEnetHarqAllDm}{0.63}  % DM vs the reference arm
\newcommand{\rhabEnetHarqAllP}{0.531}  % p-value
\newcommand{\rhabEnetHarqCloseQlike}{0.1267}  % QLIKE enet harq c16
\newcommand{\rhabEnetHarqCloseMse}{1.80}  % MSE x 1e10 enet harq c16
\newcommand{\rhabEnetHarqClosePct}{$-$0.1}  % QLIKE change % vs har
\newcommand{\rhabEnetHarqCloseDm}{$-$1.31}  % DM vs the reference arm
\newcommand{\rhabEnetHarqCloseP}{0.191}  % p-value
\newcommand{\rhabEnetHarqCloseLateQlike}{0.1134}  % QLIKE enet harq c16_2020
\newcommand{\rhabEnetHarqCloseLateMse}{3.84}  % MSE x 1e10 enet harq c16_2020
\newcommand{\rhabEnetHarqCloseLatePct}{$-$0.3}  % QLIKE change % vs har
\newcommand{\rhabEnetHarqCloseLateDm}{$-$1.15}  % DM vs the reference arm
\newcommand{\rhabEnetHarqCloseLateP}{0.251}  % p-value
\newcommand{\rhabEnetLharAllQlike}{0.1841}  % QLIKE enet lhar all
\newcommand{\rhabEnetLharAllMse}{1.71}  % MSE x 1e10 enet lhar all
\newcommand{\rhabEnetLharAllPct}{$-$2.7}  % QLIKE change % vs har
\newcommand{\rhabEnetLharAllDm}{$-$6.66}  % DM vs the reference arm
\newcommand{\rhabEnetLharAllP}{\ensuremath{<}0.001}  % p-value
\newcommand{\rhabEnetLharCloseQlike}{0.1217}  % QLIKE enet lhar c16
\newcommand{\rhabEnetLharCloseMse}{2.76}  % MSE x 1e10 enet lhar c16
\newcommand{\rhabEnetLharClosePct}{$-$4.1}  % QLIKE change % vs har
\newcommand{\rhabEnetLharCloseDm}{$-$3.16}  % DM vs the reference arm
\newcommand{\rhabEnetLharCloseP}{0.002}  % p-value
\newcommand{\rhabEnetLharCloseLateQlike}{0.1124}  % QLIKE enet lhar c16_2020
\newcommand{\rhabEnetLharCloseLateMse}{6.19}  % MSE x 1e10 enet lhar c16_2020
\newcommand{\rhabEnetLharCloseLatePct}{$-$1.1}  % QLIKE change % vs har
\newcommand{\rhabEnetLharCloseLateDm}{$-$0.45}  % DM vs the reference arm
\newcommand{\rhabEnetLharCloseLateP}{0.656}  % p-value
\newcommand{\rhabEnetLoEwAllQlike}{0.1776}  % QLIKE enet lo_ew all
\newcommand{\rhabEnetLoEwAllMse}{2.54}  % MSE x 1e10 enet lo_ew all
\newcommand{\rhabEnetLoEwAllPct}{$-$0.1}  % QLIKE change % vs full
\newcommand{\rhabEnetLoEwAllDm}{$-$1.41}  % DM vs the reference arm
\newcommand{\rhabEnetLoEwAllP}{0.157}  % p-value
\newcommand{\rhabEnetLoEwCloseQlike}{0.1267}  % QLIKE enet lo_ew c16
\newcommand{\rhabEnetLoEwCloseMse}{4.02}  % MSE x 1e10 enet lo_ew c16
\newcommand{\rhabEnetLoEwClosePct}{0.3}  % QLIKE change % vs full
\newcommand{\rhabEnetLoEwCloseDm}{1.18}  % DM vs the reference arm
\newcommand{\rhabEnetLoEwCloseP}{0.238}  % p-value
\newcommand{\rhabEnetLoEwCloseLateQlike}{0.1235}  % QLIKE enet lo_ew c16_2020
\newcommand{\rhabEnetLoEwCloseLateMse}{9.21}  % MSE x 1e10 enet lo_ew c16_2020
\newcommand{\rhabEnetLoEwCloseLatePct}{0.1}  % QLIKE change % vs full
\newcommand{\rhabEnetLoEwCloseLateDm}{0.23}  % DM vs the reference arm
\newcommand{\rhabEnetLoEwCloseLateP}{0.822}  % p-value
\newcommand{\rhabEnetLoFomcAllQlike}{0.1820}  % QLIKE enet lo_fomc all
\newcommand{\rhabEnetLoFomcAllMse}{2.42}  % MSE x 1e10 enet lo_fomc all
\newcommand{\rhabEnetLoFomcAllPct}{2.3}  % QLIKE change % vs full
\newcommand{\rhabEnetLoFomcAllDm}{5.74}  % DM vs the reference arm
\newcommand{\rhabEnetLoFomcAllP}{\ensuremath{<}0.001}  % p-value
\newcommand{\rhabEnetLoFomcCloseQlike}{0.1261}  % QLIKE enet lo_fomc c16
\newcommand{\rhabEnetLoFomcCloseMse}{3.99}  % MSE x 1e10 enet lo_fomc c16
\newcommand{\rhabEnetLoFomcClosePct}{$-$0.1}  % QLIKE change % vs full
\newcommand{\rhabEnetLoFomcCloseDm}{$-$0.51}  % DM vs the reference arm
\newcommand{\rhabEnetLoFomcCloseP}{0.607}  % p-value
\newcommand{\rhabEnetLoFomcCloseLateQlike}{0.1234}  % QLIKE enet lo_fomc c16_2020
\newcommand{\rhabEnetLoFomcCloseLateMse}{9.13}  % MSE x 1e10 enet lo_fomc c16_2020
\newcommand{\rhabEnetLoFomcCloseLatePct}{0.1}  % QLIKE change % vs full
\newcommand{\rhabEnetLoFomcCloseLateDm}{0.13}  % DM vs the reference arm
\newcommand{\rhabEnetLoFomcCloseLateP}{0.894}  % p-value
\newcommand{\rhabEnetLoIvAllQlike}{0.1778}  % QLIKE enet lo_iv all
\newcommand{\rhabEnetLoIvAllMse}{2.49}  % MSE x 1e10 enet lo_iv all
\newcommand{\rhabEnetLoIvAllPct}{0.0}  % QLIKE change % vs full
\newcommand{\rhabEnetLoIvAllDm}{$-$0.13}  % DM vs the reference arm
\newcommand{\rhabEnetLoIvAllP}{0.898}  % p-value
\newcommand{\rhabEnetLoIvCloseQlike}{0.1243}  % QLIKE enet lo_iv c16
\newcommand{\rhabEnetLoIvCloseMse}{4.03}  % MSE x 1e10 enet lo_iv c16
\newcommand{\rhabEnetLoIvClosePct}{$-$1.5}  % QLIKE change % vs full
\newcommand{\rhabEnetLoIvCloseDm}{$-$1.51}  % DM vs the reference arm
\newcommand{\rhabEnetLoIvCloseP}{0.132}  % p-value
\newcommand{\rhabEnetLoIvCloseLateQlike}{0.1193}  % QLIKE enet lo_iv c16_2020
\newcommand{\rhabEnetLoIvCloseLateMse}{9.21}  % MSE x 1e10 enet lo_iv c16_2020
\newcommand{\rhabEnetLoIvCloseLatePct}{$-$3.3}  % QLIKE change % vs full
\newcommand{\rhabEnetLoIvCloseLateDm}{$-$1.72}  % DM vs the reference arm
\newcommand{\rhabEnetLoIvCloseLateP}{0.086}  % p-value
\newcommand{\rhabEnetLoLiqAllQlike}{0.1779}  % QLIKE enet lo_liq all
\newcommand{\rhabEnetLoLiqAllMse}{2.43}  % MSE x 1e10 enet lo_liq all
\newcommand{\rhabEnetLoLiqAllPct}{0.1}  % QLIKE change % vs full
\newcommand{\rhabEnetLoLiqAllDm}{0.35}  % DM vs the reference arm
\newcommand{\rhabEnetLoLiqAllP}{0.724}  % p-value
\newcommand{\rhabEnetLoLiqCloseQlike}{0.1251}  % QLIKE enet lo_liq c16
\newcommand{\rhabEnetLoLiqCloseMse}{4.29}  % MSE x 1e10 enet lo_liq c16
\newcommand{\rhabEnetLoLiqClosePct}{$-$0.9}  % QLIKE change % vs full
\newcommand{\rhabEnetLoLiqCloseDm}{$-$1.54}  % DM vs the reference arm
\newcommand{\rhabEnetLoLiqCloseP}{0.123}  % p-value
\newcommand{\rhabEnetLoLiqCloseLateQlike}{0.1223}  % QLIKE enet lo_liq c16_2020
\newcommand{\rhabEnetLoLiqCloseLateMse}{9.84}  % MSE x 1e10 enet lo_liq c16_2020
\newcommand{\rhabEnetLoLiqCloseLatePct}{$-$0.9}  % QLIKE change % vs full
\newcommand{\rhabEnetLoLiqCloseLateDm}{$-$0.92}  % DM vs the reference arm
\newcommand{\rhabEnetLoLiqCloseLateP}{0.358}  % p-value
\newcommand{\rhabEnetLoRmAllQlike}{0.1813}  % QLIKE enet lo_rm all
\newcommand{\rhabEnetLoRmAllMse}{2.00}  % MSE x 1e10 enet lo_rm all
\newcommand{\rhabEnetLoRmAllPct}{1.9}  % QLIKE change % vs full
\newcommand{\rhabEnetLoRmAllDm}{4.38}  % DM vs the reference arm
\newcommand{\rhabEnetLoRmAllP}{\ensuremath{<}0.001}  % p-value
\newcommand{\rhabEnetLoRmCloseQlike}{0.1285}  % QLIKE enet lo_rm c16
\newcommand{\rhabEnetLoRmCloseMse}{2.35}  % MSE x 1e10 enet lo_rm c16
\newcommand{\rhabEnetLoRmClosePct}{1.8}  % QLIKE change % vs full
\newcommand{\rhabEnetLoRmCloseDm}{1.31}  % DM vs the reference arm
\newcommand{\rhabEnetLoRmCloseP}{0.190}  % p-value
\newcommand{\rhabEnetLoRmCloseLateQlike}{0.1217}  % QLIKE enet lo_rm c16_2020
\newcommand{\rhabEnetLoRmCloseLateMse}{5.18}  % MSE x 1e10 enet lo_rm c16_2020
\newcommand{\rhabEnetLoRmCloseLatePct}{$-$1.3}  % QLIKE change % vs full
\newcommand{\rhabEnetLoRmCloseLateDm}{$-$0.62}  % DM vs the reference arm
\newcommand{\rhabEnetLoRmCloseLateP}{0.538}  % p-value
\newcommand{\rhabEnetLoSentAllQlike}{0.1782}  % QLIKE enet lo_sent all
\newcommand{\rhabEnetLoSentAllMse}{2.43}  % MSE x 1e10 enet lo_sent all
\newcommand{\rhabEnetLoSentAllPct}{0.2}  % QLIKE change % vs full
\newcommand{\rhabEnetLoSentAllDm}{1.49}  % DM vs the reference arm
\newcommand{\rhabEnetLoSentAllP}{0.136}  % p-value
\newcommand{\rhabEnetLoSentCloseQlike}{0.1253}  % QLIKE enet lo_sent c16
\newcommand{\rhabEnetLoSentCloseMse}{3.94}  % MSE x 1e10 enet lo_sent c16
\newcommand{\rhabEnetLoSentClosePct}{$-$0.7}  % QLIKE change % vs full
\newcommand{\rhabEnetLoSentCloseDm}{$-$2.12}  % DM vs the reference arm
\newcommand{\rhabEnetLoSentCloseP}{0.034}  % p-value
\newcommand{\rhabEnetLoSentCloseLateQlike}{0.1212}  % QLIKE enet lo_sent c16_2020
\newcommand{\rhabEnetLoSentCloseLateMse}{9.01}  % MSE x 1e10 enet lo_sent c16_2020
\newcommand{\rhabEnetLoSentCloseLatePct}{$-$1.7}  % QLIKE change % vs full
\newcommand{\rhabEnetLoSentCloseLateDm}{$-$2.21}  % DM vs the reference arm
\newcommand{\rhabEnetLoSentCloseLateP}{0.027}  % p-value
\newcommand{\rhabEnetLoVdAllQlike}{0.1777}  % QLIKE enet lo_vd all
\newcommand{\rhabEnetLoVdAllMse}{2.43}  % MSE x 1e10 enet lo_vd all
\newcommand{\rhabEnetLoVdAllPct}{$-$0.1}  % QLIKE change % vs full
\newcommand{\rhabEnetLoVdAllDm}{$-$1.13}  % DM vs the reference arm
\newcommand{\rhabEnetLoVdAllP}{0.260}  % p-value
\newcommand{\rhabEnetLoVdCloseQlike}{0.1252}  % QLIKE enet lo_vd c16
\newcommand{\rhabEnetLoVdCloseMse}{3.94}  % MSE x 1e10 enet lo_vd c16
\newcommand{\rhabEnetLoVdClosePct}{$-$0.8}  % QLIKE change % vs full
\newcommand{\rhabEnetLoVdCloseDm}{$-$1.94}  % DM vs the reference arm
\newcommand{\rhabEnetLoVdCloseP}{0.053}  % p-value
\newcommand{\rhabEnetLoVdCloseLateQlike}{0.1207}  % QLIKE enet lo_vd c16_2020
\newcommand{\rhabEnetLoVdCloseLateMse}{9.01}  % MSE x 1e10 enet lo_vd c16_2020
\newcommand{\rhabEnetLoVdCloseLatePct}{$-$2.1}  % QLIKE change % vs full
\newcommand{\rhabEnetLoVdCloseLateDm}{$-$2.18}  % DM vs the reference arm
\newcommand{\rhabEnetLoVdCloseLateP}{0.029}  % p-value
\newcommand{\rhabEnetLoVwAllQlike}{0.1777}  % QLIKE enet lo_vw all
\newcommand{\rhabEnetLoVwAllMse}{2.08}  % MSE x 1e10 enet lo_vw all
\newcommand{\rhabEnetLoVwAllPct}{$-$0.1}  % QLIKE change % vs full
\newcommand{\rhabEnetLoVwAllDm}{$-$0.65}  % DM vs the reference arm
\newcommand{\rhabEnetLoVwAllP}{0.516}  % p-value
\newcommand{\rhabEnetLoVwCloseQlike}{0.1254}  % QLIKE enet lo_vw c16
\newcommand{\rhabEnetLoVwCloseMse}{3.77}  % MSE x 1e10 enet lo_vw c16
\newcommand{\rhabEnetLoVwClosePct}{$-$0.7}  % QLIKE change % vs full
\newcommand{\rhabEnetLoVwCloseDm}{$-$1.66}  % DM vs the reference arm
\newcommand{\rhabEnetLoVwCloseP}{0.097}  % p-value
\newcommand{\rhabEnetLoVwCloseLateQlike}{0.1218}  % QLIKE enet lo_vw c16_2020
\newcommand{\rhabEnetLoVwCloseLateMse}{8.59}  % MSE x 1e10 enet lo_vw c16_2020
\newcommand{\rhabEnetLoVwCloseLatePct}{$-$1.3}  % QLIKE change % vs full
\newcommand{\rhabEnetLoVwCloseLateDm}{$-$2.65}  % DM vs the reference arm
\newcommand{\rhabEnetLoVwCloseLateP}{0.008}  % p-value
\newcommand{\rhabEnetSharAllQlike}{0.1877}  % QLIKE enet shar all
\newcommand{\rhabEnetSharAllMse}{1.38}  % MSE x 1e10 enet shar all
\newcommand{\rhabEnetSharAllPct}{$-$0.9}  % QLIKE change % vs har
\newcommand{\rhabEnetSharAllDm}{$-$3.68}  % DM vs the reference arm
\newcommand{\rhabEnetSharAllP}{\ensuremath{<}0.001}  % p-value
\newcommand{\rhabEnetSharCloseQlike}{0.1265}  % QLIKE enet shar c16
\newcommand{\rhabEnetSharCloseMse}{1.83}  % MSE x 1e10 enet shar c16
\newcommand{\rhabEnetSharClosePct}{$-$0.3}  % QLIKE change % vs har
\newcommand{\rhabEnetSharCloseDm}{$-$0.54}  % DM vs the reference arm
\newcommand{\rhabEnetSharCloseP}{0.586}  % p-value
\newcommand{\rhabEnetSharCloseLateQlike}{0.1141}  % QLIKE enet shar c16_2020
\newcommand{\rhabEnetSharCloseLateMse}{3.92}  % MSE x 1e10 enet shar c16_2020
\newcommand{\rhabEnetSharCloseLatePct}{0.3}  % QLIKE change % vs har
\newcommand{\rhabEnetSharCloseLateDm}{0.34}  % DM vs the reference arm
\newcommand{\rhabEnetSharCloseLateP}{0.731}  % p-value
\newcommand{\rhabEnetSharqlAllQlike}{0.1840}  % QLIKE enet sharql all
\newcommand{\rhabEnetSharqlAllMse}{1.69}  % MSE x 1e10 enet sharql all
\newcommand{\rhabEnetSharqlAllPct}{$-$2.8}  % QLIKE change % vs har
\newcommand{\rhabEnetSharqlAllDm}{$-$6.95}  % DM vs the reference arm
\newcommand{\rhabEnetSharqlAllP}{\ensuremath{<}0.001}  % p-value
\newcommand{\rhabEnetSharqlCloseQlike}{0.1216}  % QLIKE enet sharql c16
\newcommand{\rhabEnetSharqlCloseMse}{2.64}  % MSE x 1e10 enet sharql c16
\newcommand{\rhabEnetSharqlClosePct}{$-$4.1}  % QLIKE change % vs har
\newcommand{\rhabEnetSharqlCloseDm}{$-$3.46}  % DM vs the reference arm
\newcommand{\rhabEnetSharqlCloseP}{\ensuremath{<}0.001}  % p-value
\newcommand{\rhabEnetSharqlCloseLateQlike}{0.1121}  % QLIKE enet sharql c16_2020
\newcommand{\rhabEnetSharqlCloseLateMse}{5.90}  % MSE x 1e10 enet sharql c16_2020
\newcommand{\rhabEnetSharqlCloseLatePct}{$-$1.4}  % QLIKE change % vs har
\newcommand{\rhabEnetSharqlCloseLateDm}{$-$0.64}  % DM vs the reference arm
\newcommand{\rhabEnetSharqlCloseLateP}{0.524}  % p-value
\newcommand{\rhabLassoAddEwAllQlike}{0.1877}  % QLIKE lasso add_ew all
\newcommand{\rhabLassoAddEwAllMse}{1.47}  % MSE x 1e10 lasso add_ew all
\newcommand{\rhabLassoAddEwAllPct}{$-$0.8}  % QLIKE change % vs har
\newcommand{\rhabLassoAddEwAllDm}{$-$2.01}  % DM vs the reference arm
\newcommand{\rhabLassoAddEwAllP}{0.045}  % p-value
\newcommand{\rhabLassoAddEwCloseQlike}{0.1257}  % QLIKE lasso add_ew c16
\newcommand{\rhabLassoAddEwCloseMse}{2.13}  % MSE x 1e10 lasso add_ew c16
\newcommand{\rhabLassoAddEwClosePct}{$-$1.0}  % QLIKE change % vs har
\newcommand{\rhabLassoAddEwCloseDm}{$-$0.81}  % DM vs the reference arm
\newcommand{\rhabLassoAddEwCloseP}{0.417}  % p-value
\newcommand{\rhabLassoAddEwCloseLateQlike}{0.1123}  % QLIKE lasso add_ew c16_2020
\newcommand{\rhabLassoAddEwCloseLateMse}{4.65}  % MSE x 1e10 lasso add_ew c16_2020
\newcommand{\rhabLassoAddEwCloseLatePct}{$-$1.2}  % QLIKE change % vs har
\newcommand{\rhabLassoAddEwCloseLateDm}{$-$0.70}  % DM vs the reference arm
\newcommand{\rhabLassoAddEwCloseLateP}{0.485}  % p-value
\newcommand{\rhabLassoAddFomcAllQlike}{0.1820}  % QLIKE lasso add_fomc all
\newcommand{\rhabLassoAddFomcAllMse}{1.13}  % MSE x 1e10 lasso add_fomc all
\newcommand{\rhabLassoAddFomcAllPct}{$-$3.8}  % QLIKE change % vs har
\newcommand{\rhabLassoAddFomcAllDm}{$-$5.76}  % DM vs the reference arm
\newcommand{\rhabLassoAddFomcAllP}{\ensuremath{<}0.001}  % p-value
\newcommand{\rhabLassoAddFomcCloseQlike}{0.1296}  % QLIKE lasso add_fomc c16
\newcommand{\rhabLassoAddFomcCloseMse}{1.83}  % MSE x 1e10 lasso add_fomc c16
\newcommand{\rhabLassoAddFomcClosePct}{2.2}  % QLIKE change % vs har
\newcommand{\rhabLassoAddFomcCloseDm}{4.33}  % DM vs the reference arm
\newcommand{\rhabLassoAddFomcCloseP}{\ensuremath{<}0.001}  % p-value
\newcommand{\rhabLassoAddFomcCloseLateQlike}{0.1135}  % QLIKE lasso add_fomc c16_2020
\newcommand{\rhabLassoAddFomcCloseLateMse}{3.92}  % MSE x 1e10 lasso add_fomc c16_2020
\newcommand{\rhabLassoAddFomcCloseLatePct}{$-$0.1}  % QLIKE change % vs har
\newcommand{\rhabLassoAddFomcCloseLateDm}{$-$0.16}  % DM vs the reference arm
\newcommand{\rhabLassoAddFomcCloseLateP}{0.875}  % p-value
\newcommand{\rhabLassoAddIvAllQlike}{0.1887}  % QLIKE lasso add_iv all
\newcommand{\rhabLassoAddIvAllMse}{1.91}  % MSE x 1e10 lasso add_iv all
\newcommand{\rhabLassoAddIvAllPct}{$-$0.3}  % QLIKE change % vs har
\newcommand{\rhabLassoAddIvAllDm}{$-$0.79}  % DM vs the reference arm
\newcommand{\rhabLassoAddIvAllP}{0.430}  % p-value
\newcommand{\rhabLassoAddIvCloseQlike}{0.1263}  % QLIKE lasso add_iv c16
\newcommand{\rhabLassoAddIvCloseMse}{1.82}  % MSE x 1e10 lasso add_iv c16
\newcommand{\rhabLassoAddIvClosePct}{$-$0.4}  % QLIKE change % vs har
\newcommand{\rhabLassoAddIvCloseDm}{$-$0.38}  % DM vs the reference arm
\newcommand{\rhabLassoAddIvCloseP}{0.705}  % p-value
\newcommand{\rhabLassoAddIvCloseLateQlike}{0.1144}  % QLIKE lasso add_iv c16_2020
\newcommand{\rhabLassoAddIvCloseLateMse}{3.92}  % MSE x 1e10 lasso add_iv c16_2020
\newcommand{\rhabLassoAddIvCloseLatePct}{0.6}  % QLIKE change % vs har
\newcommand{\rhabLassoAddIvCloseLateDm}{0.29}  % DM vs the reference arm
\newcommand{\rhabLassoAddIvCloseLateP}{0.768}  % p-value
\newcommand{\rhabLassoAddLiqAllQlike}{0.1888}  % QLIKE lasso add_liq all
\newcommand{\rhabLassoAddLiqAllMse}{1.17}  % MSE x 1e10 lasso add_liq all
\newcommand{\rhabLassoAddLiqAllPct}{$-$0.3}  % QLIKE change % vs har
\newcommand{\rhabLassoAddLiqAllDm}{$-$0.55}  % DM vs the reference arm
\newcommand{\rhabLassoAddLiqAllP}{0.583}  % p-value
\newcommand{\rhabLassoAddLiqCloseQlike}{0.1305}  % QLIKE lasso add_liq c16
\newcommand{\rhabLassoAddLiqCloseMse}{2.02}  % MSE x 1e10 lasso add_liq c16
\newcommand{\rhabLassoAddLiqClosePct}{2.9}  % QLIKE change % vs har
\newcommand{\rhabLassoAddLiqCloseDm}{1.86}  % DM vs the reference arm
\newcommand{\rhabLassoAddLiqCloseP}{0.063}  % p-value
\newcommand{\rhabLassoAddLiqCloseLateQlike}{0.1206}  % QLIKE lasso add_liq c16_2020
\newcommand{\rhabLassoAddLiqCloseLateMse}{4.37}  % MSE x 1e10 lasso add_liq c16_2020
\newcommand{\rhabLassoAddLiqCloseLatePct}{6.1}  % QLIKE change % vs har
\newcommand{\rhabLassoAddLiqCloseLateDm}{2.00}  % DM vs the reference arm
\newcommand{\rhabLassoAddLiqCloseLateP}{0.045}  % p-value
\newcommand{\rhabLassoAddRmAllQlike}{0.1852}  % QLIKE lasso add_rm all
\newcommand{\rhabLassoAddRmAllMse}{1.81}  % MSE x 1e10 lasso add_rm all
\newcommand{\rhabLassoAddRmAllPct}{$-$2.1}  % QLIKE change % vs har
\newcommand{\rhabLassoAddRmAllDm}{$-$2.64}  % DM vs the reference arm
\newcommand{\rhabLassoAddRmAllP}{0.008}  % p-value
\newcommand{\rhabLassoAddRmCloseQlike}{0.1234}  % QLIKE lasso add_rm c16
\newcommand{\rhabLassoAddRmCloseMse}{3.11}  % MSE x 1e10 lasso add_rm c16
\newcommand{\rhabLassoAddRmClosePct}{$-$2.8}  % QLIKE change % vs har
\newcommand{\rhabLassoAddRmCloseDm}{$-$1.41}  % DM vs the reference arm
\newcommand{\rhabLassoAddRmCloseP}{0.160}  % p-value
\newcommand{\rhabLassoAddRmCloseLateQlike}{0.1176}  % QLIKE lasso add_rm c16_2020
\newcommand{\rhabLassoAddRmCloseLateMse}{7.01}  % MSE x 1e10 lasso add_rm c16_2020
\newcommand{\rhabLassoAddRmCloseLatePct}{3.5}  % QLIKE change % vs har
\newcommand{\rhabLassoAddRmCloseLateDm}{0.90}  % DM vs the reference arm
\newcommand{\rhabLassoAddRmCloseLateP}{0.368}  % p-value
\newcommand{\rhabLassoAddSentAllQlike}{0.1895}  % QLIKE lasso add_sent all
\newcommand{\rhabLassoAddSentAllMse}{1.11}  % MSE x 1e10 lasso add_sent all
\newcommand{\rhabLassoAddSentAllPct}{0.1}  % QLIKE change % vs har
\newcommand{\rhabLassoAddSentAllDm}{0.60}  % DM vs the reference arm
\newcommand{\rhabLassoAddSentAllP}{0.550}  % p-value
\newcommand{\rhabLassoAddSentCloseQlike}{0.1310}  % QLIKE lasso add_sent c16
\newcommand{\rhabLassoAddSentCloseMse}{1.98}  % MSE x 1e10 lasso add_sent c16
\newcommand{\rhabLassoAddSentClosePct}{3.2}  % QLIKE change % vs har
\newcommand{\rhabLassoAddSentCloseDm}{4.34}  % DM vs the reference arm
\newcommand{\rhabLassoAddSentCloseP}{\ensuremath{<}0.001}  % p-value
\newcommand{\rhabLassoAddSentCloseLateQlike}{0.1214}  % QLIKE lasso add_sent c16_2020
\newcommand{\rhabLassoAddSentCloseLateMse}{4.27}  % MSE x 1e10 lasso add_sent c16_2020
\newcommand{\rhabLassoAddSentCloseLatePct}{6.8}  % QLIKE change % vs har
\newcommand{\rhabLassoAddSentCloseLateDm}{3.64}  % DM vs the reference arm
\newcommand{\rhabLassoAddSentCloseLateP}{\ensuremath{<}0.001}  % p-value
\newcommand{\rhabLassoAddVdAllQlike}{0.1903}  % QLIKE lasso add_vd all
\newcommand{\rhabLassoAddVdAllMse}{1.12}  % MSE x 1e10 lasso add_vd all
\newcommand{\rhabLassoAddVdAllPct}{0.6}  % QLIKE change % vs har
\newcommand{\rhabLassoAddVdAllDm}{3.05}  % DM vs the reference arm
\newcommand{\rhabLassoAddVdAllP}{0.002}  % p-value
\newcommand{\rhabLassoAddVdCloseQlike}{0.1272}  % QLIKE lasso add_vd c16
\newcommand{\rhabLassoAddVdCloseMse}{1.95}  % MSE x 1e10 lasso add_vd c16
\newcommand{\rhabLassoAddVdClosePct}{0.3}  % QLIKE change % vs har
\newcommand{\rhabLassoAddVdCloseDm}{0.65}  % DM vs the reference arm
\newcommand{\rhabLassoAddVdCloseP}{0.516}  % p-value
\newcommand{\rhabLassoAddVdCloseLateQlike}{0.1169}  % QLIKE lasso add_vd c16_2020
\newcommand{\rhabLassoAddVdCloseLateMse}{4.20}  % MSE x 1e10 lasso add_vd c16_2020
\newcommand{\rhabLassoAddVdCloseLatePct}{2.9}  % QLIKE change % vs har
\newcommand{\rhabLassoAddVdCloseLateDm}{3.03}  % DM vs the reference arm
\newcommand{\rhabLassoAddVdCloseLateP}{0.002}  % p-value
\newcommand{\rhabLassoAddVwAllQlike}{0.1881}  % QLIKE lasso add_vw all
\newcommand{\rhabLassoAddVwAllMse}{2.04}  % MSE x 1e10 lasso add_vw all
\newcommand{\rhabLassoAddVwAllPct}{$-$0.6}  % QLIKE change % vs har
\newcommand{\rhabLassoAddVwAllDm}{$-$1.02}  % DM vs the reference arm
\newcommand{\rhabLassoAddVwAllP}{0.309}  % p-value
\newcommand{\rhabLassoAddVwCloseQlike}{0.1262}  % QLIKE lasso add_vw c16
\newcommand{\rhabLassoAddVwCloseMse}{2.38}  % MSE x 1e10 lasso add_vw c16
\newcommand{\rhabLassoAddVwClosePct}{$-$0.5}  % QLIKE change % vs har
\newcommand{\rhabLassoAddVwCloseDm}{$-$0.35}  % DM vs the reference arm
\newcommand{\rhabLassoAddVwCloseP}{0.726}  % p-value
\newcommand{\rhabLassoAddVwCloseLateQlike}{0.1136}  % QLIKE lasso add_vw c16_2020
\newcommand{\rhabLassoAddVwCloseLateMse}{5.24}  % MSE x 1e10 lasso add_vw c16_2020
\newcommand{\rhabLassoAddVwCloseLatePct}{$-$0.1}  % QLIKE change % vs har
\newcommand{\rhabLassoAddVwCloseLateDm}{$-$0.04}  % DM vs the reference arm
\newcommand{\rhabLassoAddVwCloseLateP}{0.968}  % p-value
\newcommand{\rhabLassoFullAllQlike}{0.1775}  % QLIKE lasso full all
\newcommand{\rhabLassoFullAllMse}{2.37}  % MSE x 1e10 lasso full all
\newcommand{\rhabLassoFullAllPct}{$-$6.2}  % QLIKE change % vs har
\newcommand{\rhabLassoFullAllDm}{$-$4.72}  % DM vs the reference arm
\newcommand{\rhabLassoFullAllP}{\ensuremath{<}0.001}  % p-value
\newcommand{\rhabLassoFullCloseQlike}{0.1310}  % QLIKE lasso full c16
\newcommand{\rhabLassoFullCloseMse}{3.57}  % MSE x 1e10 lasso full c16
\newcommand{\rhabLassoFullClosePct}{3.3}  % QLIKE change % vs har
\newcommand{\rhabLassoFullCloseDm}{1.04}  % DM vs the reference arm
\newcommand{\rhabLassoFullCloseP}{0.300}  % p-value
\newcommand{\rhabLassoFullCloseLateQlike}{0.1301}  % QLIKE lasso full c16_2020
\newcommand{\rhabLassoFullCloseLateMse}{8.13}  % MSE x 1e10 lasso full c16_2020
\newcommand{\rhabLassoFullCloseLatePct}{14.4}  % QLIKE change % vs har
\newcommand{\rhabLassoFullCloseLateDm}{2.51}  % DM vs the reference arm
\newcommand{\rhabLassoFullCloseLateP}{0.012}  % p-value
\newcommand{\rhabLassoHarAllQlike}{0.1893}  % QLIKE lasso har all
\newcommand{\rhabLassoHarAllMse}{1.11}  % MSE x 1e10 lasso har all
\newcommand{\rhabLassoHarCloseQlike}{0.1269}  % QLIKE lasso har c16
\newcommand{\rhabLassoHarCloseMse}{1.97}  % MSE x 1e10 lasso har c16
\newcommand{\rhabLassoHarCloseLateQlike}{0.1137}  % QLIKE lasso har c16_2020
\newcommand{\rhabLassoHarCloseLateMse}{4.24}  % MSE x 1e10 lasso har c16_2020
\newcommand{\rhabLassoHarqAllQlike}{0.1893}  % QLIKE lasso harq all
\newcommand{\rhabLassoHarqAllMse}{1.17}  % MSE x 1e10 lasso harq all
\newcommand{\rhabLassoHarqAllPct}{0.0}  % QLIKE change % vs har
\newcommand{\rhabLassoHarqAllDm}{0.61}  % DM vs the reference arm
\newcommand{\rhabLassoHarqAllP}{0.544}  % p-value
\newcommand{\rhabLassoHarqCloseQlike}{0.1267}  % QLIKE lasso harq c16
\newcommand{\rhabLassoHarqCloseMse}{1.78}  % MSE x 1e10 lasso harq c16
\newcommand{\rhabLassoHarqClosePct}{$-$0.2}  % QLIKE change % vs har
\newcommand{\rhabLassoHarqCloseDm}{$-$1.37}  % DM vs the reference arm
\newcommand{\rhabLassoHarqCloseP}{0.171}  % p-value
\newcommand{\rhabLassoHarqCloseLateQlike}{0.1133}  % QLIKE lasso harq c16_2020
\newcommand{\rhabLassoHarqCloseLateMse}{3.80}  % MSE x 1e10 lasso harq c16_2020
\newcommand{\rhabLassoHarqCloseLatePct}{$-$0.3}  % QLIKE change % vs har
\newcommand{\rhabLassoHarqCloseLateDm}{$-$1.18}  % DM vs the reference arm
\newcommand{\rhabLassoHarqCloseLateP}{0.238}  % p-value
\newcommand{\rhabLassoLharAllQlike}{0.1842}  % QLIKE lasso lhar all
\newcommand{\rhabLassoLharAllMse}{1.71}  % MSE x 1e10 lasso lhar all
\newcommand{\rhabLassoLharAllPct}{$-$2.7}  % QLIKE change % vs har
\newcommand{\rhabLassoLharAllDm}{$-$6.30}  % DM vs the reference arm
\newcommand{\rhabLassoLharAllP}{\ensuremath{<}0.001}  % p-value
\newcommand{\rhabLassoLharCloseQlike}{0.1219}  % QLIKE lasso lhar c16
\newcommand{\rhabLassoLharCloseMse}{2.80}  % MSE x 1e10 lasso lhar c16
\newcommand{\rhabLassoLharClosePct}{$-$3.9}  % QLIKE change % vs har
\newcommand{\rhabLassoLharCloseDm}{$-$2.77}  % DM vs the reference arm
\newcommand{\rhabLassoLharCloseP}{0.006}  % p-value
\newcommand{\rhabLassoLharCloseLateQlike}{0.1130}  % QLIKE lasso lhar c16_2020
\newcommand{\rhabLassoLharCloseLateMse}{6.28}  % MSE x 1e10 lasso lhar c16_2020
\newcommand{\rhabLassoLharCloseLatePct}{$-$0.6}  % QLIKE change % vs har
\newcommand{\rhabLassoLharCloseLateDm}{$-$0.21}  % DM vs the reference arm
\newcommand{\rhabLassoLharCloseLateP}{0.836}  % p-value
\newcommand{\rhabLassoLoEwAllQlike}{0.1770}  % QLIKE lasso lo_ew all
\newcommand{\rhabLassoLoEwAllMse}{2.54}  % MSE x 1e10 lasso lo_ew all
\newcommand{\rhabLassoLoEwAllPct}{$-$0.3}  % QLIKE change % vs full
\newcommand{\rhabLassoLoEwAllDm}{$-$1.46}  % DM vs the reference arm
\newcommand{\rhabLassoLoEwAllP}{0.144}  % p-value
\newcommand{\rhabLassoLoEwCloseQlike}{0.1310}  % QLIKE lasso lo_ew c16
\newcommand{\rhabLassoLoEwCloseMse}{3.82}  % MSE x 1e10 lasso lo_ew c16
\newcommand{\rhabLassoLoEwClosePct}{0.0}  % QLIKE change % vs full
\newcommand{\rhabLassoLoEwCloseDm}{0.00}  % DM vs the reference arm
\newcommand{\rhabLassoLoEwCloseP}{0.999}  % p-value
\newcommand{\rhabLassoLoEwCloseLateQlike}{0.1289}  % QLIKE lasso lo_ew c16_2020
\newcommand{\rhabLassoLoEwCloseLateMse}{8.72}  % MSE x 1e10 lasso lo_ew c16_2020
\newcommand{\rhabLassoLoEwCloseLatePct}{$-$0.9}  % QLIKE change % vs full
\newcommand{\rhabLassoLoEwCloseLateDm}{$-$0.76}  % DM vs the reference arm
\newcommand{\rhabLassoLoEwCloseLateP}{0.446}  % p-value
\newcommand{\rhabLassoLoFomcAllQlike}{0.1839}  % QLIKE lasso lo_fomc all
\newcommand{\rhabLassoLoFomcAllMse}{2.39}  % MSE x 1e10 lasso lo_fomc all
\newcommand{\rhabLassoLoFomcAllPct}{3.6}  % QLIKE change % vs full
\newcommand{\rhabLassoLoFomcAllDm}{5.10}  % DM vs the reference arm
\newcommand{\rhabLassoLoFomcAllP}{\ensuremath{<}0.001}  % p-value
\newcommand{\rhabLassoLoFomcCloseQlike}{0.1293}  % QLIKE lasso lo_fomc c16
\newcommand{\rhabLassoLoFomcCloseMse}{3.62}  % MSE x 1e10 lasso lo_fomc c16
\newcommand{\rhabLassoLoFomcClosePct}{$-$1.4}  % QLIKE change % vs full
\newcommand{\rhabLassoLoFomcCloseDm}{$-$3.21}  % DM vs the reference arm
\newcommand{\rhabLassoLoFomcCloseP}{0.001}  % p-value
\newcommand{\rhabLassoLoFomcCloseLateQlike}{0.1293}  % QLIKE lasso lo_fomc c16_2020
\newcommand{\rhabLassoLoFomcCloseLateMse}{8.25}  % MSE x 1e10 lasso lo_fomc c16_2020
\newcommand{\rhabLassoLoFomcCloseLatePct}{$-$0.6}  % QLIKE change % vs full
\newcommand{\rhabLassoLoFomcCloseLateDm}{$-$0.92}  % DM vs the reference arm
\newcommand{\rhabLassoLoFomcCloseLateP}{0.358}  % p-value
\newcommand{\rhabLassoLoIvAllQlike}{0.1778}  % QLIKE lasso lo_iv all
\newcommand{\rhabLassoLoIvAllMse}{2.43}  % MSE x 1e10 lasso lo_iv all
\newcommand{\rhabLassoLoIvAllPct}{0.2}  % QLIKE change % vs full
\newcommand{\rhabLassoLoIvAllDm}{0.34}  % DM vs the reference arm
\newcommand{\rhabLassoLoIvAllP}{0.731}  % p-value
\newcommand{\rhabLassoLoIvCloseQlike}{0.1284}  % QLIKE lasso lo_iv c16
\newcommand{\rhabLassoLoIvCloseMse}{3.90}  % MSE x 1e10 lasso lo_iv c16
\newcommand{\rhabLassoLoIvClosePct}{$-$2.0}  % QLIKE change % vs full
\newcommand{\rhabLassoLoIvCloseDm}{$-$1.66}  % DM vs the reference arm
\newcommand{\rhabLassoLoIvCloseP}{0.096}  % p-value
\newcommand{\rhabLassoLoIvCloseLateQlike}{0.1251}  % QLIKE lasso lo_iv c16_2020
\newcommand{\rhabLassoLoIvCloseLateMse}{8.92}  % MSE x 1e10 lasso lo_iv c16_2020
\newcommand{\rhabLassoLoIvCloseLatePct}{$-$3.8}  % QLIKE change % vs full
\newcommand{\rhabLassoLoIvCloseLateDm}{$-$1.77}  % DM vs the reference arm
\newcommand{\rhabLassoLoIvCloseLateP}{0.076}  % p-value
\newcommand{\rhabLassoLoLiqAllQlike}{0.1774}  % QLIKE lasso lo_liq all
\newcommand{\rhabLassoLoLiqAllMse}{2.39}  % MSE x 1e10 lasso lo_liq all
\newcommand{\rhabLassoLoLiqAllPct}{$-$0.1}  % QLIKE change % vs full
\newcommand{\rhabLassoLoLiqAllDm}{$-$0.21}  % DM vs the reference arm
\newcommand{\rhabLassoLoLiqAllP}{0.831}  % p-value
\newcommand{\rhabLassoLoLiqCloseQlike}{0.1285}  % QLIKE lasso lo_liq c16
\newcommand{\rhabLassoLoLiqCloseMse}{4.03}  % MSE x 1e10 lasso lo_liq c16
\newcommand{\rhabLassoLoLiqClosePct}{$-$1.9}  % QLIKE change % vs full
\newcommand{\rhabLassoLoLiqCloseDm}{$-$2.17}  % DM vs the reference arm
\newcommand{\rhabLassoLoLiqCloseP}{0.030}  % p-value
\newcommand{\rhabLassoLoLiqCloseLateQlike}{0.1272}  % QLIKE lasso lo_liq c16_2020
\newcommand{\rhabLassoLoLiqCloseLateMse}{9.23}  % MSE x 1e10 lasso lo_liq c16_2020
\newcommand{\rhabLassoLoLiqCloseLatePct}{$-$2.2}  % QLIKE change % vs full
\newcommand{\rhabLassoLoLiqCloseLateDm}{$-$1.99}  % DM vs the reference arm
\newcommand{\rhabLassoLoLiqCloseLateP}{0.047}  % p-value
\newcommand{\rhabLassoLoRmAllQlike}{0.1798}  % QLIKE lasso lo_rm all
\newcommand{\rhabLassoLoRmAllMse}{2.07}  % MSE x 1e10 lasso lo_rm all
\newcommand{\rhabLassoLoRmAllPct}{1.3}  % QLIKE change % vs full
\newcommand{\rhabLassoLoRmAllDm}{2.07}  % DM vs the reference arm
\newcommand{\rhabLassoLoRmAllP}{0.039}  % p-value
\newcommand{\rhabLassoLoRmCloseQlike}{0.1311}  % QLIKE lasso lo_rm c16
\newcommand{\rhabLassoLoRmCloseMse}{2.27}  % MSE x 1e10 lasso lo_rm c16
\newcommand{\rhabLassoLoRmClosePct}{0.0}  % QLIKE change % vs full
\newcommand{\rhabLassoLoRmCloseDm}{0.01}  % DM vs the reference arm
\newcommand{\rhabLassoLoRmCloseP}{0.990}  % p-value
\newcommand{\rhabLassoLoRmCloseLateQlike}{0.1255}  % QLIKE lasso lo_rm c16_2020
\newcommand{\rhabLassoLoRmCloseLateMse}{4.99}  % MSE x 1e10 lasso lo_rm c16_2020
\newcommand{\rhabLassoLoRmCloseLatePct}{$-$3.5}  % QLIKE change % vs full
\newcommand{\rhabLassoLoRmCloseLateDm}{$-$1.28}  % DM vs the reference arm
\newcommand{\rhabLassoLoRmCloseLateP}{0.200}  % p-value
\newcommand{\rhabLassoLoSentAllQlike}{0.1776}  % QLIKE lasso lo_sent all
\newcommand{\rhabLassoLoSentAllMse}{2.37}  % MSE x 1e10 lasso lo_sent all
\newcommand{\rhabLassoLoSentAllPct}{0.1}  % QLIKE change % vs full
\newcommand{\rhabLassoLoSentAllDm}{0.41}  % DM vs the reference arm
\newcommand{\rhabLassoLoSentAllP}{0.679}  % p-value
\newcommand{\rhabLassoLoSentCloseQlike}{0.1294}  % QLIKE lasso lo_sent c16
\newcommand{\rhabLassoLoSentCloseMse}{3.60}  % MSE x 1e10 lasso lo_sent c16
\newcommand{\rhabLassoLoSentClosePct}{$-$1.3}  % QLIKE change % vs full
\newcommand{\rhabLassoLoSentCloseDm}{$-$2.54}  % DM vs the reference arm
\newcommand{\rhabLassoLoSentCloseP}{0.011}  % p-value
\newcommand{\rhabLassoLoSentCloseLateQlike}{0.1262}  % QLIKE lasso lo_sent c16_2020
\newcommand{\rhabLassoLoSentCloseLateMse}{8.21}  % MSE x 1e10 lasso lo_sent c16_2020
\newcommand{\rhabLassoLoSentCloseLatePct}{$-$3.0}  % QLIKE change % vs full
\newcommand{\rhabLassoLoSentCloseLateDm}{$-$2.62}  % DM vs the reference arm
\newcommand{\rhabLassoLoSentCloseLateP}{0.009}  % p-value
\newcommand{\rhabLassoLoVdAllQlike}{0.1773}  % QLIKE lasso lo_vd all
\newcommand{\rhabLassoLoVdAllMse}{2.35}  % MSE x 1e10 lasso lo_vd all
\newcommand{\rhabLassoLoVdAllPct}{$-$0.1}  % QLIKE change % vs full
\newcommand{\rhabLassoLoVdAllDm}{$-$0.63}  % DM vs the reference arm
\newcommand{\rhabLassoLoVdAllP}{0.531}  % p-value
\newcommand{\rhabLassoLoVdCloseQlike}{0.1299}  % QLIKE lasso lo_vd c16
\newcommand{\rhabLassoLoVdCloseMse}{3.56}  % MSE x 1e10 lasso lo_vd c16
\newcommand{\rhabLassoLoVdClosePct}{$-$0.9}  % QLIKE change % vs full
\newcommand{\rhabLassoLoVdCloseDm}{$-$2.10}  % DM vs the reference arm
\newcommand{\rhabLassoLoVdCloseP}{0.036}  % p-value
\newcommand{\rhabLassoLoVdCloseLateQlike}{0.1275}  % QLIKE lasso lo_vd c16_2020
\newcommand{\rhabLassoLoVdCloseLateMse}{8.12}  % MSE x 1e10 lasso lo_vd c16_2020
\newcommand{\rhabLassoLoVdCloseLatePct}{$-$2.0}  % QLIKE change % vs full
\newcommand{\rhabLassoLoVdCloseLateDm}{$-$2.12}  % DM vs the reference arm
\newcommand{\rhabLassoLoVdCloseLateP}{0.034}  % p-value
\newcommand{\rhabLassoLoVwAllQlike}{0.1769}  % QLIKE lasso lo_vw all
\newcommand{\rhabLassoLoVwAllMse}{2.07}  % MSE x 1e10 lasso lo_vw all
\newcommand{\rhabLassoLoVwAllPct}{$-$0.3}  % QLIKE change % vs full
\newcommand{\rhabLassoLoVwAllDm}{$-$1.09}  % DM vs the reference arm
\newcommand{\rhabLassoLoVwAllP}{0.274}  % p-value
\newcommand{\rhabLassoLoVwCloseQlike}{0.1302}  % QLIKE lasso lo_vw c16
\newcommand{\rhabLassoLoVwCloseMse}{3.61}  % MSE x 1e10 lasso lo_vw c16
\newcommand{\rhabLassoLoVwClosePct}{$-$0.6}  % QLIKE change % vs full
\newcommand{\rhabLassoLoVwCloseDm}{$-$1.06}  % DM vs the reference arm
\newcommand{\rhabLassoLoVwCloseP}{0.289}  % p-value
\newcommand{\rhabLassoLoVwCloseLateQlike}{0.1277}  % QLIKE lasso lo_vw c16_2020
\newcommand{\rhabLassoLoVwCloseLateMse}{8.22}  % MSE x 1e10 lasso lo_vw c16_2020
\newcommand{\rhabLassoLoVwCloseLatePct}{$-$1.8}  % QLIKE change % vs full
\newcommand{\rhabLassoLoVwCloseLateDm}{$-$1.96}  % DM vs the reference arm
\newcommand{\rhabLassoLoVwCloseLateP}{0.050}  % p-value
\newcommand{\rhabLassoSharAllQlike}{0.1877}  % QLIKE lasso shar all
\newcommand{\rhabLassoSharAllMse}{1.39}  % MSE x 1e10 lasso shar all
\newcommand{\rhabLassoSharAllPct}{$-$0.8}  % QLIKE change % vs har
\newcommand{\rhabLassoSharAllDm}{$-$3.52}  % DM vs the reference arm
\newcommand{\rhabLassoSharAllP}{\ensuremath{<}0.001}  % p-value
\newcommand{\rhabLassoSharCloseQlike}{0.1265}  % QLIKE lasso shar c16
\newcommand{\rhabLassoSharCloseMse}{1.84}  % MSE x 1e10 lasso shar c16
\newcommand{\rhabLassoSharClosePct}{$-$0.3}  % QLIKE change % vs har
\newcommand{\rhabLassoSharCloseDm}{$-$0.44}  % DM vs the reference arm
\newcommand{\rhabLassoSharCloseP}{0.657}  % p-value
\newcommand{\rhabLassoSharCloseLateQlike}{0.1142}  % QLIKE lasso shar c16_2020
\newcommand{\rhabLassoSharCloseLateMse}{3.95}  % MSE x 1e10 lasso shar c16_2020
\newcommand{\rhabLassoSharCloseLatePct}{0.4}  % QLIKE change % vs har
\newcommand{\rhabLassoSharCloseLateDm}{0.42}  % DM vs the reference arm
\newcommand{\rhabLassoSharCloseLateP}{0.673}  % p-value
\newcommand{\rhabLassoSharqlAllQlike}{0.1840}  % QLIKE lasso sharql all
\newcommand{\rhabLassoSharqlAllMse}{1.69}  % MSE x 1e10 lasso sharql all
\newcommand{\rhabLassoSharqlAllPct}{$-$2.8}  % QLIKE change % vs har
\newcommand{\rhabLassoSharqlAllDm}{$-$6.63}  % DM vs the reference arm
\newcommand{\rhabLassoSharqlAllP}{\ensuremath{<}0.001}  % p-value
\newcommand{\rhabLassoSharqlCloseQlike}{0.1218}  % QLIKE lasso sharql c16
\newcommand{\rhabLassoSharqlCloseMse}{2.68}  % MSE x 1e10 lasso sharql c16
\newcommand{\rhabLassoSharqlClosePct}{$-$4.0}  % QLIKE change % vs har
\newcommand{\rhabLassoSharqlCloseDm}{$-$3.14}  % DM vs the reference arm
\newcommand{\rhabLassoSharqlCloseP}{0.002}  % p-value
\newcommand{\rhabLassoSharqlCloseLateQlike}{0.1125}  % QLIKE lasso sharql c16_2020
\newcommand{\rhabLassoSharqlCloseLateMse}{5.98}  % MSE x 1e10 lasso sharql c16_2020
\newcommand{\rhabLassoSharqlCloseLatePct}{$-$1.0}  % QLIKE change % vs har
\newcommand{\rhabLassoSharqlCloseLateDm}{$-$0.44}  % DM vs the reference arm
\newcommand{\rhabLassoSharqlCloseLateP}{0.662}  % p-value
\newcommand{\rhabLgbmFullAllQlike}{0.1613}  % QLIKE lgbm full all
\newcommand{\rhabLgbmFullAllMse}{1.41}  % MSE x 1e10 lgbm full all
\newcommand{\rhabLgbmFullAllPct}{$-$14.0}  % QLIKE change % vs har
\newcommand{\rhabLgbmFullAllDm}{$-$10.32}  % DM vs the reference arm
\newcommand{\rhabLgbmFullAllP}{\ensuremath{<}0.001}  % p-value
\newcommand{\rhabLgbmFullCloseQlike}{0.1201}  % QLIKE lgbm full c16
\newcommand{\rhabLgbmFullCloseMse}{2.34}  % MSE x 1e10 lgbm full c16
\newcommand{\rhabLgbmFullClosePct}{$-$2.7}  % QLIKE change % vs har
\newcommand{\rhabLgbmFullCloseDm}{$-$1.81}  % DM vs the reference arm
\newcommand{\rhabLgbmFullCloseP}{0.071}  % p-value
\newcommand{\rhabLgbmFullCloseLateQlike}{0.1089}  % QLIKE lgbm full c16_2020
\newcommand{\rhabLgbmFullCloseLateMse}{5.09}  % MSE x 1e10 lgbm full c16_2020
\newcommand{\rhabLgbmFullCloseLatePct}{$-$1.1}  % QLIKE change % vs har
\newcommand{\rhabLgbmFullCloseLateDm}{$-$0.39}  % DM vs the reference arm
\newcommand{\rhabLgbmFullCloseLateP}{0.697}  % p-value
\newcommand{\rhabLgbmHarAllQlike}{0.1875}  % QLIKE lgbm har all
\newcommand{\rhabLgbmHarAllMse}{1.18}  % MSE x 1e10 lgbm har all
\newcommand{\rhabLgbmHarCloseQlike}{0.1235}  % QLIKE lgbm har c16
\newcommand{\rhabLgbmHarCloseMse}{1.77}  % MSE x 1e10 lgbm har c16
\newcommand{\rhabLgbmHarCloseLateQlike}{0.1101}  % QLIKE lgbm har c16_2020
\newcommand{\rhabLgbmHarCloseLateMse}{3.68}  % MSE x 1e10 lgbm har c16_2020
\newcommand{\rhabOlsAddRmAllQlike}{0.1858}  % QLIKE ols add_rm all
\newcommand{\rhabOlsAddRmAllMse}{1.79}  % MSE x 1e10 ols add_rm all
\newcommand{\rhabOlsAddRmAllPct}{$-$1.9}  % QLIKE change % vs har
\newcommand{\rhabOlsAddRmAllDm}{$-$2.00}  % DM vs the reference arm
\newcommand{\rhabOlsAddRmAllP}{0.045}  % p-value
\newcommand{\rhabOlsAddRmCloseQlike}{0.1252}  % QLIKE ols add_rm c16
\newcommand{\rhabOlsAddRmCloseMse}{3.02}  % MSE x 1e10 ols add_rm c16
\newcommand{\rhabOlsAddRmClosePct}{$-$1.3}  % QLIKE change % vs har
\newcommand{\rhabOlsAddRmCloseDm}{$-$0.56}  % DM vs the reference arm
\newcommand{\rhabOlsAddRmCloseP}{0.576}  % p-value
\newcommand{\rhabOlsAddRmCloseLateQlike}{0.1221}  % QLIKE ols add_rm c16_2020
\newcommand{\rhabOlsAddRmCloseLateMse}{6.80}  % MSE x 1e10 ols add_rm c16_2020
\newcommand{\rhabOlsAddRmCloseLatePct}{7.5}  % QLIKE change % vs har
\newcommand{\rhabOlsAddRmCloseLateDm}{1.43}  % DM vs the reference arm
\newcommand{\rhabOlsAddRmCloseLateP}{0.152}  % p-value
\newcommand{\rhabOlsHarAllQlike}{0.1893}  % QLIKE ols har all
\newcommand{\rhabOlsHarAllMse}{1.11}  % MSE x 1e10 ols har all
\newcommand{\rhabOlsHarCloseQlike}{0.1269}  % QLIKE ols har c16
\newcommand{\rhabOlsHarCloseMse}{1.97}  % MSE x 1e10 ols har c16
\newcommand{\rhabOlsHarCloseLateQlike}{0.1137}  % QLIKE ols har c16_2020
\newcommand{\rhabOlsHarCloseLateMse}{4.24}  % MSE x 1e10 ols har c16_2020
\newcommand{\rhabOlsHarqAllQlike}{0.1893}  % QLIKE ols harq all
\newcommand{\rhabOlsHarqAllMse}{1.17}  % MSE x 1e10 ols harq all
\newcommand{\rhabOlsHarqAllPct}{0.0}  % QLIKE change % vs har
\newcommand{\rhabOlsHarqAllDm}{0.61}  % DM vs the reference arm
\newcommand{\rhabOlsHarqAllP}{0.543}  % p-value
\newcommand{\rhabOlsHarqCloseQlike}{0.1267}  % QLIKE ols harq c16
\newcommand{\rhabOlsHarqCloseMse}{1.78}  % MSE x 1e10 ols harq c16
\newcommand{\rhabOlsHarqClosePct}{$-$0.2}  % QLIKE change % vs har
\newcommand{\rhabOlsHarqCloseDm}{$-$1.38}  % DM vs the reference arm
\newcommand{\rhabOlsHarqCloseP}{0.166}  % p-value
\newcommand{\rhabOlsHarqCloseLateQlike}{0.1133}  % QLIKE ols harq c16_2020
\newcommand{\rhabOlsHarqCloseLateMse}{3.80}  % MSE x 1e10 ols harq c16_2020
\newcommand{\rhabOlsHarqCloseLatePct}{$-$0.3}  % QLIKE change % vs har
\newcommand{\rhabOlsHarqCloseLateDm}{$-$1.19}  % DM vs the reference arm
\newcommand{\rhabOlsHarqCloseLateP}{0.235}  % p-value
\newcommand{\rhabOlsLharAllQlike}{0.1842}  % QLIKE ols lhar all
\newcommand{\rhabOlsLharAllMse}{1.71}  % MSE x 1e10 ols lhar all
\newcommand{\rhabOlsLharAllPct}{$-$2.7}  % QLIKE change % vs har
\newcommand{\rhabOlsLharAllDm}{$-$6.19}  % DM vs the reference arm
\newcommand{\rhabOlsLharAllP}{\ensuremath{<}0.001}  % p-value
\newcommand{\rhabOlsLharCloseQlike}{0.1220}  % QLIKE ols lhar c16
\newcommand{\rhabOlsLharCloseMse}{2.81}  % MSE x 1e10 ols lhar c16
\newcommand{\rhabOlsLharClosePct}{$-$3.9}  % QLIKE change % vs har
\newcommand{\rhabOlsLharCloseDm}{$-$2.66}  % DM vs the reference arm
\newcommand{\rhabOlsLharCloseP}{0.008}  % p-value
\newcommand{\rhabOlsLharCloseLateQlike}{0.1132}  % QLIKE ols lhar c16_2020
\newcommand{\rhabOlsLharCloseLateMse}{6.30}  % MSE x 1e10 ols lhar c16_2020
\newcommand{\rhabOlsLharCloseLatePct}{$-$0.4}  % QLIKE change % vs har
\newcommand{\rhabOlsLharCloseLateDm}{$-$0.15}  % DM vs the reference arm
\newcommand{\rhabOlsLharCloseLateP}{0.883}  % p-value
\newcommand{\rhabOlsSharAllQlike}{0.1877}  % QLIKE ols shar all
\newcommand{\rhabOlsSharAllMse}{1.40}  % MSE x 1e10 ols shar all
\newcommand{\rhabOlsSharAllPct}{$-$0.8}  % QLIKE change % vs har
\newcommand{\rhabOlsSharAllDm}{$-$3.48}  % DM vs the reference arm
\newcommand{\rhabOlsSharAllP}{\ensuremath{<}0.001}  % p-value
\newcommand{\rhabOlsSharCloseQlike}{0.1265}  % QLIKE ols shar c16
\newcommand{\rhabOlsSharCloseMse}{1.84}  % MSE x 1e10 ols shar c16
\newcommand{\rhabOlsSharClosePct}{$-$0.3}  % QLIKE change % vs har
\newcommand{\rhabOlsSharCloseDm}{$-$0.42}  % DM vs the reference arm
\newcommand{\rhabOlsSharCloseP}{0.674}  % p-value
\newcommand{\rhabOlsSharCloseLateQlike}{0.1142}  % QLIKE ols shar c16_2020
\newcommand{\rhabOlsSharCloseLateMse}{3.95}  % MSE x 1e10 ols shar c16_2020
\newcommand{\rhabOlsSharCloseLatePct}{0.5}  % QLIKE change % vs har
\newcommand{\rhabOlsSharCloseLateDm}{0.44}  % DM vs the reference arm
\newcommand{\rhabOlsSharCloseLateP}{0.659}  % p-value
\newcommand{\rhabOlsSharqlAllQlike}{0.1841}  % QLIKE ols sharql all
\newcommand{\rhabOlsSharqlAllMse}{1.69}  % MSE x 1e10 ols sharql all
\newcommand{\rhabOlsSharqlAllPct}{$-$2.7}  % QLIKE change % vs har
\newcommand{\rhabOlsSharqlAllDm}{$-$6.53}  % DM vs the reference arm
\newcommand{\rhabOlsSharqlAllP}{\ensuremath{<}0.001}  % p-value
\newcommand{\rhabOlsSharqlCloseQlike}{0.1218}  % QLIKE ols sharql c16
\newcommand{\rhabOlsSharqlCloseMse}{2.68}  % MSE x 1e10 ols sharql c16
\newcommand{\rhabOlsSharqlClosePct}{$-$4.0}  % QLIKE change % vs har
\newcommand{\rhabOlsSharqlCloseDm}{$-$3.05}  % DM vs the reference arm
\newcommand{\rhabOlsSharqlCloseP}{0.002}  % p-value
\newcommand{\rhabOlsSharqlCloseLateQlike}{0.1126}  % QLIKE ols sharql c16_2020
\newcommand{\rhabOlsSharqlCloseLateMse}{5.99}  % MSE x 1e10 ols sharql c16_2020
\newcommand{\rhabOlsSharqlCloseLatePct}{$-$0.9}  % QLIKE change % vs har
\newcommand{\rhabOlsSharqlCloseLateDm}{$-$0.38}  % DM vs the reference arm
\newcommand{\rhabOlsSharqlCloseLateP}{0.703}  % p-value
\newcommand{\rhabRfFullAllQlike}{0.1685}  % QLIKE rf full all
\newcommand{\rhabRfFullAllMse}{1.32}  % MSE x 1e10 rf full all
\newcommand{\rhabRfFullAllPct}{$-$11.1}  % QLIKE change % vs har
\newcommand{\rhabRfFullAllDm}{$-$9.85}  % DM vs the reference arm
\newcommand{\rhabRfFullAllP}{\ensuremath{<}0.001}  % p-value
\newcommand{\rhabRfFullCloseQlike}{0.1248}  % QLIKE rf full c16
\newcommand{\rhabRfFullCloseMse}{1.99}  % MSE x 1e10 rf full c16
\newcommand{\rhabRfFullClosePct}{$-$1.4}  % QLIKE change % vs har
\newcommand{\rhabRfFullCloseDm}{$-$0.92}  % DM vs the reference arm
\newcommand{\rhabRfFullCloseP}{0.357}  % p-value
\newcommand{\rhabRfFullCloseLateQlike}{0.1123}  % QLIKE rf full c16_2020
\newcommand{\rhabRfFullCloseLateMse}{4.26}  % MSE x 1e10 rf full c16_2020
\newcommand{\rhabRfFullCloseLatePct}{$-$0.7}  % QLIKE change % vs har
\newcommand{\rhabRfFullCloseLateDm}{$-$0.26}  % DM vs the reference arm
\newcommand{\rhabRfFullCloseLateP}{0.794}  % p-value
\newcommand{\rhabRfHarAllQlike}{0.1895}  % QLIKE rf har all
\newcommand{\rhabRfHarAllMse}{1.19}  % MSE x 1e10 rf har all
\newcommand{\rhabRfHarCloseQlike}{0.1266}  % QLIKE rf har c16
\newcommand{\rhabRfHarCloseMse}{1.63}  % MSE x 1e10 rf har c16
\newcommand{\rhabRfHarCloseLateQlike}{0.1131}  % QLIKE rf har c16_2020
\newcommand{\rhabRfHarCloseLateMse}{3.39}  % MSE x 1e10 rf har c16_2020
\newcommand{\rhabRidgeAddEwAllQlike}{0.1876}  % QLIKE ridge add_ew all
\newcommand{\rhabRidgeAddEwAllMse}{1.42}  % MSE x 1e10 ridge add_ew all
\newcommand{\rhabRidgeAddEwAllPct}{$-$0.9}  % QLIKE change % vs har
\newcommand{\rhabRidgeAddEwAllDm}{$-$2.26}  % DM vs the reference arm
\newcommand{\rhabRidgeAddEwAllP}{0.024}  % p-value
\newcommand{\rhabRidgeAddEwCloseQlike}{0.1256}  % QLIKE ridge add_ew c16
\newcommand{\rhabRidgeAddEwCloseMse}{2.08}  % MSE x 1e10 ridge add_ew c16
\newcommand{\rhabRidgeAddEwClosePct}{$-$1.0}  % QLIKE change % vs har
\newcommand{\rhabRidgeAddEwCloseDm}{$-$0.89}  % DM vs the reference arm
\newcommand{\rhabRidgeAddEwCloseP}{0.374}  % p-value
\newcommand{\rhabRidgeAddEwCloseLateQlike}{0.1125}  % QLIKE ridge add_ew c16_2020
\newcommand{\rhabRidgeAddEwCloseLateMse}{4.51}  % MSE x 1e10 ridge add_ew c16_2020
\newcommand{\rhabRidgeAddEwCloseLatePct}{$-$1.0}  % QLIKE change % vs har
\newcommand{\rhabRidgeAddEwCloseLateDm}{$-$0.62}  % DM vs the reference arm
\newcommand{\rhabRidgeAddEwCloseLateP}{0.537}  % p-value
\newcommand{\rhabRidgeAddFomcAllQlike}{0.1847}  % QLIKE ridge add_fomc all
\newcommand{\rhabRidgeAddFomcAllMse}{1.13}  % MSE x 1e10 ridge add_fomc all
\newcommand{\rhabRidgeAddFomcAllPct}{$-$2.4}  % QLIKE change % vs har
\newcommand{\rhabRidgeAddFomcAllDm}{$-$5.79}  % DM vs the reference arm
\newcommand{\rhabRidgeAddFomcAllP}{\ensuremath{<}0.001}  % p-value
\newcommand{\rhabRidgeAddFomcCloseQlike}{0.1276}  % QLIKE ridge add_fomc c16
\newcommand{\rhabRidgeAddFomcCloseMse}{1.84}  % MSE x 1e10 ridge add_fomc c16
\newcommand{\rhabRidgeAddFomcClosePct}{0.6}  % QLIKE change % vs har
\newcommand{\rhabRidgeAddFomcCloseDm}{2.17}  % DM vs the reference arm
\newcommand{\rhabRidgeAddFomcCloseP}{0.030}  % p-value
\newcommand{\rhabRidgeAddFomcCloseLateQlike}{0.1135}  % QLIKE ridge add_fomc c16_2020
\newcommand{\rhabRidgeAddFomcCloseLateMse}{3.94}  % MSE x 1e10 ridge add_fomc c16_2020
\newcommand{\rhabRidgeAddFomcCloseLatePct}{$-$0.2}  % QLIKE change % vs har
\newcommand{\rhabRidgeAddFomcCloseLateDm}{$-$0.35}  % DM vs the reference arm
\newcommand{\rhabRidgeAddFomcCloseLateP}{0.729}  % p-value
\newcommand{\rhabRidgeAddIvAllQlike}{0.1897}  % QLIKE ridge add_iv all
\newcommand{\rhabRidgeAddIvAllMse}{1.28}  % MSE x 1e10 ridge add_iv all
\newcommand{\rhabRidgeAddIvAllPct}{0.2}  % QLIKE change % vs har
\newcommand{\rhabRidgeAddIvAllDm}{0.67}  % DM vs the reference arm
\newcommand{\rhabRidgeAddIvAllP}{0.506}  % p-value
\newcommand{\rhabRidgeAddIvCloseQlike}{0.1266}  % QLIKE ridge add_iv c16
\newcommand{\rhabRidgeAddIvCloseMse}{1.79}  % MSE x 1e10 ridge add_iv c16
\newcommand{\rhabRidgeAddIvClosePct}{$-$0.2}  % QLIKE change % vs har
\newcommand{\rhabRidgeAddIvCloseDm}{$-$0.20}  % DM vs the reference arm
\newcommand{\rhabRidgeAddIvCloseP}{0.845}  % p-value
\newcommand{\rhabRidgeAddIvCloseLateQlike}{0.1153}  % QLIKE ridge add_iv c16_2020
\newcommand{\rhabRidgeAddIvCloseLateMse}{3.82}  % MSE x 1e10 ridge add_iv c16_2020
\newcommand{\rhabRidgeAddIvCloseLatePct}{1.4}  % QLIKE change % vs har
\newcommand{\rhabRidgeAddIvCloseLateDm}{0.70}  % DM vs the reference arm
\newcommand{\rhabRidgeAddIvCloseLateP}{0.482}  % p-value
\newcommand{\rhabRidgeAddLiqAllQlike}{0.1885}  % QLIKE ridge add_liq all
\newcommand{\rhabRidgeAddLiqAllMse}{1.18}  % MSE x 1e10 ridge add_liq all
\newcommand{\rhabRidgeAddLiqAllPct}{$-$0.4}  % QLIKE change % vs har
\newcommand{\rhabRidgeAddLiqAllDm}{$-$0.98}  % DM vs the reference arm
\newcommand{\rhabRidgeAddLiqAllP}{0.326}  % p-value
\newcommand{\rhabRidgeAddLiqCloseQlike}{0.1304}  % QLIKE ridge add_liq c16
\newcommand{\rhabRidgeAddLiqCloseMse}{2.01}  % MSE x 1e10 ridge add_liq c16
\newcommand{\rhabRidgeAddLiqClosePct}{2.8}  % QLIKE change % vs har
\newcommand{\rhabRidgeAddLiqCloseDm}{1.91}  % DM vs the reference arm
\newcommand{\rhabRidgeAddLiqCloseP}{0.056}  % p-value
\newcommand{\rhabRidgeAddLiqCloseLateQlike}{0.1204}  % QLIKE ridge add_liq c16_2020
\newcommand{\rhabRidgeAddLiqCloseLateMse}{4.35}  % MSE x 1e10 ridge add_liq c16_2020
\newcommand{\rhabRidgeAddLiqCloseLatePct}{6.0}  % QLIKE change % vs har
\newcommand{\rhabRidgeAddLiqCloseLateDm}{2.09}  % DM vs the reference arm
\newcommand{\rhabRidgeAddLiqCloseLateP}{0.036}  % p-value
\newcommand{\rhabRidgeAddRmAllQlike}{0.1843}  % QLIKE ridge add_rm all
\newcommand{\rhabRidgeAddRmAllMse}{1.69}  % MSE x 1e10 ridge add_rm all
\newcommand{\rhabRidgeAddRmAllPct}{$-$2.6}  % QLIKE change % vs har
\newcommand{\rhabRidgeAddRmAllDm}{$-$4.14}  % DM vs the reference arm
\newcommand{\rhabRidgeAddRmAllP}{\ensuremath{<}0.001}  % p-value
\newcommand{\rhabRidgeAddRmCloseQlike}{0.1230}  % QLIKE ridge add_rm c16
\newcommand{\rhabRidgeAddRmCloseMse}{2.81}  % MSE x 1e10 ridge add_rm c16
\newcommand{\rhabRidgeAddRmClosePct}{$-$3.1}  % QLIKE change % vs har
\newcommand{\rhabRidgeAddRmCloseDm}{$-$1.80}  % DM vs the reference arm
\newcommand{\rhabRidgeAddRmCloseP}{0.071}  % p-value
\newcommand{\rhabRidgeAddRmCloseLateQlike}{0.1156}  % QLIKE ridge add_rm c16_2020
\newcommand{\rhabRidgeAddRmCloseLateMse}{6.29}  % MSE x 1e10 ridge add_rm c16_2020
\newcommand{\rhabRidgeAddRmCloseLatePct}{1.7}  % QLIKE change % vs har
\newcommand{\rhabRidgeAddRmCloseLateDm}{0.56}  % DM vs the reference arm
\newcommand{\rhabRidgeAddRmCloseLateP}{0.577}  % p-value
\newcommand{\rhabRidgeAddSentAllQlike}{0.1897}  % QLIKE ridge add_sent all
\newcommand{\rhabRidgeAddSentAllMse}{1.11}  % MSE x 1e10 ridge add_sent all
\newcommand{\rhabRidgeAddSentAllPct}{0.2}  % QLIKE change % vs har
\newcommand{\rhabRidgeAddSentAllDm}{0.82}  % DM vs the reference arm
\newcommand{\rhabRidgeAddSentAllP}{0.411}  % p-value
\newcommand{\rhabRidgeAddSentCloseQlike}{0.1310}  % QLIKE ridge add_sent c16
\newcommand{\rhabRidgeAddSentCloseMse}{1.98}  % MSE x 1e10 ridge add_sent c16
\newcommand{\rhabRidgeAddSentClosePct}{3.2}  % QLIKE change % vs har
\newcommand{\rhabRidgeAddSentCloseDm}{4.41}  % DM vs the reference arm
\newcommand{\rhabRidgeAddSentCloseP}{\ensuremath{<}0.001}  % p-value
\newcommand{\rhabRidgeAddSentCloseLateQlike}{0.1214}  % QLIKE ridge add_sent c16_2020
\newcommand{\rhabRidgeAddSentCloseLateMse}{4.27}  % MSE x 1e10 ridge add_sent c16_2020
\newcommand{\rhabRidgeAddSentCloseLatePct}{6.8}  % QLIKE change % vs har
\newcommand{\rhabRidgeAddSentCloseLateDm}{3.65}  % DM vs the reference arm
\newcommand{\rhabRidgeAddSentCloseLateP}{\ensuremath{<}0.001}  % p-value
\newcommand{\rhabRidgeAddVdAllQlike}{0.1903}  % QLIKE ridge add_vd all
\newcommand{\rhabRidgeAddVdAllMse}{1.11}  % MSE x 1e10 ridge add_vd all
\newcommand{\rhabRidgeAddVdAllPct}{0.5}  % QLIKE change % vs har
\newcommand{\rhabRidgeAddVdAllDm}{3.10}  % DM vs the reference arm
\newcommand{\rhabRidgeAddVdAllP}{0.002}  % p-value
\newcommand{\rhabRidgeAddVdCloseQlike}{0.1276}  % QLIKE ridge add_vd c16
\newcommand{\rhabRidgeAddVdCloseMse}{1.95}  % MSE x 1e10 ridge add_vd c16
\newcommand{\rhabRidgeAddVdClosePct}{0.6}  % QLIKE change % vs har
\newcommand{\rhabRidgeAddVdCloseDm}{1.55}  % DM vs the reference arm
\newcommand{\rhabRidgeAddVdCloseP}{0.120}  % p-value
\newcommand{\rhabRidgeAddVdCloseLateQlike}{0.1169}  % QLIKE ridge add_vd c16_2020
\newcommand{\rhabRidgeAddVdCloseLateMse}{4.21}  % MSE x 1e10 ridge add_vd c16_2020
\newcommand{\rhabRidgeAddVdCloseLatePct}{2.8}  % QLIKE change % vs har
\newcommand{\rhabRidgeAddVdCloseLateDm}{3.10}  % DM vs the reference arm
\newcommand{\rhabRidgeAddVdCloseLateP}{0.002}  % p-value
\newcommand{\rhabRidgeAddVwAllQlike}{0.1876}  % QLIKE ridge add_vw all
\newcommand{\rhabRidgeAddVwAllMse}{1.87}  % MSE x 1e10 ridge add_vw all
\newcommand{\rhabRidgeAddVwAllPct}{$-$0.9}  % QLIKE change % vs har
\newcommand{\rhabRidgeAddVwAllDm}{$-$1.58}  % DM vs the reference arm
\newcommand{\rhabRidgeAddVwAllP}{0.114}  % p-value
\newcommand{\rhabRidgeAddVwCloseQlike}{0.1264}  % QLIKE ridge add_vw c16
\newcommand{\rhabRidgeAddVwCloseMse}{2.27}  % MSE x 1e10 ridge add_vw c16
\newcommand{\rhabRidgeAddVwClosePct}{$-$0.4}  % QLIKE change % vs har
\newcommand{\rhabRidgeAddVwCloseDm}{$-$0.28}  % DM vs the reference arm
\newcommand{\rhabRidgeAddVwCloseP}{0.781}  % p-value
\newcommand{\rhabRidgeAddVwCloseLateQlike}{0.1139}  % QLIKE ridge add_vw c16_2020
\newcommand{\rhabRidgeAddVwCloseLateMse}{5.00}  % MSE x 1e10 ridge add_vw c16_2020
\newcommand{\rhabRidgeAddVwCloseLatePct}{0.2}  % QLIKE change % vs har
\newcommand{\rhabRidgeAddVwCloseLateDm}{0.12}  % DM vs the reference arm
\newcommand{\rhabRidgeAddVwCloseLateP}{0.907}  % p-value
\newcommand{\rhabRidgeFullAllQlike}{0.1791}  % QLIKE ridge full all
\newcommand{\rhabRidgeFullAllMse}{2.24}  % MSE x 1e10 ridge full all
\newcommand{\rhabRidgeFullAllPct}{$-$5.4}  % QLIKE change % vs har
\newcommand{\rhabRidgeFullAllDm}{$-$5.63}  % DM vs the reference arm
\newcommand{\rhabRidgeFullAllP}{\ensuremath{<}0.001}  % p-value
\newcommand{\rhabRidgeFullCloseQlike}{0.1279}  % QLIKE ridge full c16
\newcommand{\rhabRidgeFullCloseMse}{3.52}  % MSE x 1e10 ridge full c16
\newcommand{\rhabRidgeFullClosePct}{0.8}  % QLIKE change % vs har
\newcommand{\rhabRidgeFullCloseDm}{0.31}  % DM vs the reference arm
\newcommand{\rhabRidgeFullCloseP}{0.759}  % p-value
\newcommand{\rhabRidgeFullCloseLateQlike}{0.1260}  % QLIKE ridge full c16_2020
\newcommand{\rhabRidgeFullCloseLateMse}{7.99}  % MSE x 1e10 ridge full c16_2020
\newcommand{\rhabRidgeFullCloseLatePct}{10.9}  % QLIKE change % vs har
\newcommand{\rhabRidgeFullCloseLateDm}{2.31}  % DM vs the reference arm
\newcommand{\rhabRidgeFullCloseLateP}{0.021}  % p-value
\newcommand{\rhabRidgeHarAllQlike}{0.1893}  % QLIKE ridge har all
\newcommand{\rhabRidgeHarAllMse}{1.11}  % MSE x 1e10 ridge har all
\newcommand{\rhabRidgeHarCloseQlike}{0.1269}  % QLIKE ridge har c16
\newcommand{\rhabRidgeHarCloseMse}{1.97}  % MSE x 1e10 ridge har c16
\newcommand{\rhabRidgeHarCloseLateQlike}{0.1137}  % QLIKE ridge har c16_2020
\newcommand{\rhabRidgeHarCloseLateMse}{4.24}  % MSE x 1e10 ridge har c16_2020
\newcommand{\rhabRidgeHarqAllQlike}{0.1893}  % QLIKE ridge harq all
\newcommand{\rhabRidgeHarqAllMse}{1.16}  % MSE x 1e10 ridge harq all
\newcommand{\rhabRidgeHarqAllPct}{0.0}  % QLIKE change % vs har
\newcommand{\rhabRidgeHarqAllDm}{0.76}  % DM vs the reference arm
\newcommand{\rhabRidgeHarqAllP}{0.446}  % p-value
\newcommand{\rhabRidgeHarqCloseQlike}{0.1267}  % QLIKE ridge harq c16
\newcommand{\rhabRidgeHarqCloseMse}{1.79}  % MSE x 1e10 ridge harq c16
\newcommand{\rhabRidgeHarqClosePct}{$-$0.2}  % QLIKE change % vs har
\newcommand{\rhabRidgeHarqCloseDm}{$-$1.42}  % DM vs the reference arm
\newcommand{\rhabRidgeHarqCloseP}{0.157}  % p-value
\newcommand{\rhabRidgeHarqCloseLateQlike}{0.1133}  % QLIKE ridge harq c16_2020
\newcommand{\rhabRidgeHarqCloseLateMse}{3.81}  % MSE x 1e10 ridge harq c16_2020
\newcommand{\rhabRidgeHarqCloseLatePct}{$-$0.3}  % QLIKE change % vs har
\newcommand{\rhabRidgeHarqCloseLateDm}{$-$1.20}  % DM vs the reference arm
\newcommand{\rhabRidgeHarqCloseLateP}{0.228}  % p-value
\newcommand{\rhabRidgeLharAllQlike}{0.1841}  % QLIKE ridge lhar all
\newcommand{\rhabRidgeLharAllMse}{1.72}  % MSE x 1e10 ridge lhar all
\newcommand{\rhabRidgeLharAllPct}{$-$2.7}  % QLIKE change % vs har
\newcommand{\rhabRidgeLharAllDm}{$-$6.56}  % DM vs the reference arm
\newcommand{\rhabRidgeLharAllP}{\ensuremath{<}0.001}  % p-value
\newcommand{\rhabRidgeLharCloseQlike}{0.1216}  % QLIKE ridge lhar c16
\newcommand{\rhabRidgeLharCloseMse}{2.78}  % MSE x 1e10 ridge lhar c16
\newcommand{\rhabRidgeLharClosePct}{$-$4.2}  % QLIKE change % vs har
\newcommand{\rhabRidgeLharCloseDm}{$-$3.28}  % DM vs the reference arm
\newcommand{\rhabRidgeLharCloseP}{0.001}  % p-value
\newcommand{\rhabRidgeLharCloseLateQlike}{0.1121}  % QLIKE ridge lhar c16_2020
\newcommand{\rhabRidgeLharCloseLateMse}{6.23}  % MSE x 1e10 ridge lhar c16_2020
\newcommand{\rhabRidgeLharCloseLatePct}{$-$1.4}  % QLIKE change % vs har
\newcommand{\rhabRidgeLharCloseLateDm}{$-$0.59}  % DM vs the reference arm
\newcommand{\rhabRidgeLharCloseLateP}{0.554}  % p-value
\newcommand{\rhabRidgeLoEwAllQlike}{0.1789}  % QLIKE ridge lo_ew all
\newcommand{\rhabRidgeLoEwAllMse}{2.38}  % MSE x 1e10 ridge lo_ew all
\newcommand{\rhabRidgeLoEwAllPct}{$-$0.1}  % QLIKE change % vs full
\newcommand{\rhabRidgeLoEwAllDm}{$-$1.20}  % DM vs the reference arm
\newcommand{\rhabRidgeLoEwAllP}{0.230}  % p-value
\newcommand{\rhabRidgeLoEwCloseQlike}{0.1285}  % QLIKE ridge lo_ew c16
\newcommand{\rhabRidgeLoEwCloseMse}{3.74}  % MSE x 1e10 ridge lo_ew c16
\newcommand{\rhabRidgeLoEwClosePct}{0.5}  % QLIKE change % vs full
\newcommand{\rhabRidgeLoEwCloseDm}{1.23}  % DM vs the reference arm
\newcommand{\rhabRidgeLoEwCloseP}{0.218}  % p-value
\newcommand{\rhabRidgeLoEwCloseLateQlike}{0.1262}  % QLIKE ridge lo_ew c16_2020
\newcommand{\rhabRidgeLoEwCloseLateMse}{8.53}  % MSE x 1e10 ridge lo_ew c16_2020
\newcommand{\rhabRidgeLoEwCloseLatePct}{0.2}  % QLIKE change % vs full
\newcommand{\rhabRidgeLoEwCloseLateDm}{0.22}  % DM vs the reference arm
\newcommand{\rhabRidgeLoEwCloseLateP}{0.828}  % p-value
\newcommand{\rhabRidgeLoFomcAllQlike}{0.1833}  % QLIKE ridge lo_fomc all
\newcommand{\rhabRidgeLoFomcAllMse}{2.25}  % MSE x 1e10 ridge lo_fomc all
\newcommand{\rhabRidgeLoFomcAllPct}{2.3}  % QLIKE change % vs full
\newcommand{\rhabRidgeLoFomcAllDm}{5.55}  % DM vs the reference arm
\newcommand{\rhabRidgeLoFomcAllP}{\ensuremath{<}0.001}  % p-value
\newcommand{\rhabRidgeLoFomcCloseQlike}{0.1275}  % QLIKE ridge lo_fomc c16
\newcommand{\rhabRidgeLoFomcCloseMse}{3.58}  % MSE x 1e10 ridge lo_fomc c16
\newcommand{\rhabRidgeLoFomcClosePct}{$-$0.3}  % QLIKE change % vs full
\newcommand{\rhabRidgeLoFomcCloseDm}{$-$1.23}  % DM vs the reference arm
\newcommand{\rhabRidgeLoFomcCloseP}{0.219}  % p-value
\newcommand{\rhabRidgeLoFomcCloseLateQlike}{0.1257}  % QLIKE ridge lo_fomc c16_2020
\newcommand{\rhabRidgeLoFomcCloseLateMse}{8.15}  % MSE x 1e10 ridge lo_fomc c16_2020
\newcommand{\rhabRidgeLoFomcCloseLatePct}{$-$0.2}  % QLIKE change % vs full
\newcommand{\rhabRidgeLoFomcCloseLateDm}{$-$0.49}  % DM vs the reference arm
\newcommand{\rhabRidgeLoFomcCloseLateP}{0.622}  % p-value
\newcommand{\rhabRidgeLoIvAllQlike}{0.1791}  % QLIKE ridge lo_iv all
\newcommand{\rhabRidgeLoIvAllMse}{2.30}  % MSE x 1e10 ridge lo_iv all
\newcommand{\rhabRidgeLoIvAllPct}{0.0}  % QLIKE change % vs full
\newcommand{\rhabRidgeLoIvAllDm}{0.03}  % DM vs the reference arm
\newcommand{\rhabRidgeLoIvAllP}{0.980}  % p-value
\newcommand{\rhabRidgeLoIvCloseQlike}{0.1267}  % QLIKE ridge lo_iv c16
\newcommand{\rhabRidgeLoIvCloseMse}{3.61}  % MSE x 1e10 ridge lo_iv c16
\newcommand{\rhabRidgeLoIvClosePct}{$-$1.0}  % QLIKE change % vs full
\newcommand{\rhabRidgeLoIvCloseDm}{$-$1.02}  % DM vs the reference arm
\newcommand{\rhabRidgeLoIvCloseP}{0.309}  % p-value
\newcommand{\rhabRidgeLoIvCloseLateQlike}{0.1228}  % QLIKE ridge lo_iv c16_2020
\newcommand{\rhabRidgeLoIvCloseLateMse}{8.20}  % MSE x 1e10 ridge lo_iv c16_2020
\newcommand{\rhabRidgeLoIvCloseLatePct}{$-$2.5}  % QLIKE change % vs full
\newcommand{\rhabRidgeLoIvCloseLateDm}{$-$1.41}  % DM vs the reference arm
\newcommand{\rhabRidgeLoIvCloseLateP}{0.159}  % p-value
\newcommand{\rhabRidgeLoLiqAllQlike}{0.1794}  % QLIKE ridge lo_liq all
\newcommand{\rhabRidgeLoLiqAllMse}{2.29}  % MSE x 1e10 ridge lo_liq all
\newcommand{\rhabRidgeLoLiqAllPct}{0.2}  % QLIKE change % vs full
\newcommand{\rhabRidgeLoLiqAllDm}{0.75}  % DM vs the reference arm
\newcommand{\rhabRidgeLoLiqAllP}{0.454}  % p-value
\newcommand{\rhabRidgeLoLiqCloseQlike}{0.1269}  % QLIKE ridge lo_liq c16
\newcommand{\rhabRidgeLoLiqCloseMse}{3.93}  % MSE x 1e10 ridge lo_liq c16
\newcommand{\rhabRidgeLoLiqClosePct}{$-$0.8}  % QLIKE change % vs full
\newcommand{\rhabRidgeLoLiqCloseDm}{$-$1.14}  % DM vs the reference arm
\newcommand{\rhabRidgeLoLiqCloseP}{0.256}  % p-value
\newcommand{\rhabRidgeLoLiqCloseLateQlike}{0.1251}  % QLIKE ridge lo_liq c16_2020
\newcommand{\rhabRidgeLoLiqCloseLateMse}{8.99}  % MSE x 1e10 ridge lo_liq c16_2020
\newcommand{\rhabRidgeLoLiqCloseLatePct}{$-$0.7}  % QLIKE change % vs full
\newcommand{\rhabRidgeLoLiqCloseLateDm}{$-$0.77}  % DM vs the reference arm
\newcommand{\rhabRidgeLoLiqCloseLateP}{0.442}  % p-value
\newcommand{\rhabRidgeLoRmAllQlike}{0.1822}  % QLIKE ridge lo_rm all
\newcommand{\rhabRidgeLoRmAllMse}{1.82}  % MSE x 1e10 ridge lo_rm all
\newcommand{\rhabRidgeLoRmAllPct}{1.7}  % QLIKE change % vs full
\newcommand{\rhabRidgeLoRmAllDm}{3.75}  % DM vs the reference arm
\newcommand{\rhabRidgeLoRmAllP}{\ensuremath{<}0.001}  % p-value
\newcommand{\rhabRidgeLoRmCloseQlike}{0.1292}  % QLIKE ridge lo_rm c16
\newcommand{\rhabRidgeLoRmCloseMse}{2.26}  % MSE x 1e10 ridge lo_rm c16
\newcommand{\rhabRidgeLoRmClosePct}{1.0}  % QLIKE change % vs full
\newcommand{\rhabRidgeLoRmCloseDm}{0.73}  % DM vs the reference arm
\newcommand{\rhabRidgeLoRmCloseP}{0.466}  % p-value
\newcommand{\rhabRidgeLoRmCloseLateQlike}{0.1245}  % QLIKE ridge lo_rm c16_2020
\newcommand{\rhabRidgeLoRmCloseLateMse}{4.97}  % MSE x 1e10 ridge lo_rm c16_2020
\newcommand{\rhabRidgeLoRmCloseLatePct}{$-$1.2}  % QLIKE change % vs full
\newcommand{\rhabRidgeLoRmCloseLateDm}{$-$0.55}  % DM vs the reference arm
\newcommand{\rhabRidgeLoRmCloseLateP}{0.579}  % p-value
\newcommand{\rhabRidgeLoSentAllQlike}{0.1793}  % QLIKE ridge lo_sent all
\newcommand{\rhabRidgeLoSentAllMse}{2.24}  % MSE x 1e10 ridge lo_sent all
\newcommand{\rhabRidgeLoSentAllPct}{0.1}  % QLIKE change % vs full
\newcommand{\rhabRidgeLoSentAllDm}{0.55}  % DM vs the reference arm
\newcommand{\rhabRidgeLoSentAllP}{0.584}  % p-value
\newcommand{\rhabRidgeLoSentCloseQlike}{0.1261}  % QLIKE ridge lo_sent c16
\newcommand{\rhabRidgeLoSentCloseMse}{3.51}  % MSE x 1e10 ridge lo_sent c16
\newcommand{\rhabRidgeLoSentClosePct}{$-$1.4}  % QLIKE change % vs full
\newcommand{\rhabRidgeLoSentCloseDm}{$-$3.16}  % DM vs the reference arm
\newcommand{\rhabRidgeLoSentCloseP}{0.002}  % p-value
\newcommand{\rhabRidgeLoSentCloseLateQlike}{0.1218}  % QLIKE ridge lo_sent c16_2020
\newcommand{\rhabRidgeLoSentCloseLateMse}{7.97}  % MSE x 1e10 ridge lo_sent c16_2020
\newcommand{\rhabRidgeLoSentCloseLatePct}{$-$3.3}  % QLIKE change % vs full
\newcommand{\rhabRidgeLoSentCloseLateDm}{$-$3.21}  % DM vs the reference arm
\newcommand{\rhabRidgeLoSentCloseLateP}{0.001}  % p-value
\newcommand{\rhabRidgeLoVdAllQlike}{0.1789}  % QLIKE ridge lo_vd all
\newcommand{\rhabRidgeLoVdAllMse}{2.24}  % MSE x 1e10 ridge lo_vd all
\newcommand{\rhabRidgeLoVdAllPct}{$-$0.1}  % QLIKE change % vs full
\newcommand{\rhabRidgeLoVdAllDm}{$-$0.98}  % DM vs the reference arm
\newcommand{\rhabRidgeLoVdAllP}{0.326}  % p-value
\newcommand{\rhabRidgeLoVdCloseQlike}{0.1266}  % QLIKE ridge lo_vd c16
\newcommand{\rhabRidgeLoVdCloseMse}{3.51}  % MSE x 1e10 ridge lo_vd c16
\newcommand{\rhabRidgeLoVdClosePct}{$-$1.0}  % QLIKE change % vs full
\newcommand{\rhabRidgeLoVdCloseDm}{$-$2.17}  % DM vs the reference arm
\newcommand{\rhabRidgeLoVdCloseP}{0.030}  % p-value
\newcommand{\rhabRidgeLoVdCloseLateQlike}{0.1232}  % QLIKE ridge lo_vd c16_2020
\newcommand{\rhabRidgeLoVdCloseLateMse}{7.98}  % MSE x 1e10 ridge lo_vd c16_2020
\newcommand{\rhabRidgeLoVdCloseLatePct}{$-$2.2}  % QLIKE change % vs full
\newcommand{\rhabRidgeLoVdCloseLateDm}{$-$2.26}  % DM vs the reference arm
\newcommand{\rhabRidgeLoVdCloseLateP}{0.024}  % p-value
\newcommand{\rhabRidgeLoVwAllQlike}{0.1789}  % QLIKE ridge lo_vw all
\newcommand{\rhabRidgeLoVwAllMse}{1.97}  % MSE x 1e10 ridge lo_vw all
\newcommand{\rhabRidgeLoVwAllPct}{$-$0.1}  % QLIKE change % vs full
\newcommand{\rhabRidgeLoVwAllDm}{$-$0.70}  % DM vs the reference arm
\newcommand{\rhabRidgeLoVwAllP}{0.484}  % p-value
\newcommand{\rhabRidgeLoVwCloseQlike}{0.1274}  % QLIKE ridge lo_vw c16
\newcommand{\rhabRidgeLoVwCloseMse}{3.50}  % MSE x 1e10 ridge lo_vw c16
\newcommand{\rhabRidgeLoVwClosePct}{$-$0.4}  % QLIKE change % vs full
\newcommand{\rhabRidgeLoVwCloseDm}{$-$0.90}  % DM vs the reference arm
\newcommand{\rhabRidgeLoVwCloseP}{0.366}  % p-value
\newcommand{\rhabRidgeLoVwCloseLateQlike}{0.1246}  % QLIKE ridge lo_vw c16_2020
\newcommand{\rhabRidgeLoVwCloseLateMse}{7.95}  % MSE x 1e10 ridge lo_vw c16_2020
\newcommand{\rhabRidgeLoVwCloseLatePct}{$-$1.1}  % QLIKE change % vs full
\newcommand{\rhabRidgeLoVwCloseLateDm}{$-$1.96}  % DM vs the reference arm
\newcommand{\rhabRidgeLoVwCloseLateP}{0.050}  % p-value
\newcommand{\rhabRidgeSharAllQlike}{0.1876}  % QLIKE ridge shar all
\newcommand{\rhabRidgeSharAllMse}{1.35}  % MSE x 1e10 ridge shar all
\newcommand{\rhabRidgeSharAllPct}{$-$0.9}  % QLIKE change % vs har
\newcommand{\rhabRidgeSharAllDm}{$-$3.95}  % DM vs the reference arm
\newcommand{\rhabRidgeSharAllP}{\ensuremath{<}0.001}  % p-value
\newcommand{\rhabRidgeSharCloseQlike}{0.1264}  % QLIKE ridge shar c16
\newcommand{\rhabRidgeSharCloseMse}{1.81}  % MSE x 1e10 ridge shar c16
\newcommand{\rhabRidgeSharClosePct}{$-$0.4}  % QLIKE change % vs har
\newcommand{\rhabRidgeSharCloseDm}{$-$0.69}  % DM vs the reference arm
\newcommand{\rhabRidgeSharCloseP}{0.488}  % p-value
\newcommand{\rhabRidgeSharCloseLateQlike}{0.1139}  % QLIKE ridge shar c16_2020
\newcommand{\rhabRidgeSharCloseLateMse}{3.87}  % MSE x 1e10 ridge shar c16_2020
\newcommand{\rhabRidgeSharCloseLatePct}{0.2}  % QLIKE change % vs har
\newcommand{\rhabRidgeSharCloseLateDm}{0.23}  % DM vs the reference arm
\newcommand{\rhabRidgeSharCloseLateP}{0.821}  % p-value
\newcommand{\rhabRidgeSharqlAllQlike}{0.1840}  % QLIKE ridge sharql all
\newcommand{\rhabRidgeSharqlAllMse}{1.68}  % MSE x 1e10 ridge sharql all
\newcommand{\rhabRidgeSharqlAllPct}{$-$2.8}  % QLIKE change % vs har
\newcommand{\rhabRidgeSharqlAllDm}{$-$6.82}  % DM vs the reference arm
\newcommand{\rhabRidgeSharqlAllP}{\ensuremath{<}0.001}  % p-value
\newcommand{\rhabRidgeSharqlCloseQlike}{0.1216}  % QLIKE ridge sharql c16
\newcommand{\rhabRidgeSharqlCloseMse}{2.65}  % MSE x 1e10 ridge sharql c16
\newcommand{\rhabRidgeSharqlClosePct}{$-$4.2}  % QLIKE change % vs har
\newcommand{\rhabRidgeSharqlCloseDm}{$-$3.49}  % DM vs the reference arm
\newcommand{\rhabRidgeSharqlCloseP}{\ensuremath{<}0.001}  % p-value
\newcommand{\rhabRidgeSharqlCloseLateQlike}{0.1119}  % QLIKE ridge sharql c16_2020
\newcommand{\rhabRidgeSharqlCloseLateMse}{5.92}  % MSE x 1e10 ridge sharql c16_2020
\newcommand{\rhabRidgeSharqlCloseLatePct}{$-$1.6}  % QLIKE change % vs har
\newcommand{\rhabRidgeSharqlCloseLateDm}{$-$0.73}  % DM vs the reference arm
\newcommand{\rhabRidgeSharqlCloseLateP}{0.468}  % p-value
\newcommand{\rhabXgbFullAllQlike}{0.1627}  % QLIKE xgb full all
\newcommand{\rhabXgbFullAllMse}{1.34}  % MSE x 1e10 xgb full all
\newcommand{\rhabXgbFullAllPct}{$-$13.8}  % QLIKE change % vs har
\newcommand{\rhabXgbFullAllDm}{$-$9.94}  % DM vs the reference arm
\newcommand{\rhabXgbFullAllP}{\ensuremath{<}0.001}  % p-value
\newcommand{\rhabXgbFullCloseQlike}{0.1192}  % QLIKE xgb full c16
\newcommand{\rhabXgbFullCloseMse}{2.20}  % MSE x 1e10 xgb full c16
\newcommand{\rhabXgbFullClosePct}{$-$3.4}  % QLIKE change % vs har
\newcommand{\rhabXgbFullCloseDm}{$-$2.18}  % DM vs the reference arm
\newcommand{\rhabXgbFullCloseP}{0.030}  % p-value
\newcommand{\rhabXgbFullCloseLateQlike}{0.1079}  % QLIKE xgb full c16_2020
\newcommand{\rhabXgbFullCloseLateMse}{4.79}  % MSE x 1e10 xgb full c16_2020
\newcommand{\rhabXgbFullCloseLatePct}{$-$2.1}  % QLIKE change % vs har
\newcommand{\rhabXgbFullCloseLateDm}{$-$0.76}  % DM vs the reference arm
\newcommand{\rhabXgbFullCloseLateP}{0.444}  % p-value
\newcommand{\rhabXgbHarAllQlike}{0.1888}  % QLIKE xgb har all
\newcommand{\rhabXgbHarAllMse}{1.16}  % MSE x 1e10 xgb har all
\newcommand{\rhabXgbHarCloseQlike}{0.1234}  % QLIKE xgb har c16
\newcommand{\rhabXgbHarCloseMse}{1.73}  % MSE x 1e10 xgb har c16
\newcommand{\rhabXgbHarCloseLateQlike}{0.1102}  % QLIKE xgb har c16_2020
\newcommand{\rhabXgbHarCloseLateMse}{3.62}  % MSE x 1e10 xgb har c16_2020
\newcommand{\rhabRidgeFullLiveSentQlike}{0.1833}  % QLIKE ridge full live_sent (subsample)
\newcommand{\rhabRidgeFullLiveSentPct}{$-$1.9}  % QLIKE change % vs the reference (subsample)
\newcommand{\rhabRidgeFullLiveSentDm}{$-$0.98}  % DM vs the reference arm (subsample)
\newcommand{\rhabRidgeFullLiveSentP}{0.327}  % p-value
\newcommand{\rhabRidgeLoSentLiveSentQlike}{0.1842}  % QLIKE ridge lo_sent live_sent (subsample)
\newcommand{\rhabRidgeLoSentLiveSentPct}{0.5}  % QLIKE change % vs the reference (subsample)
\newcommand{\rhabRidgeLoSentLiveSentDm}{0.71}  % DM vs the reference arm (subsample)
\newcommand{\rhabRidgeLoSentLiveSentP}{0.477}  % p-value
\newcommand{\rhabRidgeAddSentLiveSentQlike}{0.1846}  % QLIKE ridge add_sent live_sent (subsample)
\newcommand{\rhabRidgeAddSentLiveSentPct}{$-$1.2}  % QLIKE change % vs the reference (subsample)
\newcommand{\rhabRidgeAddSentLiveSentDm}{$-$1.18}  % DM vs the reference arm (subsample)
\newcommand{\rhabRidgeAddSentLiveSentP}{0.238}  % p-value
\newcommand{\rhabRidgeHarLiveSentQlike}{0.1868}  % QLIKE ridge har live_sent (subsample)
\newcommand{\rhabLassoFullLiveSentQlike}{0.1780}  % QLIKE lasso full live_sent (subsample)
\newcommand{\rhabLassoFullLiveSentPct}{$-$4.7}  % QLIKE change % vs the reference (subsample)
\newcommand{\rhabLassoFullLiveSentDm}{$-$1.83}  % DM vs the reference arm (subsample)
\newcommand{\rhabLassoFullLiveSentP}{0.068}  % p-value
\newcommand{\rhabLassoLoSentLiveSentQlike}{0.1786}  % QLIKE lasso lo_sent live_sent (subsample)
\newcommand{\rhabLassoLoSentLiveSentPct}{0.3}  % QLIKE change % vs the reference (subsample)
\newcommand{\rhabLassoLoSentLiveSentDm}{0.41}  % DM vs the reference arm (subsample)
\newcommand{\rhabLassoLoSentLiveSentP}{0.685}  % p-value
\newcommand{\rhabLassoAddSentLiveSentQlike}{0.1850}  % QLIKE lasso add_sent live_sent (subsample)
\newcommand{\rhabLassoAddSentLiveSentPct}{$-$1.0}  % QLIKE change % vs the reference (subsample)
\newcommand{\rhabLassoAddSentLiveSentDm}{$-$0.96}  % DM vs the reference arm (subsample)
\newcommand{\rhabLassoAddSentLiveSentP}{0.339}  % p-value
\newcommand{\rhabLassoHarLiveSentQlike}{0.1868}  % QLIKE lasso har live_sent (subsample)
\newcommand{\rhabEnetFullLiveSentQlike}{0.1818}  % QLIKE enet full live_sent (subsample)
\newcommand{\rhabEnetFullLiveSentPct}{$-$2.7}  % QLIKE change % vs the reference (subsample)
\newcommand{\rhabEnetFullLiveSentDm}{$-$1.55}  % DM vs the reference arm (subsample)
\newcommand{\rhabEnetFullLiveSentP}{0.121}  % p-value
\newcommand{\rhabEnetLoSentLiveSentQlike}{0.1833}  % QLIKE enet lo_sent live_sent (subsample)
\newcommand{\rhabEnetLoSentLiveSentPct}{0.8}  % QLIKE change % vs the reference (subsample)
\newcommand{\rhabEnetLoSentLiveSentDm}{1.43}  % DM vs the reference arm (subsample)
\newcommand{\rhabEnetLoSentLiveSentP}{0.154}  % p-value
\newcommand{\rhabEnetAddSentLiveSentQlike}{0.1845}  % QLIKE enet add_sent live_sent (subsample)
\newcommand{\rhabEnetAddSentLiveSentPct}{$-$1.3}  % QLIKE change % vs the reference (subsample)
\newcommand{\rhabEnetAddSentLiveSentDm}{$-$1.48}  % DM vs the reference arm (subsample)
\newcommand{\rhabEnetAddSentLiveSentP}{0.139}  % p-value
\newcommand{\rhabEnetHarLiveSentQlike}{0.1868}  % QLIKE enet har live_sent (subsample)
\newcommand{\rhabLgbmFullLiveSentQlike}{0.1564}  % QLIKE lgbm full live_sent (subsample)
\newcommand{\rhabLgbmFullLiveSentPct}{$-$13.8}  % QLIKE change % vs the reference (subsample)
\newcommand{\rhabLgbmFullLiveSentDm}{$-$3.76}  % DM vs the reference arm (subsample)
\newcommand{\rhabLgbmFullLiveSentP}{\ensuremath{<}0.001}  % p-value
\newcommand{\rhabLgbmHarLiveSentQlike}{0.1814}  % QLIKE lgbm har live_sent (subsample)
\newcommand{\rhabXgbFullLiveSentQlike}{0.1559}  % QLIKE xgb full live_sent (subsample)
\newcommand{\rhabXgbFullLiveSentPct}{$-$15.6}  % QLIKE change % vs the reference (subsample)
\newcommand{\rhabXgbFullLiveSentDm}{$-$3.69}  % DM vs the reference arm (subsample)
\newcommand{\rhabXgbFullLiveSentP}{\ensuremath{<}0.001}  % p-value
\newcommand{\rhabXgbHarLiveSentQlike}{0.1846}  % QLIKE xgb har live_sent (subsample)
\newcommand{\rhabRfFullLiveSentQlike}{0.1676}  % QLIKE rf full live_sent (subsample)
\newcommand{\rhabRfFullLiveSentPct}{$-$10.8}  % QLIKE change % vs the reference (subsample)
\newcommand{\rhabRfFullLiveSentDm}{$-$4.09}  % DM vs the reference arm (subsample)
\newcommand{\rhabRfFullLiveSentP}{\ensuremath{<}0.001}  % p-value
\newcommand{\rhabRfHarLiveSentQlike}{0.1878}  % QLIKE rf har live_sent (subsample)
\newcommand{\rhabRidgeFullLiveSentCloseQlike}{0.1222}  % QLIKE ridge full live16_sent (subsample)
\newcommand{\rhabRidgeFullLiveSentClosePct}{12.8}  % QLIKE change % vs the reference (subsample)
\newcommand{\rhabRidgeFullLiveSentCloseDm}{2.03}  % DM vs the reference arm (subsample)
\newcommand{\rhabRidgeFullLiveSentCloseP}{0.042}  % p-value
\newcommand{\rhabRidgeLoSentLiveSentCloseQlike}{0.1141}  % QLIKE ridge lo_sent live16_sent (subsample)
\newcommand{\rhabRidgeLoSentLiveSentClosePct}{$-$6.6}  % QLIKE change % vs the reference (subsample)
\newcommand{\rhabRidgeLoSentLiveSentCloseDm}{$-$3.48}  % DM vs the reference arm (subsample)
\newcommand{\rhabRidgeLoSentLiveSentCloseP}{\ensuremath{<}0.001}  % p-value
\newcommand{\rhabRidgeAddSentLiveSentCloseQlike}{0.1188}  % QLIKE ridge add_sent live16_sent (subsample)
\newcommand{\rhabRidgeAddSentLiveSentClosePct}{9.7}  % QLIKE change % vs the reference (subsample)
\newcommand{\rhabRidgeAddSentLiveSentCloseDm}{2.63}  % DM vs the reference arm (subsample)
\newcommand{\rhabRidgeAddSentLiveSentCloseP}{0.009}  % p-value
\newcommand{\rhabRidgeHarLiveSentCloseQlike}{0.1083}  % QLIKE ridge har live16_sent (subsample)
\newcommand{\rhabLassoFullLiveSentCloseQlike}{0.1234}  % QLIKE lasso full live16_sent (subsample)
\newcommand{\rhabLassoFullLiveSentClosePct}{13.9}  % QLIKE change % vs the reference (subsample)
\newcommand{\rhabLassoFullLiveSentCloseDm}{2.17}  % DM vs the reference arm (subsample)
\newcommand{\rhabLassoFullLiveSentCloseP}{0.030}  % p-value
\newcommand{\rhabLassoLoSentLiveSentCloseQlike}{0.1155}  % QLIKE lasso lo_sent live16_sent (subsample)
\newcommand{\rhabLassoLoSentLiveSentClosePct}{$-$6.3}  % QLIKE change % vs the reference (subsample)
\newcommand{\rhabLassoLoSentLiveSentCloseDm}{$-$2.92}  % DM vs the reference arm (subsample)
\newcommand{\rhabLassoLoSentLiveSentCloseP}{0.003}  % p-value
\newcommand{\rhabLassoAddSentLiveSentCloseQlike}{0.1195}  % QLIKE lasso add_sent live16_sent (subsample)
\newcommand{\rhabLassoAddSentLiveSentClosePct}{10.3}  % QLIKE change % vs the reference (subsample)
\newcommand{\rhabLassoAddSentLiveSentCloseDm}{2.76}  % DM vs the reference arm (subsample)
\newcommand{\rhabLassoAddSentLiveSentCloseP}{0.006}  % p-value
\newcommand{\rhabLassoHarLiveSentCloseQlike}{0.1083}  % QLIKE lasso har live16_sent (subsample)
\newcommand{\rhabEnetFullLiveSentCloseQlike}{0.1171}  % QLIKE enet full live16_sent (subsample)
\newcommand{\rhabEnetFullLiveSentClosePct}{8.1}  % QLIKE change % vs the reference (subsample)
\newcommand{\rhabEnetFullLiveSentCloseDm}{1.35}  % DM vs the reference arm (subsample)
\newcommand{\rhabEnetFullLiveSentCloseP}{0.176}  % p-value
\newcommand{\rhabEnetLoSentLiveSentCloseQlike}{0.1134}  % QLIKE enet lo_sent live16_sent (subsample)
\newcommand{\rhabEnetLoSentLiveSentClosePct}{$-$3.2}  % QLIKE change % vs the reference (subsample)
\newcommand{\rhabEnetLoSentLiveSentCloseDm}{$-$2.23}  % DM vs the reference arm (subsample)
\newcommand{\rhabEnetLoSentLiveSentCloseP}{0.026}  % p-value
\newcommand{\rhabEnetAddSentLiveSentCloseQlike}{0.1156}  % QLIKE enet add_sent live16_sent (subsample)
\newcommand{\rhabEnetAddSentLiveSentClosePct}{6.7}  % QLIKE change % vs the reference (subsample)
\newcommand{\rhabEnetAddSentLiveSentCloseDm}{1.91}  % DM vs the reference arm (subsample)
\newcommand{\rhabEnetAddSentLiveSentCloseP}{0.056}  % p-value
\newcommand{\rhabEnetHarLiveSentCloseQlike}{0.1083}  % QLIKE enet har live16_sent (subsample)
\newcommand{\rhabLgbmFullLiveSentCloseQlike}{0.0995}  % QLIKE lgbm full live16_sent (subsample)
\newcommand{\rhabLgbmFullLiveSentClosePct}{$-$2.5}  % QLIKE change % vs the reference (subsample)
\newcommand{\rhabLgbmFullLiveSentCloseDm}{$-$0.57}  % DM vs the reference arm (subsample)
\newcommand{\rhabLgbmFullLiveSentCloseP}{0.565}  % p-value
\newcommand{\rhabLgbmHarLiveSentCloseQlike}{0.1021}  % QLIKE lgbm har live16_sent (subsample)
\newcommand{\rhabXgbFullLiveSentCloseQlike}{0.0997}  % QLIKE xgb full live16_sent (subsample)
\newcommand{\rhabXgbFullLiveSentClosePct}{$-$2.8}  % QLIKE change % vs the reference (subsample)
\newcommand{\rhabXgbFullLiveSentCloseDm}{$-$0.61}  % DM vs the reference arm (subsample)
\newcommand{\rhabXgbFullLiveSentCloseP}{0.541}  % p-value
\newcommand{\rhabXgbHarLiveSentCloseQlike}{0.1025}  % QLIKE xgb har live16_sent (subsample)
\newcommand{\rhabRfFullLiveSentCloseQlike}{0.1026}  % QLIKE rf full live16_sent (subsample)
\newcommand{\rhabRfFullLiveSentClosePct}{$-$2.6}  % QLIKE change % vs the reference (subsample)
\newcommand{\rhabRfFullLiveSentCloseDm}{$-$0.65}  % DM vs the reference arm (subsample)
\newcommand{\rhabRfFullLiveSentCloseP}{0.513}  % p-value
\newcommand{\rhabRfHarLiveSentCloseQlike}{0.1053}  % QLIKE rf har live16_sent (subsample)
\newcommand{\rhabRidgeFullLiveIvQlike}{0.1940}  % QLIKE ridge full live_iv (subsample)
\newcommand{\rhabRidgeFullLiveIvPct}{$-$4.7}  % QLIKE change % vs the reference (subsample)
\newcommand{\rhabRidgeFullLiveIvDm}{$-$3.62}  % DM vs the reference arm (subsample)
\newcommand{\rhabRidgeFullLiveIvP}{\ensuremath{<}0.001}  % p-value
\newcommand{\rhabRidgeLoIvLiveIvQlike}{0.1942}  % QLIKE ridge lo_iv live_iv (subsample)
\newcommand{\rhabRidgeLoIvLiveIvPct}{0.1}  % QLIKE change % vs the reference (subsample)
\newcommand{\rhabRidgeLoIvLiveIvDm}{0.25}  % DM vs the reference arm (subsample)
\newcommand{\rhabRidgeLoIvLiveIvP}{0.806}  % p-value
\newcommand{\rhabRidgeAddIvLiveIvQlike}{0.2037}  % QLIKE ridge add_iv live_iv (subsample)
\newcommand{\rhabRidgeAddIvLiveIvPct}{0.0}  % QLIKE change % vs the reference (subsample)
\newcommand{\rhabRidgeAddIvLiveIvDm}{0.02}  % DM vs the reference arm (subsample)
\newcommand{\rhabRidgeAddIvLiveIvP}{0.986}  % p-value
\newcommand{\rhabRidgeHarLiveIvQlike}{0.2037}  % QLIKE ridge har live_iv (subsample)
\newcommand{\rhabLassoFullLiveIvQlike}{0.1924}  % QLIKE lasso full live_iv (subsample)
\newcommand{\rhabLassoFullLiveIvPct}{$-$5.5}  % QLIKE change % vs the reference (subsample)
\newcommand{\rhabLassoFullLiveIvDm}{$-$2.98}  % DM vs the reference arm (subsample)
\newcommand{\rhabLassoFullLiveIvP}{0.003}  % p-value
\newcommand{\rhabLassoLoIvLiveIvQlike}{0.1931}  % QLIKE lasso lo_iv live_iv (subsample)
\newcommand{\rhabLassoLoIvLiveIvPct}{0.4}  % QLIKE change % vs the reference (subsample)
\newcommand{\rhabLassoLoIvLiveIvDm}{0.47}  % DM vs the reference arm (subsample)
\newcommand{\rhabLassoLoIvLiveIvP}{0.641}  % p-value
\newcommand{\rhabLassoAddIvLiveIvQlike}{0.2019}  % QLIKE lasso add_iv live_iv (subsample)
\newcommand{\rhabLassoAddIvLiveIvPct}{$-$0.9}  % QLIKE change % vs the reference (subsample)
\newcommand{\rhabLassoAddIvLiveIvDm}{$-$1.39}  % DM vs the reference arm (subsample)
\newcommand{\rhabLassoAddIvLiveIvP}{0.165}  % p-value
\newcommand{\rhabLassoHarLiveIvQlike}{0.2037}  % QLIKE lasso har live_iv (subsample)
\newcommand{\rhabEnetFullLiveIvQlike}{0.1921}  % QLIKE enet full live_iv (subsample)
\newcommand{\rhabEnetFullLiveIvPct}{$-$5.7}  % QLIKE change % vs the reference (subsample)
\newcommand{\rhabEnetFullLiveIvDm}{$-$4.87}  % DM vs the reference arm (subsample)
\newcommand{\rhabEnetFullLiveIvP}{\ensuremath{<}0.001}  % p-value
\newcommand{\rhabEnetLoIvLiveIvQlike}{0.1924}  % QLIKE enet lo_iv live_iv (subsample)
\newcommand{\rhabEnetLoIvLiveIvPct}{0.1}  % QLIKE change % vs the reference (subsample)
\newcommand{\rhabEnetLoIvLiveIvDm}{0.23}  % DM vs the reference arm (subsample)
\newcommand{\rhabEnetLoIvLiveIvP}{0.816}  % p-value
\newcommand{\rhabEnetAddIvLiveIvQlike}{0.2033}  % QLIKE enet add_iv live_iv (subsample)
\newcommand{\rhabEnetAddIvLiveIvPct}{$-$0.2}  % QLIKE change % vs the reference (subsample)
\newcommand{\rhabEnetAddIvLiveIvDm}{$-$0.36}  % DM vs the reference arm (subsample)
\newcommand{\rhabEnetAddIvLiveIvP}{0.722}  % p-value
\newcommand{\rhabEnetHarLiveIvQlike}{0.2037}  % QLIKE enet har live_iv (subsample)
\newcommand{\rhabLgbmFullLiveIvQlike}{0.1739}  % QLIKE lgbm full live_iv (subsample)
\newcommand{\rhabLgbmFullLiveIvPct}{$-$13.3}  % QLIKE change % vs the reference (subsample)
\newcommand{\rhabLgbmFullLiveIvDm}{$-$7.31}  % DM vs the reference arm (subsample)
\newcommand{\rhabLgbmFullLiveIvP}{\ensuremath{<}0.001}  % p-value
\newcommand{\rhabLgbmHarLiveIvQlike}{0.2006}  % QLIKE lgbm har live_iv (subsample)
\newcommand{\rhabXgbFullLiveIvQlike}{0.1739}  % QLIKE xgb full live_iv (subsample)
\newcommand{\rhabXgbFullLiveIvPct}{$-$14.1}  % QLIKE change % vs the reference (subsample)
\newcommand{\rhabXgbFullLiveIvDm}{$-$7.16}  % DM vs the reference arm (subsample)
\newcommand{\rhabXgbFullLiveIvP}{\ensuremath{<}0.001}  % p-value
\newcommand{\rhabXgbHarLiveIvQlike}{0.2024}  % QLIKE xgb har live_iv (subsample)
\newcommand{\rhabRfFullLiveIvQlike}{0.1818}  % QLIKE rf full live_iv (subsample)
\newcommand{\rhabRfFullLiveIvPct}{$-$10.4}  % QLIKE change % vs the reference (subsample)
\newcommand{\rhabRfFullLiveIvDm}{$-$7.47}  % DM vs the reference arm (subsample)
\newcommand{\rhabRfFullLiveIvP}{\ensuremath{<}0.001}  % p-value
\newcommand{\rhabRfHarLiveIvQlike}{0.2029}  % QLIKE rf har live_iv (subsample)
\newcommand{\rhabRidgeFullLiveIvCloseQlike}{0.1380}  % QLIKE ridge full live16_iv (subsample)
\newcommand{\rhabRidgeFullLiveIvClosePct}{6.0}  % QLIKE change % vs the reference (subsample)
\newcommand{\rhabRidgeFullLiveIvCloseDm}{1.50}  % DM vs the reference arm (subsample)
\newcommand{\rhabRidgeFullLiveIvCloseP}{0.134}  % p-value
\newcommand{\rhabRidgeLoIvLiveIvCloseQlike}{0.1355}  % QLIKE ridge lo_iv live16_iv (subsample)
\newcommand{\rhabRidgeLoIvLiveIvClosePct}{$-$1.8}  % QLIKE change % vs the reference (subsample)
\newcommand{\rhabRidgeLoIvLiveIvCloseDm}{$-$1.22}  % DM vs the reference arm (subsample)
\newcommand{\rhabRidgeLoIvLiveIvCloseP}{0.222}  % p-value
\newcommand{\rhabRidgeAddIvLiveIvCloseQlike}{0.1306}  % QLIKE ridge add_iv live16_iv (subsample)
\newcommand{\rhabRidgeAddIvLiveIvClosePct}{0.3}  % QLIKE change % vs the reference (subsample)
\newcommand{\rhabRidgeAddIvLiveIvCloseDm}{0.19}  % DM vs the reference arm (subsample)
\newcommand{\rhabRidgeAddIvLiveIvCloseP}{0.853}  % p-value
\newcommand{\rhabRidgeHarLiveIvCloseQlike}{0.1302}  % QLIKE ridge har live16_iv (subsample)
\newcommand{\rhabLassoFullLiveIvCloseQlike}{0.1425}  % QLIKE lasso full live16_iv (subsample)
\newcommand{\rhabLassoFullLiveIvClosePct}{9.4}  % QLIKE change % vs the reference (subsample)
\newcommand{\rhabLassoFullLiveIvCloseDm}{1.93}  % DM vs the reference arm (subsample)
\newcommand{\rhabLassoFullLiveIvCloseP}{0.054}  % p-value
\newcommand{\rhabLassoLoIvLiveIvCloseQlike}{0.1378}  % QLIKE lasso lo_iv live16_iv (subsample)
\newcommand{\rhabLassoLoIvLiveIvClosePct}{$-$3.3}  % QLIKE change % vs the reference (subsample)
\newcommand{\rhabLassoLoIvLiveIvCloseDm}{$-$1.78}  % DM vs the reference arm (subsample)
\newcommand{\rhabLassoLoIvLiveIvCloseP}{0.075}  % p-value
\newcommand{\rhabLassoAddIvLiveIvCloseQlike}{0.1300}  % QLIKE lasso add_iv live16_iv (subsample)
\newcommand{\rhabLassoAddIvLiveIvClosePct}{$-$0.1}  % QLIKE change % vs the reference (subsample)
\newcommand{\rhabLassoAddIvLiveIvCloseDm}{$-$0.08}  % DM vs the reference arm (subsample)
\newcommand{\rhabLassoAddIvLiveIvCloseP}{0.940}  % p-value
\newcommand{\rhabLassoHarLiveIvCloseQlike}{0.1302}  % QLIKE lasso har live16_iv (subsample)
\newcommand{\rhabEnetFullLiveIvCloseQlike}{0.1355}  % QLIKE enet full live16_iv (subsample)
\newcommand{\rhabEnetFullLiveIvClosePct}{4.1}  % QLIKE change % vs the reference (subsample)
\newcommand{\rhabEnetFullLiveIvCloseDm}{1.12}  % DM vs the reference arm (subsample)
\newcommand{\rhabEnetFullLiveIvCloseP}{0.262}  % p-value
\newcommand{\rhabEnetLoIvLiveIvCloseQlike}{0.1321}  % QLIKE enet lo_iv live16_iv (subsample)
\newcommand{\rhabEnetLoIvLiveIvClosePct}{$-$2.5}  % QLIKE change % vs the reference (subsample)
\newcommand{\rhabEnetLoIvLiveIvCloseDm}{$-$1.58}  % DM vs the reference arm (subsample)
\newcommand{\rhabEnetLoIvLiveIvCloseP}{0.115}  % p-value
\newcommand{\rhabEnetAddIvLiveIvCloseQlike}{0.1308}  % QLIKE enet add_iv live16_iv (subsample)
\newcommand{\rhabEnetAddIvLiveIvClosePct}{0.5}  % QLIKE change % vs the reference (subsample)
\newcommand{\rhabEnetAddIvLiveIvCloseDm}{0.31}  % DM vs the reference arm (subsample)
\newcommand{\rhabEnetAddIvLiveIvCloseP}{0.756}  % p-value
\newcommand{\rhabEnetHarLiveIvCloseQlike}{0.1302}  % QLIKE enet har live16_iv (subsample)
\newcommand{\rhabLgbmFullLiveIvCloseQlike}{0.1242}  % QLIKE lgbm full live16_iv (subsample)
\newcommand{\rhabLgbmFullLiveIvClosePct}{$-$1.6}  % QLIKE change % vs the reference (subsample)
\newcommand{\rhabLgbmFullLiveIvCloseDm}{$-$0.74}  % DM vs the reference arm (subsample)
\newcommand{\rhabLgbmFullLiveIvCloseP}{0.461}  % p-value
\newcommand{\rhabLgbmHarLiveIvCloseQlike}{0.1263}  % QLIKE lgbm har live16_iv (subsample)
\newcommand{\rhabXgbFullLiveIvCloseQlike}{0.1234}  % QLIKE xgb full live16_iv (subsample)
\newcommand{\rhabXgbFullLiveIvClosePct}{$-$2.3}  % QLIKE change % vs the reference (subsample)
\newcommand{\rhabXgbFullLiveIvCloseDm}{$-$1.01}  % DM vs the reference arm (subsample)
\newcommand{\rhabXgbFullLiveIvCloseP}{0.312}  % p-value
\newcommand{\rhabXgbHarLiveIvCloseQlike}{0.1263}  % QLIKE xgb har live16_iv (subsample)
\newcommand{\rhabRfFullLiveIvCloseQlike}{0.1304}  % QLIKE rf full live16_iv (subsample)
\newcommand{\rhabRfFullLiveIvClosePct}{$-$0.4}  % QLIKE change % vs the reference (subsample)
\newcommand{\rhabRfFullLiveIvCloseDm}{$-$0.16}  % DM vs the reference arm (subsample)
\newcommand{\rhabRfFullLiveIvCloseP}{0.870}  % p-value
\newcommand{\rhabRfHarLiveIvCloseQlike}{0.1309}  % QLIKE rf har live16_iv (subsample)
\newcommand{\rhabRidgeFullLiveIvbQlike}{0.1757}  % QLIKE ridge full live_ivb (subsample)
\newcommand{\rhabRidgeFullLiveIvbPct}{$-$3.6}  % QLIKE change % vs the reference (subsample)
\newcommand{\rhabRidgeFullLiveIvbDm}{$-$2.12}  % DM vs the reference arm (subsample)
\newcommand{\rhabRidgeFullLiveIvbP}{0.034}  % p-value
\newcommand{\rhabRidgeLoIvLiveIvbQlike}{0.1746}  % QLIKE ridge lo_iv live_ivb (subsample)
\newcommand{\rhabRidgeLoIvLiveIvbPct}{$-$0.6}  % QLIKE change % vs the reference (subsample)
\newcommand{\rhabRidgeLoIvLiveIvbDm}{$-$0.79}  % DM vs the reference arm (subsample)
\newcommand{\rhabRidgeLoIvLiveIvbP}{0.431}  % p-value
\newcommand{\rhabRidgeAddIvLiveIvbQlike}{0.1850}  % QLIKE ridge add_iv live_ivb (subsample)
\newcommand{\rhabRidgeAddIvLiveIvbPct}{1.4}  % QLIKE change % vs the reference (subsample)
\newcommand{\rhabRidgeAddIvLiveIvbDm}{1.50}  % DM vs the reference arm (subsample)
\newcommand{\rhabRidgeAddIvLiveIvbP}{0.134}  % p-value
\newcommand{\rhabRidgeHarLiveIvbQlike}{0.1824}  % QLIKE ridge har live_ivb (subsample)
\newcommand{\rhabLassoFullLiveIvbQlike}{0.1718}  % QLIKE lasso full live_ivb (subsample)
\newcommand{\rhabLassoFullLiveIvbPct}{$-$5.8}  % QLIKE change % vs the reference (subsample)
\newcommand{\rhabLassoFullLiveIvbDm}{$-$2.58}  % DM vs the reference arm (subsample)
\newcommand{\rhabLassoFullLiveIvbP}{0.010}  % p-value
\newcommand{\rhabLassoLoIvLiveIvbQlike}{0.1700}  % QLIKE lasso lo_iv live_ivb (subsample)
\newcommand{\rhabLassoLoIvLiveIvbPct}{$-$1.0}  % QLIKE change % vs the reference (subsample)
\newcommand{\rhabLassoLoIvLiveIvbDm}{$-$0.96}  % DM vs the reference arm (subsample)
\newcommand{\rhabLassoLoIvLiveIvbP}{0.339}  % p-value
\newcommand{\rhabLassoAddIvLiveIvbQlike}{0.1844}  % QLIKE lasso add_iv live_ivb (subsample)
\newcommand{\rhabLassoAddIvLiveIvbPct}{1.1}  % QLIKE change % vs the reference (subsample)
\newcommand{\rhabLassoAddIvLiveIvbDm}{1.04}  % DM vs the reference arm (subsample)
\newcommand{\rhabLassoAddIvLiveIvbP}{0.300}  % p-value
\newcommand{\rhabLassoHarLiveIvbQlike}{0.1824}  % QLIKE lasso har live_ivb (subsample)
\newcommand{\rhabEnetFullLiveIvbQlike}{0.1745}  % QLIKE enet full live_ivb (subsample)
\newcommand{\rhabEnetFullLiveIvbPct}{$-$4.3}  % QLIKE change % vs the reference (subsample)
\newcommand{\rhabEnetFullLiveIvbDm}{$-$2.67}  % DM vs the reference arm (subsample)
\newcommand{\rhabEnetFullLiveIvbP}{0.008}  % p-value
\newcommand{\rhabEnetLoIvLiveIvbQlike}{0.1732}  % QLIKE enet lo_iv live_ivb (subsample)
\newcommand{\rhabEnetLoIvLiveIvbPct}{$-$0.7}  % QLIKE change % vs the reference (subsample)
\newcommand{\rhabEnetLoIvLiveIvbDm}{$-$0.88}  % DM vs the reference arm (subsample)
\newcommand{\rhabEnetLoIvLiveIvbP}{0.377}  % p-value
\newcommand{\rhabEnetAddIvLiveIvbQlike}{0.1844}  % QLIKE enet add_iv live_ivb (subsample)
\newcommand{\rhabEnetAddIvLiveIvbPct}{1.1}  % QLIKE change % vs the reference (subsample)
\newcommand{\rhabEnetAddIvLiveIvbDm}{1.30}  % DM vs the reference arm (subsample)
\newcommand{\rhabEnetAddIvLiveIvbP}{0.195}  % p-value
\newcommand{\rhabEnetHarLiveIvbQlike}{0.1824}  % QLIKE enet har live_ivb (subsample)
\newcommand{\rhabLgbmFullLiveIvbQlike}{0.1518}  % QLIKE lgbm full live_ivb (subsample)
\newcommand{\rhabLgbmFullLiveIvbPct}{$-$14.6}  % QLIKE change % vs the reference (subsample)
\newcommand{\rhabLgbmFullLiveIvbDm}{$-$5.14}  % DM vs the reference arm (subsample)
\newcommand{\rhabLgbmFullLiveIvbP}{\ensuremath{<}0.001}  % p-value
\newcommand{\rhabLgbmHarLiveIvbQlike}{0.1778}  % QLIKE lgbm har live_ivb (subsample)
\newcommand{\rhabXgbFullLiveIvbQlike}{0.1511}  % QLIKE xgb full live_ivb (subsample)
\newcommand{\rhabXgbFullLiveIvbPct}{$-$16.3}  % QLIKE change % vs the reference (subsample)
\newcommand{\rhabXgbFullLiveIvbDm}{$-$4.98}  % DM vs the reference arm (subsample)
\newcommand{\rhabXgbFullLiveIvbP}{\ensuremath{<}0.001}  % p-value
\newcommand{\rhabXgbHarLiveIvbQlike}{0.1805}  % QLIKE xgb har live_ivb (subsample)
\newcommand{\rhabRfFullLiveIvbQlike}{0.1621}  % QLIKE rf full live_ivb (subsample)
\newcommand{\rhabRfFullLiveIvbPct}{$-$11.3}  % QLIKE change % vs the reference (subsample)
\newcommand{\rhabRfFullLiveIvbDm}{$-$5.43}  % DM vs the reference arm (subsample)
\newcommand{\rhabRfFullLiveIvbP}{\ensuremath{<}0.001}  % p-value
\newcommand{\rhabRfHarLiveIvbQlike}{0.1829}  % QLIKE rf har live_ivb (subsample)
\newcommand{\rhabRidgeFullLiveIvbCloseQlike}{0.1180}  % QLIKE ridge full live16_ivb (subsample)
\newcommand{\rhabRidgeFullLiveIvbClosePct}{15.1}  % QLIKE change % vs the reference (subsample)
\newcommand{\rhabRidgeFullLiveIvbCloseDm}{2.54}  % DM vs the reference arm (subsample)
\newcommand{\rhabRidgeFullLiveIvbCloseP}{0.011}  % p-value
\newcommand{\rhabRidgeLoIvLiveIvbCloseQlike}{0.1123}  % QLIKE ridge lo_iv live16_ivb (subsample)
\newcommand{\rhabRidgeLoIvLiveIvbClosePct}{$-$4.9}  % QLIKE change % vs the reference (subsample)
\newcommand{\rhabRidgeLoIvLiveIvbCloseDm}{$-$1.73}  % DM vs the reference arm (subsample)
\newcommand{\rhabRidgeLoIvLiveIvbCloseP}{0.083}  % p-value
\newcommand{\rhabRidgeAddIvLiveIvbCloseQlike}{0.1073}  % QLIKE ridge add_iv live16_ivb (subsample)
\newcommand{\rhabRidgeAddIvLiveIvbClosePct}{4.6}  % QLIKE change % vs the reference (subsample)
\newcommand{\rhabRidgeAddIvLiveIvbCloseDm}{1.47}  % DM vs the reference arm (subsample)
\newcommand{\rhabRidgeAddIvLiveIvbCloseP}{0.142}  % p-value
\newcommand{\rhabRidgeHarLiveIvbCloseQlike}{0.1026}  % QLIKE ridge har live16_ivb (subsample)
\newcommand{\rhabLassoFullLiveIvbCloseQlike}{0.1202}  % QLIKE lasso full live16_ivb (subsample)
\newcommand{\rhabLassoFullLiveIvbClosePct}{17.2}  % QLIKE change % vs the reference (subsample)
\newcommand{\rhabLassoFullLiveIvbCloseDm}{2.74}  % DM vs the reference arm (subsample)
\newcommand{\rhabLassoFullLiveIvbCloseP}{0.006}  % p-value
\newcommand{\rhabLassoLoIvLiveIvbCloseQlike}{0.1120}  % QLIKE lasso lo_iv live16_ivb (subsample)
\newcommand{\rhabLassoLoIvLiveIvbClosePct}{$-$6.8}  % QLIKE change % vs the reference (subsample)
\newcommand{\rhabLassoLoIvLiveIvbCloseDm}{$-$2.01}  % DM vs the reference arm (subsample)
\newcommand{\rhabLassoLoIvLiveIvbCloseP}{0.044}  % p-value
\newcommand{\rhabLassoAddIvLiveIvbCloseQlike}{0.1066}  % QLIKE lasso add_iv live16_ivb (subsample)
\newcommand{\rhabLassoAddIvLiveIvbClosePct}{3.9}  % QLIKE change % vs the reference (subsample)
\newcommand{\rhabLassoAddIvLiveIvbCloseDm}{1.22}  % DM vs the reference arm (subsample)
\newcommand{\rhabLassoAddIvLiveIvbCloseP}{0.221}  % p-value
\newcommand{\rhabLassoHarLiveIvbCloseQlike}{0.1026}  % QLIKE lasso har live16_ivb (subsample)
\newcommand{\rhabEnetFullLiveIvbCloseQlike}{0.1147}  % QLIKE enet full live16_ivb (subsample)
\newcommand{\rhabEnetFullLiveIvbClosePct}{11.8}  % QLIKE change % vs the reference (subsample)
\newcommand{\rhabEnetFullLiveIvbCloseDm}{2.03}  % DM vs the reference arm (subsample)
\newcommand{\rhabEnetFullLiveIvbCloseP}{0.043}  % p-value
\newcommand{\rhabEnetLoIvLiveIvbCloseQlike}{0.1073}  % QLIKE enet lo_iv live16_ivb (subsample)
\newcommand{\rhabEnetLoIvLiveIvbClosePct}{$-$6.5}  % QLIKE change % vs the reference (subsample)
\newcommand{\rhabEnetLoIvLiveIvbCloseDm}{$-$2.14}  % DM vs the reference arm (subsample)
\newcommand{\rhabEnetLoIvLiveIvbCloseP}{0.033}  % p-value
\newcommand{\rhabEnetAddIvLiveIvbCloseQlike}{0.1076}  % QLIKE enet add_iv live16_ivb (subsample)
\newcommand{\rhabEnetAddIvLiveIvbClosePct}{4.9}  % QLIKE change % vs the reference (subsample)
\newcommand{\rhabEnetAddIvLiveIvbCloseDm}{1.60}  % DM vs the reference arm (subsample)
\newcommand{\rhabEnetAddIvLiveIvbCloseP}{0.110}  % p-value
\newcommand{\rhabEnetHarLiveIvbCloseQlike}{0.1026}  % QLIKE enet har live16_ivb (subsample)
\newcommand{\rhabLgbmFullLiveIvbCloseQlike}{0.0964}  % QLIKE lgbm full live16_ivb (subsample)
\newcommand{\rhabLgbmFullLiveIvbClosePct}{$-$1.0}  % QLIKE change % vs the reference (subsample)
\newcommand{\rhabLgbmFullLiveIvbCloseDm}{$-$0.29}  % DM vs the reference arm (subsample)
\newcommand{\rhabLgbmFullLiveIvbCloseP}{0.775}  % p-value
\newcommand{\rhabLgbmHarLiveIvbCloseQlike}{0.0974}  % QLIKE lgbm har live16_ivb (subsample)
\newcommand{\rhabXgbFullLiveIvbCloseQlike}{0.0962}  % QLIKE xgb full live16_ivb (subsample)
\newcommand{\rhabXgbFullLiveIvbClosePct}{$-$1.5}  % QLIKE change % vs the reference (subsample)
\newcommand{\rhabXgbFullLiveIvbCloseDm}{$-$0.40}  % DM vs the reference arm (subsample)
\newcommand{\rhabXgbFullLiveIvbCloseP}{0.689}  % p-value
\newcommand{\rhabXgbHarLiveIvbCloseQlike}{0.0976}  % QLIKE xgb har live16_ivb (subsample)
\newcommand{\rhabRfFullLiveIvbCloseQlike}{0.1005}  % QLIKE rf full live16_ivb (subsample)
\newcommand{\rhabRfFullLiveIvbClosePct}{0.2}  % QLIKE change % vs the reference (subsample)
\newcommand{\rhabRfFullLiveIvbCloseDm}{0.05}  % DM vs the reference arm (subsample)
\newcommand{\rhabRfFullLiveIvbCloseP}{0.958}  % p-value
\newcommand{\rhabRfHarLiveIvbCloseQlike}{0.1003}  % QLIKE rf har live16_ivb (subsample)
\newcommand{\rhabRidgeFullLiveVdQlike}{0.1942}  % QLIKE ridge full live_vd (subsample)
\newcommand{\rhabRidgeFullLiveVdPct}{4.2}  % QLIKE change % vs the reference (subsample)
\newcommand{\rhabRidgeFullLiveVdDm}{1.33}  % DM vs the reference arm (subsample)
\newcommand{\rhabRidgeFullLiveVdP}{0.183}  % p-value
\newcommand{\rhabRidgeLoVdLiveVdQlike}{0.1887}  % QLIKE ridge lo_vd live_vd (subsample)
\newcommand{\rhabRidgeLoVdLiveVdPct}{$-$2.8}  % QLIKE change % vs the reference (subsample)
\newcommand{\rhabRidgeLoVdLiveVdDm}{$-$3.17}  % DM vs the reference arm (subsample)
\newcommand{\rhabRidgeLoVdLiveVdP}{0.002}  % p-value
\newcommand{\rhabRidgeAddVdLiveVdQlike}{0.1944}  % QLIKE ridge add_vd live_vd (subsample)
\newcommand{\rhabRidgeAddVdLiveVdPct}{4.3}  % QLIKE change % vs the reference (subsample)
\newcommand{\rhabRidgeAddVdLiveVdDm}{2.52}  % DM vs the reference arm (subsample)
\newcommand{\rhabRidgeAddVdLiveVdP}{0.012}  % p-value
\newcommand{\rhabRidgeHarLiveVdQlike}{0.1864}  % QLIKE ridge har live_vd (subsample)
\newcommand{\rhabLassoFullLiveVdQlike}{0.1898}  % QLIKE lasso full live_vd (subsample)
\newcommand{\rhabLassoFullLiveVdPct}{1.9}  % QLIKE change % vs the reference (subsample)
\newcommand{\rhabLassoFullLiveVdDm}{0.45}  % DM vs the reference arm (subsample)
\newcommand{\rhabLassoFullLiveVdP}{0.651}  % p-value
\newcommand{\rhabLassoLoVdLiveVdQlike}{0.1846}  % QLIKE lasso lo_vd live_vd (subsample)
\newcommand{\rhabLassoLoVdLiveVdPct}{$-$2.7}  % QLIKE change % vs the reference (subsample)
\newcommand{\rhabLassoLoVdLiveVdDm}{$-$2.51}  % DM vs the reference arm (subsample)
\newcommand{\rhabLassoLoVdLiveVdP}{0.012}  % p-value
\newcommand{\rhabLassoAddVdLiveVdQlike}{0.1941}  % QLIKE lasso add_vd live_vd (subsample)
\newcommand{\rhabLassoAddVdLiveVdPct}{4.2}  % QLIKE change % vs the reference (subsample)
\newcommand{\rhabLassoAddVdLiveVdDm}{2.49}  % DM vs the reference arm (subsample)
\newcommand{\rhabLassoAddVdLiveVdP}{0.013}  % p-value
\newcommand{\rhabLassoHarLiveVdQlike}{0.1864}  % QLIKE lasso har live_vd (subsample)
\newcommand{\rhabEnetFullLiveVdQlike}{0.1908}  % QLIKE enet full live_vd (subsample)
\newcommand{\rhabEnetFullLiveVdPct}{2.4}  % QLIKE change % vs the reference (subsample)
\newcommand{\rhabEnetFullLiveVdDm}{0.85}  % DM vs the reference arm (subsample)
\newcommand{\rhabEnetFullLiveVdP}{0.397}  % p-value
\newcommand{\rhabEnetLoVdLiveVdQlike}{0.1869}  % QLIKE enet lo_vd live_vd (subsample)
\newcommand{\rhabEnetLoVdLiveVdPct}{$-$2.1}  % QLIKE change % vs the reference (subsample)
\newcommand{\rhabEnetLoVdLiveVdDm}{$-$3.47}  % DM vs the reference arm (subsample)
\newcommand{\rhabEnetLoVdLiveVdP}{\ensuremath{<}0.001}  % p-value
\newcommand{\rhabEnetAddVdLiveVdQlike}{0.1916}  % QLIKE enet add_vd live_vd (subsample)
\newcommand{\rhabEnetAddVdLiveVdPct}{2.8}  % QLIKE change % vs the reference (subsample)
\newcommand{\rhabEnetAddVdLiveVdDm}{2.20}  % DM vs the reference arm (subsample)
\newcommand{\rhabEnetAddVdLiveVdP}{0.028}  % p-value
\newcommand{\rhabEnetHarLiveVdQlike}{0.1864}  % QLIKE enet har live_vd (subsample)
\newcommand{\rhabLgbmFullLiveVdQlike}{0.1612}  % QLIKE lgbm full live_vd (subsample)
\newcommand{\rhabLgbmFullLiveVdPct}{$-$10.9}  % QLIKE change % vs the reference (subsample)
\newcommand{\rhabLgbmFullLiveVdDm}{$-$2.01}  % DM vs the reference arm (subsample)
\newcommand{\rhabLgbmFullLiveVdP}{0.044}  % p-value
\newcommand{\rhabLgbmHarLiveVdQlike}{0.1809}  % QLIKE lgbm har live_vd (subsample)
\newcommand{\rhabXgbFullLiveVdQlike}{0.1616}  % QLIKE xgb full live_vd (subsample)
\newcommand{\rhabXgbFullLiveVdPct}{$-$12.1}  % QLIKE change % vs the reference (subsample)
\newcommand{\rhabXgbFullLiveVdDm}{$-$2.10}  % DM vs the reference arm (subsample)
\newcommand{\rhabXgbFullLiveVdP}{0.036}  % p-value
\newcommand{\rhabXgbHarLiveVdQlike}{0.1839}  % QLIKE xgb har live_vd (subsample)
\newcommand{\rhabRfFullLiveVdQlike}{0.1704}  % QLIKE rf full live_vd (subsample)
\newcommand{\rhabRfFullLiveVdPct}{$-$9.1}  % QLIKE change % vs the reference (subsample)
\newcommand{\rhabRfFullLiveVdDm}{$-$2.40}  % DM vs the reference arm (subsample)
\newcommand{\rhabRfFullLiveVdP}{0.016}  % p-value
\newcommand{\rhabRfHarLiveVdQlike}{0.1874}  % QLIKE rf har live_vd (subsample)
\newcommand{\rhabRidgeFullLiveVdCloseQlike}{0.1269}  % QLIKE ridge full live16_vd (subsample)
\newcommand{\rhabRidgeFullLiveVdClosePct}{45.8}  % QLIKE change % vs the reference (subsample)
\newcommand{\rhabRidgeFullLiveVdCloseDm}{3.73}  % DM vs the reference arm (subsample)
\newcommand{\rhabRidgeFullLiveVdCloseP}{\ensuremath{<}0.001}  % p-value
\newcommand{\rhabRidgeLoVdLiveVdCloseQlike}{0.1048}  % QLIKE ridge lo_vd live16_vd (subsample)
\newcommand{\rhabRidgeLoVdLiveVdClosePct}{$-$17.4}  % QLIKE change % vs the reference (subsample)
\newcommand{\rhabRidgeLoVdLiveVdCloseDm}{$-$4.97}  % DM vs the reference arm (subsample)
\newcommand{\rhabRidgeLoVdLiveVdCloseP}{\ensuremath{<}0.001}  % p-value
\newcommand{\rhabRidgeAddVdLiveVdCloseQlike}{0.1018}  % QLIKE ridge add_vd live16_vd (subsample)
\newcommand{\rhabRidgeAddVdLiveVdClosePct}{17.0}  % QLIKE change % vs the reference (subsample)
\newcommand{\rhabRidgeAddVdLiveVdCloseDm}{3.66}  % DM vs the reference arm (subsample)
\newcommand{\rhabRidgeAddVdLiveVdCloseP}{\ensuremath{<}0.001}  % p-value
\newcommand{\rhabRidgeHarLiveVdCloseQlike}{0.0870}  % QLIKE ridge har live16_vd (subsample)
\newcommand{\rhabLassoFullLiveVdCloseQlike}{0.1271}  % QLIKE lasso full live16_vd (subsample)
\newcommand{\rhabLassoFullLiveVdClosePct}{46.0}  % QLIKE change % vs the reference (subsample)
\newcommand{\rhabLassoFullLiveVdCloseDm}{3.87}  % DM vs the reference arm (subsample)
\newcommand{\rhabLassoFullLiveVdCloseP}{\ensuremath{<}0.001}  % p-value
\newcommand{\rhabLassoLoVdLiveVdCloseQlike}{0.1097}  % QLIKE lasso lo_vd live16_vd (subsample)
\newcommand{\rhabLassoLoVdLiveVdClosePct}{$-$13.7}  % QLIKE change % vs the reference (subsample)
\newcommand{\rhabLassoLoVdLiveVdCloseDm}{$-$4.18}  % DM vs the reference arm (subsample)
\newcommand{\rhabLassoLoVdLiveVdCloseP}{\ensuremath{<}0.001}  % p-value
\newcommand{\rhabLassoAddVdLiveVdCloseQlike}{0.1018}  % QLIKE lasso add_vd live16_vd (subsample)
\newcommand{\rhabLassoAddVdLiveVdClosePct}{17.0}  % QLIKE change % vs the reference (subsample)
\newcommand{\rhabLassoAddVdLiveVdCloseDm}{3.75}  % DM vs the reference arm (subsample)
\newcommand{\rhabLassoAddVdLiveVdCloseP}{\ensuremath{<}0.001}  % p-value
\newcommand{\rhabLassoHarLiveVdCloseQlike}{0.0870}  % QLIKE lasso har live16_vd (subsample)
\newcommand{\rhabEnetFullLiveVdCloseQlike}{0.1227}  % QLIKE enet full live16_vd (subsample)
\newcommand{\rhabEnetFullLiveVdClosePct}{41.0}  % QLIKE change % vs the reference (subsample)
\newcommand{\rhabEnetFullLiveVdCloseDm}{3.59}  % DM vs the reference arm (subsample)
\newcommand{\rhabEnetFullLiveVdCloseP}{\ensuremath{<}0.001}  % p-value
\newcommand{\rhabEnetLoVdLiveVdCloseQlike}{0.1029}  % QLIKE enet lo_vd live16_vd (subsample)
\newcommand{\rhabEnetLoVdLiveVdClosePct}{$-$16.2}  % QLIKE change % vs the reference (subsample)
\newcommand{\rhabEnetLoVdLiveVdCloseDm}{$-$4.00}  % DM vs the reference arm (subsample)
\newcommand{\rhabEnetLoVdLiveVdCloseP}{\ensuremath{<}0.001}  % p-value
\newcommand{\rhabEnetAddVdLiveVdCloseQlike}{0.1005}  % QLIKE enet add_vd live16_vd (subsample)
\newcommand{\rhabEnetAddVdLiveVdClosePct}{15.5}  % QLIKE change % vs the reference (subsample)
\newcommand{\rhabEnetAddVdLiveVdCloseDm}{3.55}  % DM vs the reference arm (subsample)
\newcommand{\rhabEnetAddVdLiveVdCloseP}{\ensuremath{<}0.001}  % p-value
\newcommand{\rhabEnetHarLiveVdCloseQlike}{0.0870}  % QLIKE enet har live16_vd (subsample)
\newcommand{\rhabLgbmFullLiveVdCloseQlike}{0.0802}  % QLIKE lgbm full live16_vd (subsample)
\newcommand{\rhabLgbmFullLiveVdClosePct}{$-$1.5}  % QLIKE change % vs the reference (subsample)
\newcommand{\rhabLgbmFullLiveVdCloseDm}{$-$0.29}  % DM vs the reference arm (subsample)
\newcommand{\rhabLgbmFullLiveVdCloseP}{0.769}  % p-value
\newcommand{\rhabLgbmHarLiveVdCloseQlike}{0.0814}  % QLIKE lgbm har live16_vd (subsample)
\newcommand{\rhabXgbFullLiveVdCloseQlike}{0.0785}  % QLIKE xgb full live16_vd (subsample)
\newcommand{\rhabXgbFullLiveVdClosePct}{$-$4.3}  % QLIKE change % vs the reference (subsample)
\newcommand{\rhabXgbFullLiveVdCloseDm}{$-$0.89}  % DM vs the reference arm (subsample)
\newcommand{\rhabXgbFullLiveVdCloseP}{0.375}  % p-value
\newcommand{\rhabXgbHarLiveVdCloseQlike}{0.0820}  % QLIKE xgb har live16_vd (subsample)
\newcommand{\rhabRfFullLiveVdCloseQlike}{0.0822}  % QLIKE rf full live16_vd (subsample)
\newcommand{\rhabRfFullLiveVdClosePct}{$-$1.0}  % QLIKE change % vs the reference (subsample)
\newcommand{\rhabRfFullLiveVdCloseDm}{$-$0.21}  % DM vs the reference arm (subsample)
\newcommand{\rhabRfFullLiveVdCloseP}{0.833}  % p-value
\newcommand{\rhabRfHarLiveVdCloseQlike}{0.0831}  % QLIKE rf har live16_vd (subsample)
\newcommand{\rhabLiveSentStart}{19~April 2022}  % first bar of the live window sent
\newcommand{\rhabLiveSentEnd}{30~April 2024}  % last bar of the live window sent
\newcommand{\rhabLiveSentDays}{508}  % trading days of the live window sent
\newcommand{\rhabLiveIvStart}{19~April 2018}  % first bar of the live window iv
\newcommand{\rhabLiveIvEnd}{30~April 2024}  % last bar of the live window iv
\newcommand{\rhabLiveIvDays}{1{,}506}  % trading days of the live window iv
\newcommand{\rhabLiveIvbStart}{6~July 2021}  % first bar of the live window ivb
\newcommand{\rhabLiveIvbEnd}{30~April 2024}  % last bar of the live window ivb
\newcommand{\rhabLiveIvbDays}{706}  % trading days of the live window ivb
\newcommand{\rhabLiveVdStart}{16~November 2022}  % first bar of the live window vd
\newcommand{\rhabLiveVdEnd}{31~August 2023}  % last bar of the live window vd
\newcommand{\rhabLiveVdDays}{196}  % trading days of the live window vd
\newcommand{\rhabRidgeFullVsOlsAddRmAllDm}{$-$4.70}  % DM ridge full vs OLS add_rm
\newcommand{\rhabRidgeFullVsOlsAddRmAllP}{\ensuremath{<}0.001}  % p-value
\newcommand{\rhabRidgeFullVsOlsAddRmCloseDm}{1.16}  % DM ridge full vs OLS add_rm
\newcommand{\rhabRidgeFullVsOlsAddRmCloseP}{0.244}  % p-value
\newcommand{\rhabRidgeFullVsOlsAddRmCloseLateDm}{0.95}  % DM ridge full vs OLS add_rm
\newcommand{\rhabRidgeFullVsOlsAddRmCloseLateP}{0.341}  % p-value
\newcommand{\rhabRidgeFullVsOlsSharAllDm}{$-$4.77}  % DM ridge full vs OLS shar
\newcommand{\rhabRidgeFullVsOlsSharAllP}{\ensuremath{<}0.001}  % p-value
\newcommand{\rhabRidgeFullVsOlsSharCloseDm}{0.41}  % DM ridge full vs OLS shar
\newcommand{\rhabRidgeFullVsOlsSharCloseP}{0.682}  % p-value
\newcommand{\rhabRidgeFullVsOlsSharCloseLateDm}{2.34}  % DM ridge full vs OLS shar
\newcommand{\rhabRidgeFullVsOlsSharCloseLateP}{0.019}  % p-value
\newcommand{\rhabRidgeFullVsOlsHarqAllDm}{$-$5.66}  % DM ridge full vs OLS harq
\newcommand{\rhabRidgeFullVsOlsHarqAllP}{\ensuremath{<}0.001}  % p-value
\newcommand{\rhabRidgeFullVsOlsHarqCloseDm}{0.37}  % DM ridge full vs OLS harq
\newcommand{\rhabRidgeFullVsOlsHarqCloseP}{0.715}  % p-value
\newcommand{\rhabRidgeFullVsOlsHarqCloseLateDm}{2.35}  % DM ridge full vs OLS harq
\newcommand{\rhabRidgeFullVsOlsHarqCloseLateP}{0.019}  % p-value
\newcommand{\rhabRidgeFullVsOlsLharAllDm}{$-$3.23}  % DM ridge full vs OLS lhar
\newcommand{\rhabRidgeFullVsOlsLharAllP}{0.001}  % p-value
\newcommand{\rhabRidgeFullVsOlsLharCloseDm}{2.40}  % DM ridge full vs OLS lhar
\newcommand{\rhabRidgeFullVsOlsLharCloseP}{0.016}  % p-value
\newcommand{\rhabRidgeFullVsOlsLharCloseLateDm}{2.96}  % DM ridge full vs OLS lhar
\newcommand{\rhabRidgeFullVsOlsLharCloseLateP}{0.003}  % p-value
\newcommand{\rhabRidgeFullVsOlsSharqlAllDm}{$-$3.04}  % DM ridge full vs OLS sharql
\newcommand{\rhabRidgeFullVsOlsSharqlAllP}{0.002}  % p-value
\newcommand{\rhabRidgeFullVsOlsSharqlCloseDm}{2.44}  % DM ridge full vs OLS sharql
\newcommand{\rhabRidgeFullVsOlsSharqlCloseP}{0.015}  % p-value
\newcommand{\rhabRidgeFullVsOlsSharqlCloseLateDm}{3.13}  % DM ridge full vs OLS sharql
\newcommand{\rhabRidgeFullVsOlsSharqlCloseLateP}{0.002}  % p-value
\newcommand{\rhabRidgeFullVsOlsHarAllDm}{$-$5.63}  % DM ridge full vs OLS har
\newcommand{\rhabRidgeFullVsOlsHarAllP}{\ensuremath{<}0.001}  % p-value
\newcommand{\rhabRidgeFullVsOlsHarCloseDm}{0.31}  % DM ridge full vs OLS har
\newcommand{\rhabRidgeFullVsOlsHarCloseP}{0.759}  % p-value
\newcommand{\rhabRidgeFullVsOlsHarCloseLateDm}{2.31}  % DM ridge full vs OLS har
\newcommand{\rhabRidgeFullVsOlsHarCloseLateP}{0.021}  % p-value
\newcommand{\rhabLassoFullVsOlsAddRmAllDm}{$-$4.53}  % DM lasso full vs OLS add_rm
\newcommand{\rhabLassoFullVsOlsAddRmAllP}{\ensuremath{<}0.001}  % p-value
\newcommand{\rhabLassoFullVsOlsAddRmCloseDm}{2.40}  % DM lasso full vs OLS add_rm
\newcommand{\rhabLassoFullVsOlsAddRmCloseP}{0.016}  % p-value
\newcommand{\rhabLassoFullVsOlsAddRmCloseLateDm}{2.26}  % DM lasso full vs OLS add_rm
\newcommand{\rhabLassoFullVsOlsAddRmCloseLateP}{0.024}  % p-value
\newcommand{\rhabLassoFullVsOlsSharAllDm}{$-$4.07}  % DM lasso full vs OLS shar
\newcommand{\rhabLassoFullVsOlsSharAllP}{\ensuremath{<}0.001}  % p-value
\newcommand{\rhabLassoFullVsOlsSharCloseDm}{1.12}  % DM lasso full vs OLS shar
\newcommand{\rhabLassoFullVsOlsSharCloseP}{0.261}  % p-value
\newcommand{\rhabLassoFullVsOlsSharCloseLateDm}{2.54}  % DM lasso full vs OLS shar
\newcommand{\rhabLassoFullVsOlsSharCloseLateP}{0.011}  % p-value
\newcommand{\rhabLassoFullVsOlsHarqAllDm}{$-$4.75}  % DM lasso full vs OLS harq
\newcommand{\rhabLassoFullVsOlsHarqAllP}{\ensuremath{<}0.001}  % p-value
\newcommand{\rhabLassoFullVsOlsHarqCloseDm}{1.08}  % DM lasso full vs OLS harq
\newcommand{\rhabLassoFullVsOlsHarqCloseP}{0.279}  % p-value
\newcommand{\rhabLassoFullVsOlsHarqCloseLateDm}{2.55}  % DM lasso full vs OLS harq
\newcommand{\rhabLassoFullVsOlsHarqCloseLateP}{0.011}  % p-value
\newcommand{\rhabLassoFullVsOlsLharAllDm}{$-$2.95}  % DM lasso full vs OLS lhar
\newcommand{\rhabLassoFullVsOlsLharAllP}{0.003}  % p-value
\newcommand{\rhabLassoFullVsOlsLharCloseDm}{3.00}  % DM lasso full vs OLS lhar
\newcommand{\rhabLassoFullVsOlsLharCloseP}{0.003}  % p-value
\newcommand{\rhabLassoFullVsOlsLharCloseLateDm}{3.28}  % DM lasso full vs OLS lhar
\newcommand{\rhabLassoFullVsOlsLharCloseLateP}{0.001}  % p-value
\newcommand{\rhabLassoFullVsOlsSharqlAllDm}{$-$2.81}  % DM lasso full vs OLS sharql
\newcommand{\rhabLassoFullVsOlsSharqlAllP}{0.005}  % p-value
\newcommand{\rhabLassoFullVsOlsSharqlCloseDm}{2.96}  % DM lasso full vs OLS sharql
\newcommand{\rhabLassoFullVsOlsSharqlCloseP}{0.003}  % p-value
\newcommand{\rhabLassoFullVsOlsSharqlCloseLateDm}{3.33}  % DM lasso full vs OLS sharql
\newcommand{\rhabLassoFullVsOlsSharqlCloseLateP}{\ensuremath{<}0.001}  % p-value
\newcommand{\rhabLassoFullVsOlsHarAllDm}{$-$4.72}  % DM lasso full vs OLS har
\newcommand{\rhabLassoFullVsOlsHarAllP}{\ensuremath{<}0.001}  % p-value
\newcommand{\rhabLassoFullVsOlsHarCloseDm}{1.04}  % DM lasso full vs OLS har
\newcommand{\rhabLassoFullVsOlsHarCloseP}{0.300}  % p-value
\newcommand{\rhabLassoFullVsOlsHarCloseLateDm}{2.51}  % DM lasso full vs OLS har
\newcommand{\rhabLassoFullVsOlsHarCloseLateP}{0.012}  % p-value
\newcommand{\rhabEnetFullVsOlsAddRmAllDm}{$-$5.61}  % DM enet full vs OLS add_rm
\newcommand{\rhabEnetFullVsOlsAddRmAllP}{\ensuremath{<}0.001}  % p-value
\newcommand{\rhabEnetFullVsOlsAddRmCloseDm}{0.46}  % DM enet full vs OLS add_rm
\newcommand{\rhabEnetFullVsOlsAddRmCloseP}{0.649}  % p-value
\newcommand{\rhabEnetFullVsOlsAddRmCloseLateDm}{0.29}  % DM enet full vs OLS add_rm
\newcommand{\rhabEnetFullVsOlsAddRmCloseLateP}{0.770}  % p-value
\newcommand{\rhabEnetFullVsOlsSharAllDm}{$-$5.97}  % DM enet full vs OLS shar
\newcommand{\rhabEnetFullVsOlsSharAllP}{\ensuremath{<}0.001}  % p-value
\newcommand{\rhabEnetFullVsOlsSharCloseDm}{$-$0.11}  % DM enet full vs OLS shar
\newcommand{\rhabEnetFullVsOlsSharCloseP}{0.914}  % p-value
\newcommand{\rhabEnetFullVsOlsSharCloseLateDm}{1.92}  % DM enet full vs OLS shar
\newcommand{\rhabEnetFullVsOlsSharCloseLateP}{0.055}  % p-value
\newcommand{\rhabEnetFullVsOlsHarqAllDm}{$-$7.03}  % DM enet full vs OLS harq
\newcommand{\rhabEnetFullVsOlsHarqAllP}{\ensuremath{<}0.001}  % p-value
\newcommand{\rhabEnetFullVsOlsHarqCloseDm}{$-$0.15}  % DM enet full vs OLS harq
\newcommand{\rhabEnetFullVsOlsHarqCloseP}{0.883}  % p-value
\newcommand{\rhabEnetFullVsOlsHarqCloseLateDm}{1.96}  % DM enet full vs OLS harq
\newcommand{\rhabEnetFullVsOlsHarqCloseLateP}{0.050}  % p-value
\newcommand{\rhabEnetFullVsOlsLharAllDm}{$-$4.58}  % DM enet full vs OLS lhar
\newcommand{\rhabEnetFullVsOlsLharAllP}{\ensuremath{<}0.001}  % p-value
\newcommand{\rhabEnetFullVsOlsLharCloseDm}{1.86}  % DM enet full vs OLS lhar
\newcommand{\rhabEnetFullVsOlsLharCloseP}{0.063}  % p-value
\newcommand{\rhabEnetFullVsOlsLharCloseLateDm}{2.46}  % DM enet full vs OLS lhar
\newcommand{\rhabEnetFullVsOlsLharCloseLateP}{0.014}  % p-value
\newcommand{\rhabEnetFullVsOlsSharqlAllDm}{$-$4.34}  % DM enet full vs OLS sharql
\newcommand{\rhabEnetFullVsOlsSharqlAllP}{\ensuremath{<}0.001}  % p-value
\newcommand{\rhabEnetFullVsOlsSharqlCloseDm}{1.92}  % DM enet full vs OLS sharql
\newcommand{\rhabEnetFullVsOlsSharqlCloseP}{0.055}  % p-value
\newcommand{\rhabEnetFullVsOlsSharqlCloseLateDm}{2.65}  % DM enet full vs OLS sharql
\newcommand{\rhabEnetFullVsOlsSharqlCloseLateP}{0.008}  % p-value
\newcommand{\rhabEnetFullVsOlsHarAllDm}{$-$6.98}  % DM enet full vs OLS har
\newcommand{\rhabEnetFullVsOlsHarAllP}{\ensuremath{<}0.001}  % p-value
\newcommand{\rhabEnetFullVsOlsHarCloseDm}{$-$0.21}  % DM enet full vs OLS har
\newcommand{\rhabEnetFullVsOlsHarCloseP}{0.832}  % p-value
\newcommand{\rhabEnetFullVsOlsHarCloseLateDm}{1.91}  % DM enet full vs OLS har
\newcommand{\rhabEnetFullVsOlsHarCloseLateP}{0.056}  % p-value
\newcommand{\rhabLgbmFullVsOlsAddRmAllDm}{$-$7.94}  % DM lgbm full vs OLS add_rm
\newcommand{\rhabLgbmFullVsOlsAddRmAllP}{\ensuremath{<}0.001}  % p-value
\newcommand{\rhabLgbmFullVsOlsAddRmCloseDm}{$-$1.78}  % DM lgbm full vs OLS add_rm
\newcommand{\rhabLgbmFullVsOlsAddRmCloseP}{0.076}  % p-value
\newcommand{\rhabLgbmFullVsOlsAddRmCloseLateDm}{$-$2.44}  % DM lgbm full vs OLS add_rm
\newcommand{\rhabLgbmFullVsOlsAddRmCloseLateP}{0.015}  % p-value
\newcommand{\rhabLgbmFullVsOlsSharAllDm}{$-$9.26}  % DM lgbm full vs OLS shar
\newcommand{\rhabLgbmFullVsOlsSharAllP}{\ensuremath{<}0.001}  % p-value
\newcommand{\rhabLgbmFullVsOlsSharCloseDm}{$-$3.86}  % DM lgbm full vs OLS shar
\newcommand{\rhabLgbmFullVsOlsSharCloseP}{\ensuremath{<}0.001}  % p-value
\newcommand{\rhabLgbmFullVsOlsSharCloseLateDm}{$-$2.34}  % DM lgbm full vs OLS shar
\newcommand{\rhabLgbmFullVsOlsSharCloseLateP}{0.019}  % p-value
\newcommand{\rhabLgbmFullVsOlsHarqAllDm}{$-$10.44}  % DM lgbm full vs OLS harq
\newcommand{\rhabLgbmFullVsOlsHarqAllP}{\ensuremath{<}0.001}  % p-value
\newcommand{\rhabLgbmFullVsOlsHarqCloseDm}{$-$3.77}  % DM lgbm full vs OLS harq
\newcommand{\rhabLgbmFullVsOlsHarqCloseP}{\ensuremath{<}0.001}  % p-value
\newcommand{\rhabLgbmFullVsOlsHarqCloseLateDm}{$-$1.70}  % DM lgbm full vs OLS harq
\newcommand{\rhabLgbmFullVsOlsHarqCloseLateP}{0.088}  % p-value
\newcommand{\rhabLgbmFullVsOlsLharAllDm}{$-$8.32}  % DM lgbm full vs OLS lhar
\newcommand{\rhabLgbmFullVsOlsLharAllP}{\ensuremath{<}0.001}  % p-value
\newcommand{\rhabLgbmFullVsOlsLharCloseDm}{$-$1.02}  % DM lgbm full vs OLS lhar
\newcommand{\rhabLgbmFullVsOlsLharCloseP}{0.310}  % p-value
\newcommand{\rhabLgbmFullVsOlsLharCloseLateDm}{$-$1.42}  % DM lgbm full vs OLS lhar
\newcommand{\rhabLgbmFullVsOlsLharCloseLateP}{0.155}  % p-value
\newcommand{\rhabLgbmFullVsOlsSharqlAllDm}{$-$8.22}  % DM lgbm full vs OLS sharql
\newcommand{\rhabLgbmFullVsOlsSharqlAllP}{\ensuremath{<}0.001}  % p-value
\newcommand{\rhabLgbmFullVsOlsSharqlCloseDm}{$-$1.03}  % DM lgbm full vs OLS sharql
\newcommand{\rhabLgbmFullVsOlsSharqlCloseP}{0.302}  % p-value
\newcommand{\rhabLgbmFullVsOlsSharqlCloseLateDm}{$-$1.45}  % DM lgbm full vs OLS sharql
\newcommand{\rhabLgbmFullVsOlsSharqlCloseLateP}{0.148}  % p-value
\newcommand{\rhabLgbmFullVsOlsHarAllDm}{$-$10.41}  % DM lgbm full vs OLS har
\newcommand{\rhabLgbmFullVsOlsHarAllP}{\ensuremath{<}0.001}  % p-value
\newcommand{\rhabLgbmFullVsOlsHarCloseDm}{$-$3.89}  % DM lgbm full vs OLS har
\newcommand{\rhabLgbmFullVsOlsHarCloseP}{\ensuremath{<}0.001}  % p-value
\newcommand{\rhabLgbmFullVsOlsHarCloseLateDm}{$-$1.86}  % DM lgbm full vs OLS har
\newcommand{\rhabLgbmFullVsOlsHarCloseLateP}{0.063}  % p-value
\newcommand{\rhabXgbFullVsOlsAddRmAllDm}{$-$7.47}  % DM xgb full vs OLS add_rm
\newcommand{\rhabXgbFullVsOlsAddRmAllP}{\ensuremath{<}0.001}  % p-value
\newcommand{\rhabXgbFullVsOlsAddRmCloseDm}{$-$2.16}  % DM xgb full vs OLS add_rm
\newcommand{\rhabXgbFullVsOlsAddRmCloseP}{0.031}  % p-value
\newcommand{\rhabXgbFullVsOlsAddRmCloseLateDm}{$-$2.61}  % DM xgb full vs OLS add_rm
\newcommand{\rhabXgbFullVsOlsAddRmCloseLateP}{0.009}  % p-value
\newcommand{\rhabXgbFullVsOlsSharAllDm}{$-$8.81}  % DM xgb full vs OLS shar
\newcommand{\rhabXgbFullVsOlsSharAllP}{\ensuremath{<}0.001}  % p-value
\newcommand{\rhabXgbFullVsOlsSharCloseDm}{$-$4.87}  % DM xgb full vs OLS shar
\newcommand{\rhabXgbFullVsOlsSharCloseP}{\ensuremath{<}0.001}  % p-value
\newcommand{\rhabXgbFullVsOlsSharCloseLateDm}{$-$2.99}  % DM xgb full vs OLS shar
\newcommand{\rhabXgbFullVsOlsSharCloseLateP}{0.003}  % p-value
\newcommand{\rhabXgbFullVsOlsHarqAllDm}{$-$9.96}  % DM xgb full vs OLS harq
\newcommand{\rhabXgbFullVsOlsHarqAllP}{\ensuremath{<}0.001}  % p-value
\newcommand{\rhabXgbFullVsOlsHarqCloseDm}{$-$4.77}  % DM xgb full vs OLS harq
\newcommand{\rhabXgbFullVsOlsHarqCloseP}{\ensuremath{<}0.001}  % p-value
\newcommand{\rhabXgbFullVsOlsHarqCloseLateDm}{$-$2.27}  % DM xgb full vs OLS harq
\newcommand{\rhabXgbFullVsOlsHarqCloseLateP}{0.023}  % p-value
\newcommand{\rhabXgbFullVsOlsLharAllDm}{$-$7.85}  % DM xgb full vs OLS lhar
\newcommand{\rhabXgbFullVsOlsLharAllP}{\ensuremath{<}0.001}  % p-value
\newcommand{\rhabXgbFullVsOlsLharCloseDm}{$-$1.70}  % DM xgb full vs OLS lhar
\newcommand{\rhabXgbFullVsOlsLharCloseP}{0.090}  % p-value
\newcommand{\rhabXgbFullVsOlsLharCloseLateDm}{$-$1.85}  % DM xgb full vs OLS lhar
\newcommand{\rhabXgbFullVsOlsLharCloseLateP}{0.065}  % p-value
\newcommand{\rhabXgbFullVsOlsSharqlAllDm}{$-$7.76}  % DM xgb full vs OLS sharql
\newcommand{\rhabXgbFullVsOlsSharqlAllP}{\ensuremath{<}0.001}  % p-value
\newcommand{\rhabXgbFullVsOlsSharqlCloseDm}{$-$1.82}  % DM xgb full vs OLS sharql
\newcommand{\rhabXgbFullVsOlsSharqlCloseP}{0.069}  % p-value
\newcommand{\rhabXgbFullVsOlsSharqlCloseLateDm}{$-$1.95}  % DM xgb full vs OLS sharql
\newcommand{\rhabXgbFullVsOlsSharqlCloseLateP}{0.051}  % p-value
\newcommand{\rhabXgbFullVsOlsHarAllDm}{$-$9.93}  % DM xgb full vs OLS har
\newcommand{\rhabXgbFullVsOlsHarAllP}{\ensuremath{<}0.001}  % p-value
\newcommand{\rhabXgbFullVsOlsHarCloseDm}{$-$4.88}  % DM xgb full vs OLS har
\newcommand{\rhabXgbFullVsOlsHarCloseP}{\ensuremath{<}0.001}  % p-value
\newcommand{\rhabXgbFullVsOlsHarCloseLateDm}{$-$2.44}  % DM xgb full vs OLS har
\newcommand{\rhabXgbFullVsOlsHarCloseLateP}{0.015}  % p-value
\newcommand{\rhabRfFullVsOlsAddRmAllDm}{$-$6.68}  % DM rf full vs OLS add_rm
\newcommand{\rhabRfFullVsOlsAddRmAllP}{\ensuremath{<}0.001}  % p-value
\newcommand{\rhabRfFullVsOlsAddRmCloseDm}{$-$0.14}  % DM rf full vs OLS add_rm
\newcommand{\rhabRfFullVsOlsAddRmCloseP}{0.892}  % p-value
\newcommand{\rhabRfFullVsOlsAddRmCloseLateDm}{$-$1.75}  % DM rf full vs OLS add_rm
\newcommand{\rhabRfFullVsOlsAddRmCloseLateP}{0.080}  % p-value
\newcommand{\rhabRfFullVsOlsSharAllDm}{$-$8.45}  % DM rf full vs OLS shar
\newcommand{\rhabRfFullVsOlsSharAllP}{\ensuremath{<}0.001}  % p-value
\newcommand{\rhabRfFullVsOlsSharCloseDm}{$-$0.98}  % DM rf full vs OLS shar
\newcommand{\rhabRfFullVsOlsSharCloseP}{0.328}  % p-value
\newcommand{\rhabRfFullVsOlsSharCloseLateDm}{$-$0.76}  % DM rf full vs OLS shar
\newcommand{\rhabRfFullVsOlsSharCloseLateP}{0.444}  % p-value
\newcommand{\rhabRfFullVsOlsHarqAllDm}{$-$9.73}  % DM rf full vs OLS harq
\newcommand{\rhabRfFullVsOlsHarqAllP}{\ensuremath{<}0.001}  % p-value
\newcommand{\rhabRfFullVsOlsHarqCloseDm}{$-$1.05}  % DM rf full vs OLS harq
\newcommand{\rhabRfFullVsOlsHarqCloseP}{0.293}  % p-value
\newcommand{\rhabRfFullVsOlsHarqCloseLateDm}{$-$0.36}  % DM rf full vs OLS harq
\newcommand{\rhabRfFullVsOlsHarqCloseLateP}{0.715}  % p-value
\newcommand{\rhabRfFullVsOlsLharAllDm}{$-$7.08}  % DM rf full vs OLS lhar
\newcommand{\rhabRfFullVsOlsLharAllP}{\ensuremath{<}0.001}  % p-value
\newcommand{\rhabRfFullVsOlsLharCloseDm}{1.54}  % DM rf full vs OLS lhar
\newcommand{\rhabRfFullVsOlsLharCloseP}{0.124}  % p-value
\newcommand{\rhabRfFullVsOlsLharCloseLateDm}{$-$0.26}  % DM rf full vs OLS lhar
\newcommand{\rhabRfFullVsOlsLharCloseLateP}{0.794}  % p-value
\newcommand{\rhabRfFullVsOlsSharqlAllDm}{$-$6.99}  % DM rf full vs OLS sharql
\newcommand{\rhabRfFullVsOlsSharqlAllP}{\ensuremath{<}0.001}  % p-value
\newcommand{\rhabRfFullVsOlsSharqlCloseDm}{1.76}  % DM rf full vs OLS sharql
\newcommand{\rhabRfFullVsOlsSharqlCloseP}{0.078}  % p-value
\newcommand{\rhabRfFullVsOlsSharqlCloseLateDm}{$-$0.09}  % DM rf full vs OLS sharql
\newcommand{\rhabRfFullVsOlsSharqlCloseLateP}{0.926}  % p-value
\newcommand{\rhabRfFullVsOlsHarAllDm}{$-$9.69}  % DM rf full vs OLS har
\newcommand{\rhabRfFullVsOlsHarAllP}{\ensuremath{<}0.001}  % p-value
\newcommand{\rhabRfFullVsOlsHarCloseDm}{$-$1.16}  % DM rf full vs OLS har
\newcommand{\rhabRfFullVsOlsHarCloseP}{0.246}  % p-value
\newcommand{\rhabRfFullVsOlsHarCloseLateDm}{$-$0.49}  % DM rf full vs OLS har
\newcommand{\rhabRfFullVsOlsHarCloseLateP}{0.621}  % p-value
\newcommand{\rhabNClassesAddRm}{three}  % classes with an add_rm arm
\newcommand{\rhabNClassesBeatAddRm}{none}  % full set beats its own add_rm, DM p < 0.05
\newcommand{\rhabBestOlsBench}{SHAR + HARQ + LHAR}  % lowest QLIKE OLS augmented HAR
\newcommand{\rhabLiveAttnFirst}{23~March 2009}  % first print stocktwits_attention
\newcommand{\rhabLiveAttnLast}{30~April 2024}  % last print stocktwits_attention
\newcommand{\rhabLiveAttnFromKone}{11~February 2021}  % first non-zeroed value, k = 1
\newcommand{\rhabLiveAttnFromKlast}{13~May 2021}  % the same, k = 3125
\newcommand{\rhabLiveAttnSharePct}{31}  % % of evaluation bars live, k = 1
\newcommand{\rhabLiveNSentFirst}{1~June 2010}  % first print stocktwits_sentcount
\newcommand{\rhabLiveNSentLast}{30~April 2024}  % last print stocktwits_sentcount
\newcommand{\rhabLiveNSentFromKone}{19~April 2022}  % first non-zeroed value, k = 1
\newcommand{\rhabLiveNSentFromKlast}{19~July 2022}  % the same, k = 3125
\newcommand{\rhabLiveNSentSharePct}{20}  % % of evaluation bars live, k = 1
\newcommand{\rhabLiveSentFirst}{1~June 2010}  % first print stocktwits_sentiment
\newcommand{\rhabLiveSentLast}{30~April 2024}  % last print stocktwits_sentiment
\newcommand{\rhabLiveSentFromKone}{19~April 2022}  % first non-zeroed value, k = 1
\newcommand{\rhabLiveSentFromKlast}{19~July 2022}  % the same, k = 3125
\newcommand{\rhabLiveSentSharePct}{20}  % % of evaluation bars live, k = 1
\newcommand{\rhabLiveVixFirst}{22~May 2006}  % first print vix
\newcommand{\rhabLiveVixLast}{30~April 2024}  % last print vix
\newcommand{\rhabLiveVixFromKone}{19~April 2018}  % first non-zeroed value, k = 1
\newcommand{\rhabLiveVixFromKlast}{18~July 2018}  % the same, k = 3125
\newcommand{\rhabLiveVixSharePct}{58}  % % of evaluation bars live, k = 1
\newcommand{\rhabLiveVixthreemFirst}{12~August 2009}  % first print vix3m
\newcommand{\rhabLiveVixthreemLast}{30~April 2024}  % last print vix3m
\newcommand{\rhabLiveVixthreemFromKone}{6~July 2021}  % first non-zeroed value, k = 1
\newcommand{\rhabLiveVixthreemFromKlast}{1~October 2021}  % the same, k = 3125
\newcommand{\rhabLiveVixthreemSharePct}{27}  % % of evaluation bars live, k = 1
\newcommand{\rhabLiveVvixFirst}{26~March 2012}  % first print vvix
\newcommand{\rhabLiveVvixLast}{30~April 2024}  % last print vvix
\newcommand{\rhabLiveVvixFromKone}{12~February 2024}  % first non-zeroed value, k = 1
\newcommand{\rhabLiveVvixFromKlast}{never}  % the same, k = 3125
\newcommand{\rhabLiveVvixSharePct}{2}  % % of evaluation bars live, k = 1
\newcommand{\rhabLiveVdAllOcFirst}{3~January 2011}  % first print voldemand_all_open_and_close
\newcommand{\rhabLiveVdAllOcLast}{31~August 2023}  % last print voldemand_all_open_and_close
\newcommand{\rhabLiveVdAllOcFromKone}{16~November 2022}  % first non-zeroed value, k = 1
\newcommand{\rhabLiveVdAllOcFromKlast}{20~February 2023}  % the same, k = 3125
\newcommand{\rhabLiveVdAllOcSharePct}{8}  % % of evaluation bars live, k = 1
\newcommand{\rhabLiveVdAllOFirst}{3~January 2011}  % first print voldemand_all_open_only
\newcommand{\rhabLiveVdAllOLast}{31~August 2023}  % last print voldemand_all_open_only
\newcommand{\rhabLiveVdAllOFromKone}{16~November 2022}  % first non-zeroed value, k = 1
\newcommand{\rhabLiveVdAllOFromKlast}{20~February 2023}  % the same, k = 3125
\newcommand{\rhabLiveVdAllOSharePct}{8}  % % of evaluation bars live, k = 1
\newcommand{\rhabLiveVdSpxOcFirst}{3~January 2011}  % first print voldemand_spx_open_and_close
\newcommand{\rhabLiveVdSpxOcLast}{31~August 2023}  % last print voldemand_spx_open_and_close
\newcommand{\rhabLiveVdSpxOcFromKone}{16~November 2022}  % first non-zeroed value, k = 1
\newcommand{\rhabLiveVdSpxOcFromKlast}{20~February 2023}  % the same, k = 3125
\newcommand{\rhabLiveVdSpxOcSharePct}{8}  % % of evaluation bars live, k = 1
\newcommand{\rhabLiveVdSpxOFirst}{3~January 2011}  % first print voldemand_spx_open_only
\newcommand{\rhabLiveVdSpxOLast}{31~August 2023}  % last print voldemand_spx_open_only
\newcommand{\rhabLiveVdSpxOFromKone}{16~November 2022}  % first non-zeroed value, k = 1
\newcommand{\rhabLiveVdSpxOFromKlast}{20~February 2023}  % the same, k = 3125
\newcommand{\rhabLiveVdSpxOSharePct}{8}  % % of evaluation bars live, k = 1
 in the
% preamble, the table in 06_results_floats.tex, this text in Section~\ref{sec:rh_results_predictors}.
% [RESULTS PARAGRAPHS: written after the runs finish]
% Two references are not in writeup/references.bib yet (BollerslevPattonQuaedvlieg2016 is):
% @article{PattonSheppard2015,
%   title={Good volatility, bad volatility: signed jumps and the persistence of volatility},
%   author={Patton, Andrew J. and Sheppard, Kevin}, journal={The Review of Economics and Statistics},
%   volume={97}, number={3}, pages={683--697}, year={2015}}
% @article{CorsiReno2012,
%   title={Discrete-time volatility forecasting with persistent leverage effect and the link with continuous-time
%          volatility modeling},
%   author={Corsi, Fulvio and Ren{\`o}, Roberto}, journal={Journal of Business \& Economic Statistics},
%   volume={30}, number={3}, pages={368--380}, year={2012}}

\subsection{Which predictors carry the gain?}
\label{sec:rh_results_ablation}

The comparison of Table~\ref{tab:rh_oos} sets the full predictor set against the HAR benchmark.  Two questions remain:
whether the gain comes from data beyond the futures contract itself, and whether a richer volatility model built from
the same one-minute E-mini returns would close the gap.  We answer both with the design of Table~\ref{tab:rh_oos}
unchanged---the same rows, window, re-estimation at every bar, scored bars and test---and only the columns changed.

\paragraph{Design.}  The 46 exogenous series form the eight groups of Table~\ref{tab:rh_dm_dictionary}.  For each group
$g$ we fit two models in every class: the full set without $g$ (all of $g$'s trailing averages $u^{(k)}$ and
availability shares $\bar a^{(k)}$ removed) and the HAR benchmark with $g$ added.  The first measures what $g$
contributes given everything else; the second what it contributes on its own.  Hyperparameters stay those of the full
set: the linear penalties chosen on the 2013 pre-sample fold for the full predictor set (with the HAR block
penalized at one hundredth of the rate of the exogenous columns), and the single frozen configuration of each tree
method.  The random forest, the most expensive class, is fitted for the eight additions and for the removal of the
realized moments.  Panel~D of Table~\ref{tab:rh_ablation} restates the full-set row of Table~\ref{tab:rh_oos}.

\paragraph{Augmented HAR benchmarks.}  The first group, the futures' own realized moments, is the natural
``augmented HAR'': it adds the return, absolute variation, skewness, quarticity, positive semivariance, bipower
variation and autocovariance of the same one-minute E-mini returns that define the target.  We also fit the three
classical extensions of the HAR model by OLS, each built from the same one-minute data and passed through the same
diurnal adjustment, variance-stabilizing transform and trailing averages as every other series: the semivariance HAR
of \citet{PattonSheppard2015}, which adds the positive and negative realized semivariances $\mathrm{RS}^{\pm}$ of the
last half hour; the HARQ model of \citet{BollerslevPattonQuaedvlieg2016}, which adds the interaction of the last
half hour's realized variance with the square root of its realized quarticity $\mathrm{RQ} = (N/3)\sum_k r^4_{t,i,k}$
(in the volatility units of the design, $\sqrt{\RV\,\mathrm{RQ}^{1/2}}$); and the leverage HAR of
\citet{CorsiReno2012}, which adds the negative part of the half-hour return, $\min(r_{t,i}, 0)$, at every HAR horizon.
A fourth model adds all three.  The negative semivariance is $\mathrm{RS}^{-} = \RV - \mathrm{RS}^{+}$, exact on every
bar of the data.

\paragraph{Late-starting series.}  The implied-volatility, sentiment and volatility-demand series start after the
data, and a series enters as zero until the rolling window holds 3,000 trading days of its history
(Section~\ref{sec:rh_predictors}).  In the evaluation sample this rule binds for most of the span: the VIX enters on
\rhabLiveVixFromKone{} (its longest average on \rhabLiveVixFromKlast{}), VIX3M on \rhabLiveVixthreemFromKone{},
StockTwits attention on \rhabLiveAttnFromKone{} and the other StockTwits series on \rhabLiveSentFromKone{}, the
volatility-demand series on \rhabLiveVdSpxOcFromKone{} (and their feed stops on \rhabLiveVdSpxOcLast{}... [check]),
and VVIX on \rhabLiveVvixFromKone{}.  Before those dates the groups contribute only their availability shares, so
the ablation rows of these groups measure the series over the later part of the sample only.
